% documentation_FR.tex — WAC (Windows Artefacts Collector) user manual, French edition.
% The English edition is documentation_EN.tex, in this folder: both are maintained together,
% section for section.
%
% Build (LuaLaTeX, as the editor's recipe; build artefacts removed after):
%     ../build-docs.sh documentation_FR.tex
% The cover page is ../cover_FR.png.

\documentclass[11pt,a4paper,oneside]{report}

% ---------------------------------------------------------------- language and labels
\usepackage[french]{babel}
\newcommand{\wacdoctitle}{Documentation utilisateur}
\newcommand{\wacpageof}{sur}
\newcommand{\wacnotetitle}{À retenir}
\newcommand{\wacwarningtitle}{Attention}
% Layout shared with every WAC document.
% wac-style.tex — layout shared by every WAC document (user manuals, testing
% guides), French and English editions alike: one copy, so that they cannot
% drift apart.
%
% A document sets, BEFORE % wac-style.tex — layout shared by every WAC document (user manuals, testing
% guides), French and English editions alike: one copy, so that they cannot
% drift apart.
%
% A document sets, BEFORE % wac-style.tex — layout shared by every WAC document (user manuals, testing
% guides), French and English editions alike: one copy, so that they cannot
% drift apart.
%
% A document sets, BEFORE \input{../wac-style.tex}:
%   \usepackage[<language>]{babel}
%   \newcommand{\wacdoctitle}{...}      running header: "WAC — <title>"
%   \newcommand{\wacpageof}{...}        footer: "Page N <of> M"
%   \newcommand{\wacnotetitle}{...}     default title of the "note" box
%   \newcommand{\wacwarningtitle}{...}  default title of the "attention" box
% and AFTER it its \hypersetup (PDF title, author, subject).

% ---------------------------------------------------------------- fonts, language
\usepackage{fontspec}
\setmainfont{Noto Sans}
\setmonofont{DejaVu Sans Mono}[Scale=0.84]

% ---------------------------------------------------------------- layout
\usepackage[top=26mm,bottom=26mm,left=24mm,right=22mm,headheight=14pt]{geometry}
\usepackage{graphicx}
\graphicspath{{../}}
\usepackage{xcolor}
\usepackage{tikz}
\usepackage{booktabs,tabularx,longtable,array}
\usepackage{enumitem}
\usepackage{fancyhdr}
\usepackage{titlesec}
\usepackage[most]{tcolorbox}
\usepackage{listings}
\usepackage{url}
\usepackage{lastpage}
\usepackage[hidelinks,unicode,pdfpagelabels]{hyperref}

% Colours sampled from the cover page
\definecolor{marine}{HTML}{0A1A33}
\definecolor{bleu}{HTML}{2F6DB5}
\definecolor{accent}{HTML}{1F66D1}
\definecolor{fond}{HTML}{EEF3FA}
\definecolor{alerte}{HTML}{B3261E}
\definecolor{fondalerte}{HTML}{FBEDEC}
\definecolor{gris}{HTML}{5B6573}

\setlength{\parindent}{0pt}
\setlength{\parskip}{0.55em}
\linespread{1.08}
\setlist{itemsep=0.25em,topsep=0.35em}
\urlstyle{tt}
\renewcommand{\arraystretch}{1.25}
% Left-aligned table columns: justified narrow columns opened wide gaps
% between words.
\newcolumntype{L}[1]{>{\raggedright\arraybackslash}p{#1}}
% Monospaced file and field names cannot be hyphenated: let the line's spaces
% stretch instead.
\setlength{\emergencystretch}{3em}

% Headings
\titleformat{\chapter}[display]
  {\color{marine}\bfseries\raggedright}
  {\color{bleu}\large\MakeUppercase{\chaptertitlename}~\thechapter}
  {0.3em}
  {\Huge}[{\color{accent}\titlerule[2pt]}]
\titlespacing*{\chapter}{0pt}{0pt}{1.6em}
\titleformat{\section}{\color{marine}\Large\bfseries}{\color{bleu}\thesection}{0.8em}{}
\titleformat{\subsection}{\color{marine}\large\bfseries}{\color{bleu}\thesubsection}{0.7em}{}
\titleformat{\subsubsection}{\color{bleu}\normalsize\bfseries}{}{0em}{}

% Headers and footers
\pagestyle{fancy}
\fancyhf{}
\fancyhead[L]{\small\color{gris}WAC — \wacdoctitle}
\fancyhead[R]{\small\color{gris}\leftmark}
% Footer: "Page N of M" (the word set by \wacpageof), readable. A bare, small, light-grey number went
% unnoticed.
\newcommand{\piedpage}{\small\color{marine}\textbf{Page \thepage}\ \wacpageof\ \pageref*{LastPage}}
\fancyfoot[C]{\piedpage}
\renewcommand{\footrulewidth}{0.4pt}
\renewcommand{\footrule}{\hbox to\headwidth{\color{fond}\leaders\hrule height \footrulewidth\hfill}}
\renewcommand{\headrulewidth}{0.4pt}
\renewcommand{\headrule}{\hbox to\headwidth{\color{fond}\leaders\hrule height \headrulewidth\hfill}}
\renewcommand{\chaptermark}[1]{\markboth{#1}{}}
\fancypagestyle{plain}{\fancyhf{}\fancyfoot[C]{\piedpage}\renewcommand{\headrulewidth}{0pt}}

% Boxes
\newtcolorbox{note}[1][\wacnotetitle]{enhanced,breakable,colback=fond,colframe=bleu,
  boxrule=0pt,leftrule=3pt,arc=0pt,left=8pt,right=8pt,top=6pt,bottom=6pt,
  title={\bfseries #1},coltitle=bleu,colbacktitle=fond,attach title to upper={\par\smallskip}}
\newtcolorbox{attention}[1][\wacwarningtitle]{enhanced,breakable,colback=fondalerte,colframe=alerte,
  boxrule=0pt,leftrule=3pt,arc=0pt,left=8pt,right=8pt,top=6pt,bottom=6pt,
  title={\bfseries #1},coltitle=alerte,colbacktitle=fondalerte,attach title to upper={\par\smallskip}}

% Code
\lstset{basicstyle=\ttfamily\small,columns=fullflexible,keepspaces=true,
  backgroundcolor=\color{fond},frame=l,framerule=3pt,rulecolor=\color{bleu},
  xleftmargin=9pt,framexleftmargin=6pt,aboveskip=0.8em,belowskip=0.8em,
  breaklines=true,breakatwhitespace=false,showstringspaces=false,
  literate={é}{{\'e}}1 {è}{{\`e}}1 {à}{{\`a}}1 {ê}{{\^e}}1 {ô}{{\^o}}1 {ç}{{\c c}}1
           {É}{{\'E}}1 {—}{{---}}1 {’}{{'}}1 {…}{{...}}1 {«}{{\guillemotleft}}1 {»}{{\guillemotright}}1}

% Breakable monospace. A path or a field name in \texttt is one word that no
% hyphenation pattern cuts: in a table cell it ran 190 pt past the column
% (ProgramData\Microsoft\NetFramework\BreadcrumbStore). As the \seqsplit of
% the EventsManager theme's tables, a line may break inside it, without a
% hyphen that would read as part of the name: after \ / . _ - : and between a
% lowercase letter or digit and a capital (AmcacheSha1|Covers|Bytes). Only
% inside \texttt — the classes are set once, the breaks switched on there.
\newXeTeXintercharclass\wacBreakAfter
\newXeTeXintercharclass\wacUpper
\newXeTeXintercharclass\wacLower
\XeTeXcharclass`\\=\wacBreakAfter
\XeTeXcharclass`\/=\wacBreakAfter
\XeTeXcharclass`\.=\wacBreakAfter
\XeTeXcharclass`\_=\wacBreakAfter
\XeTeXcharclass`\-=\wacBreakAfter
\XeTeXcharclass`\:=\wacBreakAfter
\count255=`A \loop\XeTeXcharclass\count255=\wacUpper\ifnum\count255<`Z\advance\count255 by 1\repeat
\count255=`a \loop\XeTeXcharclass\count255=\wacLower\ifnum\count255<`z\advance\count255 by 1\repeat
\count255=`0 \loop\XeTeXcharclass\count255=\wacLower\ifnum\count255<`9\advance\count255 by 1\repeat
\XeTeXinterchartoks\wacBreakAfter\wacUpper={\allowbreak}
\XeTeXinterchartoks\wacBreakAfter\wacLower={\allowbreak}
\XeTeXinterchartoks\wacBreakAfter 0={\allowbreak}
\XeTeXinterchartoks\wacLower\wacUpper={\allowbreak}
% \textbackslash is no character token: it carries its own break.
\newcommand{\wacbackslash}{\char92\allowbreak}
\let\wacplaintexttt\texttt
\renewcommand{\texttt}[1]{\wacplaintexttt{\XeTeXinterchartokenstate=1\let\textbackslash\wacbackslash#1}}

% Shortcuts
\newcommand{\wac}{\textsc{wac}}
\newcommand{\opt}[1]{\texttt{-{}-#1}}
\newcommand{\champ}[1]{\texttt{#1}}
\newcommand{\fichier}[1]{\texttt{#1}}

% ============================================================ COVER PAGE
% The cover images are 2:3 (1024 x 1536), taller than an A4 page once widened
% to it: centred, they lost 9 mm at the top — the CEC logo — and 9 mm at the
% bottom. Raised 4 mm above the top of the page instead: the logo, which starts
% 6 mm below the image's edge, stays whole, and the bottom row of text keeps
% 7 mm to the page's edge; only a decorative corner below it is cut.
% Usage: \coverpage{cover_FR.png}
\newcommand{\coverpage}[1]{%
  \begin{titlepage}
  \thispagestyle{empty}
  \begin{tikzpicture}[remember picture,overlay]
    \node[inner sep=0pt,anchor=north,yshift=4mm] at (current page.north)
      {\includegraphics[width=\paperwidth]{#1}};
  \end{tikzpicture}
  \end{titlepage}}
:
%   \usepackage[<language>]{babel}
%   \newcommand{\wacdoctitle}{...}      running header: "WAC — <title>"
%   \newcommand{\wacpageof}{...}        footer: "Page N <of> M"
%   \newcommand{\wacnotetitle}{...}     default title of the "note" box
%   \newcommand{\wacwarningtitle}{...}  default title of the "attention" box
% and AFTER it its \hypersetup (PDF title, author, subject).

% ---------------------------------------------------------------- fonts, language
\usepackage{fontspec}
\setmainfont{Noto Sans}
\setmonofont{DejaVu Sans Mono}[Scale=0.84]

% ---------------------------------------------------------------- layout
\usepackage[top=26mm,bottom=26mm,left=24mm,right=22mm,headheight=14pt]{geometry}
\usepackage{graphicx}
\graphicspath{{../}}
\usepackage{xcolor}
\usepackage{tikz}
\usepackage{booktabs,tabularx,longtable,array}
\usepackage{enumitem}
\usepackage{fancyhdr}
\usepackage{titlesec}
\usepackage[most]{tcolorbox}
\usepackage{listings}
\usepackage{url}
\usepackage{lastpage}
\usepackage[hidelinks,unicode,pdfpagelabels]{hyperref}

% Colours sampled from the cover page
\definecolor{marine}{HTML}{0A1A33}
\definecolor{bleu}{HTML}{2F6DB5}
\definecolor{accent}{HTML}{1F66D1}
\definecolor{fond}{HTML}{EEF3FA}
\definecolor{alerte}{HTML}{B3261E}
\definecolor{fondalerte}{HTML}{FBEDEC}
\definecolor{gris}{HTML}{5B6573}

\setlength{\parindent}{0pt}
\setlength{\parskip}{0.55em}
\linespread{1.08}
\setlist{itemsep=0.25em,topsep=0.35em}
\urlstyle{tt}
\renewcommand{\arraystretch}{1.25}
% Left-aligned table columns: justified narrow columns opened wide gaps
% between words.
\newcolumntype{L}[1]{>{\raggedright\arraybackslash}p{#1}}
% Monospaced file and field names cannot be hyphenated: let the line's spaces
% stretch instead.
\setlength{\emergencystretch}{3em}

% Headings
\titleformat{\chapter}[display]
  {\color{marine}\bfseries\raggedright}
  {\color{bleu}\large\MakeUppercase{\chaptertitlename}~\thechapter}
  {0.3em}
  {\Huge}[{\color{accent}\titlerule[2pt]}]
\titlespacing*{\chapter}{0pt}{0pt}{1.6em}
\titleformat{\section}{\color{marine}\Large\bfseries}{\color{bleu}\thesection}{0.8em}{}
\titleformat{\subsection}{\color{marine}\large\bfseries}{\color{bleu}\thesubsection}{0.7em}{}
\titleformat{\subsubsection}{\color{bleu}\normalsize\bfseries}{}{0em}{}

% Headers and footers
\pagestyle{fancy}
\fancyhf{}
\fancyhead[L]{\small\color{gris}WAC — \wacdoctitle}
\fancyhead[R]{\small\color{gris}\leftmark}
% Footer: "Page N of M" (the word set by \wacpageof), readable. A bare, small, light-grey number went
% unnoticed.
\newcommand{\piedpage}{\small\color{marine}\textbf{Page \thepage}\ \wacpageof\ \pageref*{LastPage}}
\fancyfoot[C]{\piedpage}
\renewcommand{\footrulewidth}{0.4pt}
\renewcommand{\footrule}{\hbox to\headwidth{\color{fond}\leaders\hrule height \footrulewidth\hfill}}
\renewcommand{\headrulewidth}{0.4pt}
\renewcommand{\headrule}{\hbox to\headwidth{\color{fond}\leaders\hrule height \headrulewidth\hfill}}
\renewcommand{\chaptermark}[1]{\markboth{#1}{}}
\fancypagestyle{plain}{\fancyhf{}\fancyfoot[C]{\piedpage}\renewcommand{\headrulewidth}{0pt}}

% Boxes
\newtcolorbox{note}[1][\wacnotetitle]{enhanced,breakable,colback=fond,colframe=bleu,
  boxrule=0pt,leftrule=3pt,arc=0pt,left=8pt,right=8pt,top=6pt,bottom=6pt,
  title={\bfseries #1},coltitle=bleu,colbacktitle=fond,attach title to upper={\par\smallskip}}
\newtcolorbox{attention}[1][\wacwarningtitle]{enhanced,breakable,colback=fondalerte,colframe=alerte,
  boxrule=0pt,leftrule=3pt,arc=0pt,left=8pt,right=8pt,top=6pt,bottom=6pt,
  title={\bfseries #1},coltitle=alerte,colbacktitle=fondalerte,attach title to upper={\par\smallskip}}

% Code
\lstset{basicstyle=\ttfamily\small,columns=fullflexible,keepspaces=true,
  backgroundcolor=\color{fond},frame=l,framerule=3pt,rulecolor=\color{bleu},
  xleftmargin=9pt,framexleftmargin=6pt,aboveskip=0.8em,belowskip=0.8em,
  breaklines=true,breakatwhitespace=false,showstringspaces=false,
  literate={é}{{\'e}}1 {è}{{\`e}}1 {à}{{\`a}}1 {ê}{{\^e}}1 {ô}{{\^o}}1 {ç}{{\c c}}1
           {É}{{\'E}}1 {—}{{---}}1 {’}{{'}}1 {…}{{...}}1 {«}{{\guillemotleft}}1 {»}{{\guillemotright}}1}

% Breakable monospace. A path or a field name in \texttt is one word that no
% hyphenation pattern cuts: in a table cell it ran 190 pt past the column
% (ProgramData\Microsoft\NetFramework\BreadcrumbStore). As the \seqsplit of
% the EventsManager theme's tables, a line may break inside it, without a
% hyphen that would read as part of the name: after \ / . _ - : and between a
% lowercase letter or digit and a capital (AmcacheSha1|Covers|Bytes). Only
% inside \texttt — the classes are set once, the breaks switched on there.
\newXeTeXintercharclass\wacBreakAfter
\newXeTeXintercharclass\wacUpper
\newXeTeXintercharclass\wacLower
\XeTeXcharclass`\\=\wacBreakAfter
\XeTeXcharclass`\/=\wacBreakAfter
\XeTeXcharclass`\.=\wacBreakAfter
\XeTeXcharclass`\_=\wacBreakAfter
\XeTeXcharclass`\-=\wacBreakAfter
\XeTeXcharclass`\:=\wacBreakAfter
\count255=`A \loop\XeTeXcharclass\count255=\wacUpper\ifnum\count255<`Z\advance\count255 by 1\repeat
\count255=`a \loop\XeTeXcharclass\count255=\wacLower\ifnum\count255<`z\advance\count255 by 1\repeat
\count255=`0 \loop\XeTeXcharclass\count255=\wacLower\ifnum\count255<`9\advance\count255 by 1\repeat
\XeTeXinterchartoks\wacBreakAfter\wacUpper={\allowbreak}
\XeTeXinterchartoks\wacBreakAfter\wacLower={\allowbreak}
\XeTeXinterchartoks\wacBreakAfter 0={\allowbreak}
\XeTeXinterchartoks\wacLower\wacUpper={\allowbreak}
% \textbackslash is no character token: it carries its own break.
\newcommand{\wacbackslash}{\char92\allowbreak}
\let\wacplaintexttt\texttt
\renewcommand{\texttt}[1]{\wacplaintexttt{\XeTeXinterchartokenstate=1\let\textbackslash\wacbackslash#1}}

% Shortcuts
\newcommand{\wac}{\textsc{wac}}
\newcommand{\opt}[1]{\texttt{-{}-#1}}
\newcommand{\champ}[1]{\texttt{#1}}
\newcommand{\fichier}[1]{\texttt{#1}}

% ============================================================ COVER PAGE
% The cover images are 2:3 (1024 x 1536), taller than an A4 page once widened
% to it: centred, they lost 9 mm at the top — the CEC logo — and 9 mm at the
% bottom. Raised 4 mm above the top of the page instead: the logo, which starts
% 6 mm below the image's edge, stays whole, and the bottom row of text keeps
% 7 mm to the page's edge; only a decorative corner below it is cut.
% Usage: \coverpage{cover_FR.png}
\newcommand{\coverpage}[1]{%
  \begin{titlepage}
  \thispagestyle{empty}
  \begin{tikzpicture}[remember picture,overlay]
    \node[inner sep=0pt,anchor=north,yshift=4mm] at (current page.north)
      {\includegraphics[width=\paperwidth]{#1}};
  \end{tikzpicture}
  \end{titlepage}}
:
%   \usepackage[<language>]{babel}
%   \newcommand{\wacdoctitle}{...}      running header: "WAC — <title>"
%   \newcommand{\wacpageof}{...}        footer: "Page N <of> M"
%   \newcommand{\wacnotetitle}{...}     default title of the "note" box
%   \newcommand{\wacwarningtitle}{...}  default title of the "attention" box
% and AFTER it its \hypersetup (PDF title, author, subject).

% ---------------------------------------------------------------- fonts, language
\usepackage{fontspec}
\setmainfont{Noto Sans}
\setmonofont{DejaVu Sans Mono}[Scale=0.84]

% ---------------------------------------------------------------- layout
\usepackage[top=26mm,bottom=26mm,left=24mm,right=22mm,headheight=14pt]{geometry}
\usepackage{graphicx}
\graphicspath{{../}}
\usepackage{xcolor}
\usepackage{tikz}
\usepackage{booktabs,tabularx,longtable,array}
\usepackage{enumitem}
\usepackage{fancyhdr}
\usepackage{titlesec}
\usepackage[most]{tcolorbox}
\usepackage{listings}
\usepackage{url}
\usepackage{lastpage}
\usepackage[hidelinks,unicode,pdfpagelabels]{hyperref}

% Colours sampled from the cover page
\definecolor{marine}{HTML}{0A1A33}
\definecolor{bleu}{HTML}{2F6DB5}
\definecolor{accent}{HTML}{1F66D1}
\definecolor{fond}{HTML}{EEF3FA}
\definecolor{alerte}{HTML}{B3261E}
\definecolor{fondalerte}{HTML}{FBEDEC}
\definecolor{gris}{HTML}{5B6573}

\setlength{\parindent}{0pt}
\setlength{\parskip}{0.55em}
\linespread{1.08}
\setlist{itemsep=0.25em,topsep=0.35em}
\urlstyle{tt}
\renewcommand{\arraystretch}{1.25}
% Left-aligned table columns: justified narrow columns opened wide gaps
% between words.
\newcolumntype{L}[1]{>{\raggedright\arraybackslash}p{#1}}
% Monospaced file and field names cannot be hyphenated: let the line's spaces
% stretch instead.
\setlength{\emergencystretch}{3em}

% Headings
\titleformat{\chapter}[display]
  {\color{marine}\bfseries\raggedright}
  {\color{bleu}\large\MakeUppercase{\chaptertitlename}~\thechapter}
  {0.3em}
  {\Huge}[{\color{accent}\titlerule[2pt]}]
\titlespacing*{\chapter}{0pt}{0pt}{1.6em}
\titleformat{\section}{\color{marine}\Large\bfseries}{\color{bleu}\thesection}{0.8em}{}
\titleformat{\subsection}{\color{marine}\large\bfseries}{\color{bleu}\thesubsection}{0.7em}{}
\titleformat{\subsubsection}{\color{bleu}\normalsize\bfseries}{}{0em}{}

% Headers and footers
\pagestyle{fancy}
\fancyhf{}
\fancyhead[L]{\small\color{gris}WAC — \wacdoctitle}
\fancyhead[R]{\small\color{gris}\leftmark}
% Footer: "Page N of M" (the word set by \wacpageof), readable. A bare, small, light-grey number went
% unnoticed.
\newcommand{\piedpage}{\small\color{marine}\textbf{Page \thepage}\ \wacpageof\ \pageref*{LastPage}}
\fancyfoot[C]{\piedpage}
\renewcommand{\footrulewidth}{0.4pt}
\renewcommand{\footrule}{\hbox to\headwidth{\color{fond}\leaders\hrule height \footrulewidth\hfill}}
\renewcommand{\headrulewidth}{0.4pt}
\renewcommand{\headrule}{\hbox to\headwidth{\color{fond}\leaders\hrule height \headrulewidth\hfill}}
\renewcommand{\chaptermark}[1]{\markboth{#1}{}}
\fancypagestyle{plain}{\fancyhf{}\fancyfoot[C]{\piedpage}\renewcommand{\headrulewidth}{0pt}}

% Boxes
\newtcolorbox{note}[1][\wacnotetitle]{enhanced,breakable,colback=fond,colframe=bleu,
  boxrule=0pt,leftrule=3pt,arc=0pt,left=8pt,right=8pt,top=6pt,bottom=6pt,
  title={\bfseries #1},coltitle=bleu,colbacktitle=fond,attach title to upper={\par\smallskip}}
\newtcolorbox{attention}[1][\wacwarningtitle]{enhanced,breakable,colback=fondalerte,colframe=alerte,
  boxrule=0pt,leftrule=3pt,arc=0pt,left=8pt,right=8pt,top=6pt,bottom=6pt,
  title={\bfseries #1},coltitle=alerte,colbacktitle=fondalerte,attach title to upper={\par\smallskip}}

% Code
\lstset{basicstyle=\ttfamily\small,columns=fullflexible,keepspaces=true,
  backgroundcolor=\color{fond},frame=l,framerule=3pt,rulecolor=\color{bleu},
  xleftmargin=9pt,framexleftmargin=6pt,aboveskip=0.8em,belowskip=0.8em,
  breaklines=true,breakatwhitespace=false,showstringspaces=false,
  literate={é}{{\'e}}1 {è}{{\`e}}1 {à}{{\`a}}1 {ê}{{\^e}}1 {ô}{{\^o}}1 {ç}{{\c c}}1
           {É}{{\'E}}1 {—}{{---}}1 {’}{{'}}1 {…}{{...}}1 {«}{{\guillemotleft}}1 {»}{{\guillemotright}}1}

% Breakable monospace. A path or a field name in \texttt is one word that no
% hyphenation pattern cuts: in a table cell it ran 190 pt past the column
% (ProgramData\Microsoft\NetFramework\BreadcrumbStore). As the \seqsplit of
% the EventsManager theme's tables, a line may break inside it, without a
% hyphen that would read as part of the name: after \ / . _ - : and between a
% lowercase letter or digit and a capital (AmcacheSha1|Covers|Bytes). Only
% inside \texttt — the classes are set once, the breaks switched on there.
\newXeTeXintercharclass\wacBreakAfter
\newXeTeXintercharclass\wacUpper
\newXeTeXintercharclass\wacLower
\XeTeXcharclass`\\=\wacBreakAfter
\XeTeXcharclass`\/=\wacBreakAfter
\XeTeXcharclass`\.=\wacBreakAfter
\XeTeXcharclass`\_=\wacBreakAfter
\XeTeXcharclass`\-=\wacBreakAfter
\XeTeXcharclass`\:=\wacBreakAfter
\count255=`A \loop\XeTeXcharclass\count255=\wacUpper\ifnum\count255<`Z\advance\count255 by 1\repeat
\count255=`a \loop\XeTeXcharclass\count255=\wacLower\ifnum\count255<`z\advance\count255 by 1\repeat
\count255=`0 \loop\XeTeXcharclass\count255=\wacLower\ifnum\count255<`9\advance\count255 by 1\repeat
\XeTeXinterchartoks\wacBreakAfter\wacUpper={\allowbreak}
\XeTeXinterchartoks\wacBreakAfter\wacLower={\allowbreak}
\XeTeXinterchartoks\wacBreakAfter 0={\allowbreak}
\XeTeXinterchartoks\wacLower\wacUpper={\allowbreak}
% \textbackslash is no character token: it carries its own break.
\newcommand{\wacbackslash}{\char92\allowbreak}
\let\wacplaintexttt\texttt
\renewcommand{\texttt}[1]{\wacplaintexttt{\XeTeXinterchartokenstate=1\let\textbackslash\wacbackslash#1}}

% Shortcuts
\newcommand{\wac}{\textsc{wac}}
\newcommand{\opt}[1]{\texttt{-{}-#1}}
\newcommand{\champ}[1]{\texttt{#1}}
\newcommand{\fichier}[1]{\texttt{#1}}

% ============================================================ COVER PAGE
% The cover images are 2:3 (1024 x 1536), taller than an A4 page once widened
% to it: centred, they lost 9 mm at the top — the CEC logo — and 9 mm at the
% bottom. Raised 4 mm above the top of the page instead: the logo, which starts
% 6 mm below the image's edge, stays whole, and the bottom row of text keeps
% 7 mm to the page's edge; only a decorative corner below it is cut.
% Usage: \coverpage{cover_FR.png}
\newcommand{\coverpage}[1]{%
  \begin{titlepage}
  \thispagestyle{empty}
  \begin{tikzpicture}[remember picture,overlay]
    \node[inner sep=0pt,anchor=north,yshift=4mm] at (current page.north)
      {\includegraphics[width=\paperwidth]{#1}};
  \end{tikzpicture}
  \end{titlepage}}


\hypersetup{pdftitle={WAC — Documentation utilisateur},
  pdfauthor={Centre d'Excellence Cyberdéfense Aérospatiale (CEC)},
  pdfsubject={Logiciel de prélèvement d'artefacts Windows pour l'investigation numérique}}

\begin{document}

% ============================================================ COVER PAGE
\pagenumbering{roman}
\coverpage{cover_FR.png}

% ============================================================ INFORMATION PAGE
\thispagestyle{empty}
\vspace*{\fill}
{\color{marine}\Large\bfseries WAC — Windows Artefacts Collector}\par
{\color{bleu}Documentation utilisateur}\par
\medskip
\begin{tabular}{@{}ll@{}}
Version du logiciel & 2.0.0 \\
Édition du document & septembre 2026 \\
Éditeur & Centre d'Excellence Cyberdéfense Aérospatiale (CEC) \\
        & École de l'air et de l'espace, base aérienne 701, Salon-de-Provence \\
\end{tabular}

\bigskip
{\small\color{gris}Ce document décrit l'emploi de \wac{} : préparation, exécution
d'une collecte, lecture et vérification des résultats, et traces laissées sur la
machine examinée. Les mesures citées ont été relevées sur une machine virtuelle
Windows~11 de test ; elles donnent des ordres de grandeur, pas des
engagements.}
\vspace*{2cm}
\clearpage

% Numbering starts at 1 with the table of contents; only the cover page and the
% information page, which have no footer, are left out of it.
\pagenumbering{arabic}
\tableofcontents

% ============================================================ 1
\chapter{Présentation}

\section{À quoi sert WAC}

\wac{} prélève, sur un poste Windows en fonctionnement, les artefacts utiles à
une investigation numérique : base de registre, journaux d'événements, fichiers
Prefetch, raccourcis et listes de sauts, tâches planifiées, services, comptes,
processus et sessions. Il les restitue sous forme de fichiers \textbf{JSON},
prêts à être analysés ou versés dans une chaîne de traitement.

Il s'exécute depuis une clé USB, sans installation, et écrit tout ce qu'il
produit sur cette clé.

\section{Le principe directeur : une empreinte minimale}

Examiner une machine la modifie. Chaque outil lancé laisse des traces : fichiers
ouverts dont la date de dernier accès change, services sollicités qui écrivent
dans les journaux, clés de registre lues par l'API. Ces traces se mêlent à celles
que l'on cherche.

\wac{} est conçu pour que ce qu'il laisse derrière lui soit aussi proche de rien
que la tâche le permet, et pour que \textbf{ce qu'il ne peut pas éviter soit
écrit}. Concrètement :

\begin{itemize}
  \item le disque est lu \textbf{directement}, en lecture seule, sans ouvrir aucun
        fichier par le système : aucune date d'accès n'est modifiée ;
  \item aucune clé du registre de la machine n'est ouverte : les ruches sont
        copiées, puis lues sur la copie ;
  \item aucun cliché instantané (VSS), aucun appel COM ou WMI, aucun service de
        journalisation n'est sollicité ;
  \item rien n'est écrit sur le disque examiné ;
  \item chaque opération est consignée, avec la trace qu'elle est censée laisser,
        dans \fichier{investigation.json}.
\end{itemize}

\section{Ce que WAC ne fait pas}

\begin{itemize}
  \item \wac{} \textbf{n'analyse pas} : il prélève et met en forme. L'interprétation
        revient à l'analyste.
  \item Il ne réalise pas d'image complète du disque : il extrait les fichiers
        porteurs d'artefacts.
  \item Il ne capture pas la mémoire vive. Les seules données volatiles relevées
        sont la liste des processus, les sessions ouvertes et l'état des
        services.
\end{itemize}

% ============================================================ 2
\chapter{Principes de fonctionnement}

\section{Lecture directe du disque}

\wac{} ouvre le volume (\path|\\.\C:|) en lecture seule, parcourt lui-même la
table des fichiers NTFS (\path|$MFT|) et en extrait le contenu des fichiers
utiles. Le système ne voit aucune ouverture de fichier : ni date d'accès mise à
jour, ni parcours de répertoire, ni verrou.

Cette lecture prend en charge les particularités de NTFS que Windows masque
d'ordinaire :

\begin{description}[style=nextline,leftmargin=1.2em]
  \item[Compression NTFS] Windows~11 compresse notamment le dossier des journaux
        d'événements. \wac{} les décompresse lui-même.
  \item[Compact OS (WOF)] Les binaires système de Windows~10 et~11 sont stockés
        compressés dans un flux caché ; vus par l'API, ils paraissent ordinaires.
        \wac{} les décompresse (variantes XPRESS).
  \item[Attributs éclatés] Un fichier très fragmenté voit ses informations
        réparties sur plusieurs enregistrements ; \wac{} les réunit.
  \item[Fichiers préalloués] Au-delà de la longueur de données valides, NTFS
        rend des zéros, même si le disque contient encore les restes d'anciens
        fichiers. \wac{} fait de même : le contenu extrait est celui que verrait
        toute acquisition ordinaire, et son empreinte aussi.
\end{description}

\section{Consigne et travail}

Une pièce à conviction numérique ne s'examine jamais sur elle-même : on en fait
une copie, on la scelle, et l'on travaille sur une seconde copie. Si l'analyse
abîme quelque chose, la copie scellée demeure et l'opération peut être refaite.

\wac{} applique cette procédure à \textbf{tous} les fichiers extraits, y compris
ceux qu'il ne modifie pas :

\begin{description}[style=nextline,leftmargin=1.2em]
  \item[\fichier{exhibits/}] les copies brutes, telles que lues sur le volume. Une
        fois recopié dans la copie de travail, ce répertoire n'est plus touché,
        ni en écriture ni en lecture : seule la vérification de son intégrité,
        avant la recopie d'une conversion, le relit. Il contient le manifeste
        qui identifie chaque pièce, et le sceau du manifeste.
  \item[\fichier{working/}] les copies de travail, recopiées depuis la consigne
        et vérifiées par empreinte — toutes, \fichier{\$MFT}, journal des
        modifications et binaires prélevés compris, ce qui double leur place sur
        le support. C'est ici que sont faites les seules modifications (rejeu
        des journaux de transaction des ruches), et c'est ici que lisent tous
        les collecteurs.
\end{description}

\begin{note}
Chaque fichier extrait est donc écrit \textbf{deux fois} sur le support de
collecte. Prévoyez la place correspondante (voir
section~\ref{sec:place}).
\end{note}

\section{Ruches du registre et journaux de transaction}

Une ruche copiée sur une machine en fonctionnement ne contient pas les dernières
modifications : elles sont encore dans ses journaux de transaction
(\fichier{.LOG1}, \fichier{.LOG2}). \wac{} extrait ces journaux et les
\textbf{rejoue sur la copie de travail}, pour que l'analyse porte sur l'état réel
de la machine.

Tout ce que \wac{} écrit sur une copie est réversible : avant le rejeu, le contenu
d'origine de chaque page remplacée est sauvegardé dans un journal d'annulation
(\fichier{<ruche>.undo}). La copie brute, elle, reste intacte dans la consigne.

\section{L'extraction en deux temps}

L'emplacement des ruches de chaque utilisateur est inscrit dans la ruche
\fichier{SOFTWARE}. \wac{} extrait donc d'abord les ruches de la machine
(\fichier{SYSTEM}, \fichier{SOFTWARE}, \fichier{SAM}, \fichier{Amcache}), lit la
liste des profils dans la copie de \fichier{SOFTWARE}, puis extrait les ruches de
chaque profil (\fichier{ntuser.dat}, \fichier{usrClass.dat}).

Rien n'est supposé sur \fichier{C:} : le lecteur système est détecté, et un
profil situé sur un autre disque (par exemple \path|D:\Users\...|) est lu sur son
propre volume.

\section{Ce qui reste lu « à chaud »}

Quelques informations décrivent l'instant même de la collecte et ne figurent
dans aucun fichier. Elles sont lues en direct, en une seule requête chacune, et
enregistrées telles que lues dans la consigne (\fichier{exhibits\textbackslash
live}) et en même temps dans la copie de travail, où la conversion les relit :

\begin{itemize}
  \item la liste des \textbf{processus} en cours (instantané du noyau) ;
  \item les \textbf{sessions} ouvertes ;
  \item l'\textbf{état courant des services} (la configuration, elle, est lue dans
        la ruche copiée) ;
  \item l'\textbf{horloge} et l'heure de démarrage ;
  \item les \textbf{réglages} de la machine en marche : lecteur système, page de
        code, fuseau horaire, volumes et leurs numéros de série.
\end{itemize}

% ============================================================ 3
\chapter{Préparer une collecte}

\section{Prérequis}

\begin{itemize}
  \item Windows~10 ou ultérieur, 64~bits. \wac{} le vérifie au lancement.
  \item Des \textbf{droits d'administrateur} : la lecture directe d'un volume les
        exige. \wac{} le vérifie également.
  \item Des volumes \textbf{NTFS}, \textbf{FAT} ou \textbf{exFAT}. Un fichier
        situé sur un autre système de fichiers (CD-ROM) ne peut pas être lu
        directement ; il est signalé comme non lu, jamais ouvert autrement.
\end{itemize}

\section{Le support de collecte}\label{sec:place}

\wac{} s'exécute depuis le support de collecte (en général une clé USB) et y écrit
tous ses résultats. Choisissez un support \textbf{vide}, de préférence en USB~3 :
l'écriture est le poste le plus coûteux d'une collecte.

La place nécessaire dépend du mode et des options. Deux emplacements se
remplissent : la \textbf{consigne} (copies brutes scellées) et la \textbf{copie
de travail}, qui ne reprend que ce que les collecteurs modifient ou lisent
pendant la collecte (ruches, Prefetch, raccourcis, listes de sauts, binaires
cités). Les binaires prélevés par \opt{collect} \opt{binary} et les
métadonnées NTFS ne sont jamais recopiés : la conversion les lit dans la
consigne même. Mesures sur la machine virtuelle de test (Windows~11 sans
logiciel ajouté, 1~volume) :

\begin{center}
\begin{tabular}{@{}L{7.4cm}rrr@{}}
\toprule
Collecte & Consigne & Copie de travail & Durée \\
\midrule
\opt{collect} \opt{events} \opt{binary} (tous les exécutables) & 2,7~Go & 0,3~Go & 8~min~30 \\
\bottomrule
\end{tabular}
\end{center}

Les postes les plus lourds de la consigne :

\begin{center}
\begin{tabular}{@{}L{9.6cm}r@{}}
\toprule
Contenu & Consigne \\
\midrule
Métadonnées NTFS : \path|$MFT| (380~Mo), \path|$LogFile| (64~Mo),
  \path|$UsnJrnl:$J| (35~Mo, sur 3,6~Go logiques : NTFS a libéré le début,
  qui n'est pas copié) & 480~Mo \\
Binaires copiés : 2\,853 non authentifiés sur 44\,446 lus & 1,2~Go \\
Journaux d'événements (\opt{events}) & 170~Mo \\
Ruches, Prefetch, raccourcis, listes de sauts & 180 à 230~Mo \\
\bottomrule
\end{tabular}
\end{center}

Une machine utilisée depuis longtemps, ou chargée de logiciels, dépasse ces
chiffres : la \path|$MFT| croît avec le nombre de fichiers, les binaires tiers
non authentifiés avec les logiciels installés. Prévoir \textbf{au moins 5~Go
libres} pour une collecte seule avec \opt{binary}.

\wac{} vérifie l'emplacement \textbf{avant la toute première écriture} et refuse
de démarrer :
\begin{itemize}
  \item si le répertoire de travail contient déjà des fichiers — une collecte
        antérieure serait analysée à la place de celle-ci ;
  \item si la place libre est manifestement insuffisante.
\end{itemize}
Si la place vient à manquer en cours de prélèvement des binaires, ceux-ci sont
hachés sans être copiés, et le journal d'investigation l'indique.

\section{Avant d'intervenir}

\begin{itemize}
  \item Utilisez un \textbf{dossier de sortie neuf} pour chaque machine
        (option \opt{output}).
  \item Notez l'heure de l'intervention et l'identité de l'opérateur : \wac{} les
        consigne aussi, mais une trace indépendante est utile.
  \item Gardez à l'esprit que brancher la clé et lancer \wac{} laissent des traces
        inévitables (chapitre~\ref{chap:traces}) : elles sont connues et
        documentées, afin de ne pas être confondues avec l'activité étudiée.
\end{itemize}

% ============================================================ 4
\chapter{Lancer une collecte}

\section{Ligne de commande}

Ouvrez une invite de commandes \textbf{en tant qu'administrateur}, placez-vous
sur le support de collecte, puis lancez :

\begin{lstlisting}
WAC.exe [--collect | --convert=dossier] [--events] [--binary | --binary-all]
        [--threads=N] [--dump]
        [--output=dossier] [--loglevel=N] [--debug]
WAC.exe --update-trust[=dossier]
\end{lstlisting}

Toutes les options sont facultatives et désactivées par défaut.

\section{Fichier de configuration : \fichier{wac.yml}}

\textbf{Version.} \wac{} numérote ses versions \texttt{<majeur>.<mineure>.<correctif>} :
majeur pour une refonte, mineure pour un ajout de fonctionnalités, correctif
pour des corrections. La version se lit dans les propriétés de
\fichier{WAC.exe} (onglet Détails), s'affiche au lancement, et est consignée
dans \fichier{investigation.json} (\champ{Tool.Version}) : une collecte dit par
quelle version elle a été faite. Ce document décrit la version 2.0.0.

La dernière version compilée se trouve dans \fichier{bin} du dépôt :
\fichier{WAC.exe}, le \fichier{wac.yml} de référence, les scripts qui
vérifient une collecte (section~\ref{sec:verif}) et \fichier{SHA256SUMS}
pour vérifier le tout. Copier \fichier{WAC.exe} et \fichier{wac.yml} sur la
clé de collecte suffit ; les scripts restent avec l'enquêteur.

Plutôt que de retaper les options à chaque collecte, la procédure s'écrit une
fois dans \fichier{wac.yml}, à côté de \fichier{WAC.exe} : \wac{} l'applique
d'abord, et une option donnée en ligne de commande la surcharge.
\opt{write-config} écrit le fichier de référence, commenté — la procédure
recommandée : binaires non vérifiés prélevés, journaux d'événements compris.

\begin{lstlisting}
wac_version: 2.0.0      # version de WAC pour laquelle le fichier est écrit (--write-config)
binary: unverified      # none | unverified | all
threads: 0              # fils d'analyse des binaires, 0 : automatique
output: output
log_level: 0
dump: false
language: fr            # fr | en : langue de la console ; les fichiers restent en anglais
artefacts:              # une clé par artefact ; une clé absente vaut true
  activities_cache:        true   # Chronologie Windows : applications, fichiers, pages, presse-papiers (ActivitiesCache.db)
  amcache:                 true   # Amcache.hve : applications, exécutables et pilotes vus
  bam:                     true   # dernière exécution de chaque programme par utilisateur (BAM)
  bits:                    true   # qmgr.db : travaux BITS, URL, destinations, commandes de notification
  boot_records:            true   # ce qui s'exécute avant Windows : premier Mio de chaque disque (MBR, GPT), secteurs de démarrage, fichiers de la partition EFI
  cloud_identity:          true   # inscription Entra ID de la machine et comptes professionnels des utilisateurs ; caches de jetons copiés chiffrés
  defender:                true   # Defender : journaux (MPLog, traces WPP brutes), historique des détections, quarantaine (copiée chiffrée)
  defender_configuration:  true   # exclusions et réglages de protection de Microsoft Defender
  defender_endpoint:       true   # dossier local de Defender for Endpoint, là où il est inscrit (sorties de Live Response, données cyber), copie brute seule
  deleted_files:                  # fichiers supprimés dont la lecture brute retrouve encore le contenu
    enabled:               true   # false : rien n'est lu
    level:                 2      # 1 : contenu des enregistrements supprimés de la $MFT ; 2 : et noms restés dans les index
                                  # 3 : et fichiers taillés dans l'espace libre (PE, ZIP, PDF) — tout l'espace libre lu
    max_file_mib:          64     # au-delà : relevé avec ses empreintes, pas copié
    max_total_mib:         4096   # total copié ; les plus récents d'abord
    excluded_folders:             # jamais lus, sur chaque volume
      - \Windows\WinSxS
      - \Windows\SoftwareDistribution
      - \Windows\servicing
      - \Windows\Installer\$PatchCache$
      - \Windows\System32\catroot
      - \Windows\System32\catroot2
      - \$Extend\$RmMetadata
  devices:                 true   # périphériques USB et SCSI et leurs dates, portables, noms de volumes
  dpapi_master_keys:       true   # clés maîtresses DPAPI copiées chiffrées ; JSON : propriétaire, GUID, dates, jamais la clé
  emails:                  true   # Outlook (PST, OST) et Thunderbird (mbox) : messages, destinataires, dates ; corps et pièces jointes exportés
  error_reports:           true   # rapports d'erreurs (WER) : programmes plantés, avec leur chemin
  etw_logs:                true   # fichiers de trace ETW que Windows garde sur son disque (.etl : sessions de démarrage, diagnostics, Windows Update), copie brute seule
  event_log_configuration: true   # état de chaque journal et session ETW de démarrage : actif, taille, rétention, accès (hors ligne)
  events:                  true   # journaux d'événements, avec leurs messages
  explorer_recent_docs:    true   # fichiers ouverts par chaque utilisateur (RecentDocs de l'Explorateur)
  fat_volumes:             true   # volumes FAT et exFAT, clés USB du suspect comprises : entrées, supprimées aussi
  firewall:                true   # règles et profils du pare-feu
  firewall_log:            true   # journal du pare-feu (pfirewall.log), s'il est activé : connexions acceptées ou rejetées, adresses, ports, processus
  hosts:                   true   # le fichier hosts : noms détournés vers une adresse
  installed_programs:      true   # programmes installés (Uninstall) de la machine et de chaque utilisateur
  internet_settings:       true   # proxy des utilisateurs
  intune_logs:             true   # journaux de l'extension Intune : scripts et applications poussés par le locataire, leur résultat
  jump_lists:              true   # listes de sauts : fichiers et lieux ouverts, par application
  known_networks:          true   # réseaux auxquels la machine s'est connectée (NetworkList)
  local_groups:            true   # membres de chaque groupe local (SAM)
  log_file:                true   # $LogFile et $TxfLog : journaux de transactions NTFS et TxF, copie brute seule
  mft:                     true   # $MFT : chaque fichier, présent ou supprimé, et ses dates
  mounted_devices:         true   # volumes et lettres de lecteur connus de la machine (MountedDevices)
  mount_points:            true   # volumes et partages vus par chaque utilisateur (MountPoints2)
  mru:                     true   # fichiers ouverts ou enregistrés par les boîtes de dialogue communes
  mru_apps:                true   # programmes qui ont ouvert ou enregistré ces fichiers
  muicache:                true   # noms des programmes lancés par chaque utilisateur
  network:                 true   # en direct : connexions, voisins, interfaces, routes
  network_configuration:   true   # cartes réseau, dernière configuration IPv4/IPv6, baux DHCP (hors ligne)
  network_shares:          true   # partages réseau de la machine, chemin et droits
  ntfs_volumes:            true   # $Boot et $UsnJrnl:$Max : numéro de série et géométrie, création du journal
  office_trust_records:    true   # documents Office approuvés, macros activées
  onedrive:                true   # journaux (.odl) et réglages de OneDrive de chaque profil, base de synchronisation comprise, copie brute seule
  operating_system:        true   # version, installation, nom de la machine, durée de fonctionnement
  persistence:             true   # persistance du registre : IFEO, AppInit, Winlogon, LSA, COM…
  powershell_history:      true   # PSReadLine : les commandes tapées dans PowerShell
  powershell_profiles:     true   # profils PowerShell de la machine et de chaque utilisateur : une persistance
  prefetch:                true   # Prefetch : programmes lancés, compteur, huit dernières exécutions, fichiers chargés
  processes:               true   # en direct
  program_usage:           true   # programmes lancés ou utilisés (FeatureUsage, JumplistData, compatibilité, épinglés)
  rdp_client:              true   # Bureau à distance sortant : serveurs, fichiers .rdp, tuiles du cache bitmap
  rdp_server:              true   # Bureau à distance entrant : authentifications, sessions, adresses source
  recent_documents:        true   # Recent : raccourcis vers les fichiers ouverts par chaque utilisateur
  recycle_bin:             true   # $I des corbeilles : quoi, par qui, quand
  registry_backups:        true   # RegBack : anciennes ruches de la machine, copie brute seule
  root_certificates:       true   # certificats racines de la machine et des utilisateurs
  run_keys:                true   # clés Run et RunOnce de la machine et de chaque utilisateur
  scheduled_tasks:         true   # tâches planifiées : commandes, déclencheurs, auteurs, dates
  search_index:            true   # index de recherche Windows : chemin, dates, taille, résumé de chaque élément indexé
  security:                true   # ruche SECURITY : noms des secrets LSA, connexions en cache, audit
  security_descriptors:    true   # $Secure:$SDS : propriétaire et droits de chaque fichier, chaque descripteur une fois
  services:                true   # services et pilotes
  sessions:                true   # en direct
  setupapi:                true   # setupapi.dev.log : installations de périphériques
  shadow_copies:           true   # clichés instantanés lus en brut, sans le service VSS : ruches et journaux de chaque cliché
  shellbags:               true   # dossiers parcourus par chaque utilisateur dans l'Explorateur
  shimcache:               true   # cache de compatibilité des applications (AppCompatCache)
  shortcuts:               true   # raccourcis hors de Recent : menu Démarrer, Bureaux, Quick Launch, barre des tâches
  srum:                    true   # SRUDB.dat : réseau et ressources par application et utilisateur, 30 à 60 jours
  startup_folders:         true   # dossiers de démarrage : ce qui s'exécute à l'ouverture de session
  sysmon:                  true   # Sysmon s'il est installé : sa configuration (règles telles que stockées) et les fichiers supprimés qu'il a archivés ; son journal est dans events
  thumbcache:              true   # caches des miniatures et des icônes : ce que chacun a vu, même supprimé ; images exportées
  typed_paths:             true   # chemins tapés dans l'Explorateur
  usbstor:                 true   # périphériques de stockage USB (Enum\USBSTOR)
  userassist:              true   # programmes lancés depuis l'Explorateur, avec compteur et date
  users:                   true   # comptes locaux (SAM) : connexions, échecs, indicateurs, condensats
  user_shell_folders:      true   # où sont vraiment le Bureau, les Documents, les Téléchargements de chacun
  usn_journal:             true   # $UsnJrnl:$J : créations, renommages, suppressions
  web_client:              true   # navigateurs Chromium et Firefox : visites, téléchargements, recherches ; leurs bases copiées brutes
  wmi:                     true   # dépôt WMI lu hors ligne : abonnements (filtres, consommateurs, liaisons)
  word_wheel_query:        true   # recherches de l'Explorateur
  zone_identifiers:        true   # flux Zone.Identifier : d'où viennent les fichiers téléchargés
\end{lstlisting}

\begin{itemize}
  \item \textbf{Strict} : une clé inconnue, une valeur hors de celles admises
    (\texttt{true}/\texttt{false} seulement, pas \texttt{yes}), une clé en
    double, une liste : la collecte est refusée avant de commencer, avec la
    ligne fautive. Une clé absente garde sa valeur par défaut : \texttt{true}.
  \item \textbf{Versionné} : \champ{wac\_version} porte la version de \wac{}
    pour laquelle le fichier a été écrit (\opt{write-config} l'y met). Un
    fichier d'une autre version, ou sans version, est refusé avant toute
    lecture : il pourrait manquer des réglages ajoutés depuis, ou en porter qui
    n'existent plus. Le refaire avec \opt{write-config}.
  \item \textbf{Une clé par artefact}, aucun prélevé d'office. Les ruches, les
    journaux d'événements et la \path|$MFT| ne sont pas des artefacts mais des
    \emph{dépendances} : prises dès qu'un artefact demandé les lit, et
    seulement alors. Les journaux servent à \texttt{events} et à
    \texttt{rdp\_server} ; les ruches des utilisateurs à \texttt{muicache},
    \texttt{userassist}, \texttt{run\_keys}, \texttt{shellbags},
    \texttt{mru}, \texttt{mru\_apps} et aux autres artefacts lus dans
    \fichier{NTUSER.DAT} ; \fichier{Amcache.hve} au seul \texttt{amcache}.
    SYSTEM, SOFTWARE, SAM, SECURITY et DEFAULT sont toujours prises : elles
    donnent le fuseau horaire, les profils et le nom des comptes.
  \item \textbf{Consigné} : chemin, SHA-256 et contenu du fichier vont dans
    \fichier{investigation.json} (\champ{Tool.Configuration}) — la procédure
    appliquée fait partie de la preuve.
  \item Un artefact désactivé n'est pas lu du tout, et son fichier JSON le dit
    (\champ{CollectionStatus} \texttt{NotRequested}, avec la clé) : jamais
    confondu avec un échec ou une absence de traces.
  \item \opt{config=fichier} lit un autre fichier.
  \item Sans \fichier{wac.yml}, \wac{} collecte tous les artefacts, journaux
    d'événements compris, options telles que données.
\end{itemize}

\section{Options}

\begin{longtable}{@{}>{\ttfamily}L{3.2cm}L{11.2cm}@{}}
\toprule
\normalfont\bfseries Option & \bfseries Effet \\
\midrule
\endhead
-{}-collect & Le mode par défaut, dit explicitement : collecte, sur la machine
  examinée — relevés en direct et extraction brute dans la consigne scellée,
  rien n'est converti. Voir chapitre~\ref{chap:separation}. \\
-{}-convert=\newline dossier & Conversion seule, sur un poste d'analyse, de la collecte
  faite avec \opt{collect} dans ce dossier, après vérification du sceau et de
  chaque empreinte. Voir chapitre~\ref{chap:separation}. \\
-{}-events & Prend les journaux d'événements (\fichier{.evtx}, ou
  \fichier{.evt} sous Windows~XP et Server~2003) même si
  \fichier{wac.yml} dit \texttt{events: false} ; sans cela, ils sont pris par
  défaut (clé \texttt{events}). Chaque événement reçoit son message lisible,
  reconstitué à partir des ressources de son fournisseur. Environ 250~Mo de
  copies. \\
-{}-binary & À la collecte, lit directement chaque exécutable, bibliothèque,
  pilote, script et document à macros des volumes fixes, vérifie sa signature en
  mémoire et \textbf{prélève} dans la consigne ceux qu'elle n'authentifie pas ;
  à la conversion, chaque fichier cité par un artefact reçoit les empreintes de
  sa pièce. Voir section~\ref{sec:binary}. \\
-{}-binary-all & Comme \opt{binary}, mais \textbf{tout} exécutable est prélevé,
  authentifié ou non ; la signature est tout de même vérifiée, et le manifeste
  dit pour chacun si elle est valide. Pour les dossiers où le contenu même doit
  figurer dans la consigne. \\
-{}-threads=N & Nombre de fils qui analysent les exécutables de \opt{collect}
  \opt{binary} (PE, Authenticode, catalogues, empreintes). Par défaut, le nombre
  de fils du processeur moins un, la lecture du disque gardant le sien. Le
  résultat n'en dépend pas : le manifeste est le même. \\
-{}-update-trust\newline [=dossier] & Sur le \textbf{poste d'analyse}, \textbf{avant} la collecte, et
  seule : télécharge chez Microsoft la liste des racines de confiance
  (\fichier{authroot.stl}), celle des certificats désapprouvés
  (\fichier{disallowedcert.stl}) et chacune des racines, vérifie leurs
  signatures et leurs empreintes, ainsi que les listes de pilotes vulnérables
  de Microsoft et de LOLDrivers (non signées : authentifiées par HTTPS seul)
  et les listes de révocation (CRL) des autorités de signature de code que
  recense la CCADB, de celles propres à Microsoft et de celles qu'une collecte
  précédente a trouvées manquantes (\fichier{wanted-crls.txt}), et les écrit dans \fichier{dossier} (par
  défaut \fichier{trust}, à côté de \fichier{WAC.exe}), à emporter sur la clé
  avec \fichier{WAC.exe}. Seul mode qui utilise le réseau : la bibliothèque
  réseau n'est chargée que là. Refusée avec toute autre option. Sous
  Windows~7, les sites servis en TLS~1.2 seul (listes de pilotes, CCADB)
  demandent la mise à jour KB3140245 : sans elle, ces listes manquent au jeu
  de confiance, et \fichier{investigation.json} le dit. \\
-{}-dump & Ajoute la valeur hexadécimale brute des éléments de shellbags et de
  raccourcis \fichier{.lnk}, pour un décodage ultérieur. \\
-{}-output=dossier & Dossier de sortie, relatif au répertoire courant. Par défaut :
  \fichier{output}. \\
-{}-language=fr|en & Langue de la console : étapes, aide, messages d'erreur.
  Français par défaut ; \texttt{language:} dans \fichier{wac.yml} la fixe
  aussi. Les fichiers JSON, \fichier{investigation.json} et le journal texte
  restent en anglais, quelle que soit la langue : ce sont des pièces lues par
  des outils, et une collecte donne les mêmes fichiers quel que soit
  l'opérateur. \\
-{}-loglevel=N & Active un journal texte détaillé, dans le répertoire courant,
  nommé d'après le programme lancé (\fichier{WAC.exe.log} pour \fichier{WAC.exe}) : 0~aucun, 1~par type
  d'artefact, 2~par artefact, 3~par sous-fonction (mise au point). \\
-{}-debug & Affiche sur la sortie d'erreur la trace du lecteur NTFS (résolution
  des chemins, blocs d'index, séquences de grappes). Réservé au diagnostic. \\
-{}-help, /? & Affiche l'aide. \\
\bottomrule
\end{longtable}

\section{Exemples}

Collecte standard :
\begin{lstlisting}
E:\> WAC.exe --output=poste-comptable-01
\end{lstlisting}

Collecte complète, avec journaux, binaires et journal détaillé :
\begin{lstlisting}
E:\> WAC.exe --events --binary --output=poste-comptable-01 --loglevel=2
\end{lstlisting}

\section{L'option \opt{binary}}\label{sec:binary}

Une empreinte suffit pour interroger une base publique (VirusTotal, par exemple)
\textbf{sans lui envoyer le fichier}. Mais elle ne dit rien d'un binaire que
personne ne connaît — précisément celui qui intéresse l'enquête — et un binaire
non prélevé peut avoir disparu quand une détection survient. C'est pourquoi
\opt{binary} copie aussi les charges potentielles :

\begin{itemize}
  \item \textbf{prélevés} : \fichier{exe}, \fichier{dll}, \fichier{sys},
        \fichier{ocx}, \fichier{cpl}, \fichier{scr}, \fichier{drv},
        \fichier{efi}, \fichier{com}, \fichier{msi} et les scripts
        \fichier{ps1}, \fichier{psm1}, \fichier{bat}, \fichier{cmd},
        \fichier{vbs}, \fichier{vbe}, \fichier{js}, \fichier{jse},
        \fichier{wsf}, \fichier{wsh}, \fichier{hta} ;
  \item \textbf{prélevés aussi} : les documents Office \textbf{capables de porter
        des macros}, vecteur d'intrusion aussi courant qu'un exécutable —
        formats binaires anciens (\fichier{doc}, \fichier{xls}, \fichier{ppt},
        \fichier{pps}, \fichier{dot}, \fichier{xlt}, \fichier{pot}), formats
        « m » (\fichier{docm}, \fichier{dotm}, \fichier{xlsm}, \fichier{xltm},
        \fichier{pptm}, \fichier{potm}, \fichier{ppsm}), \fichier{xlsb},
        compléments (\fichier{xla}, \fichier{xlam}, \fichier{xll},
        \fichier{wll}, \fichier{ppa}, \fichier{ppam}), Publisher
        (\fichier{pub}), Visio (\fichier{vsd}, \fichier{vsdm}, \fichier{vstm},
        \fichier{vssm}) et Access (\fichier{mdb}, \fichier{accdb},
        \fichier{accde}). Ils apparaissent dans les traces comme cibles de
        raccourcis ou de listes de sauts, ou parmi les fichiers chargés par
        Word ou Excel dans leur Prefetch ;
  \item \textbf{seulement hachés} : les autres documents et fichiers de données,
        y compris \fichier{docx}, \fichier{xlsx} et \fichier{pptx}, qui ne
        peuvent pas contenir de macros. Les copier ferait de la collecte une
        copie des fichiers de l'utilisateur.
\end{itemize}

Les fichiers cités proviennent des processus, services (binaire et DLL), tâches
planifiées, fichiers chargés par les programmes (Prefetch), Shimcache, Amcache et
cibles de raccourcis et de listes de sauts. Chaque fichier n'est lu qu'une fois, quel que soit le
nombre d'artefacts qui le citent, et \textbf{un même contenu n'est stocké
qu'une fois} : les autres chemins figurent au manifeste comme pièces partageant
ce contenu (\champ{SharedExhibit}).

\subsection*{Binaires Microsoft authentiques : hachés, non prélevés}

La plupart des binaires cités sont des composants de Windows, identiques sur
toute machine de la même version. Leur origine se prouve sur la machine
elle-même, sans liste à embarquer : les \textbf{catalogues} de Windows, signés
par Microsoft, listent l'empreinte de chaque fichier du système. \wac{} lit ceux
de \path|System32\CatRoot|, ceux des composants (\path|WinSxS\Catalogs|) et
ceux des paquets de maintenance (\path|servicing\Packages|), plus de 13\,000 sur
Windows~11. Les binaires signés individuellement (Edge, OneDrive, Office, .NET)
portent une \textbf{signature intégrée}. Un binaire authentifié ainsi est
haché sans être prélevé ; tous les autres sont prélevés — y compris un binaire
système remplacé par un attaquant, dont l'empreinte ne correspond plus.

Cette vérification ne laisse \textbf{aucune trace} : catalogues et binaires sont
lus directement sur le disque, et la vérification cryptographique est faite en
mémoire par \wac{}, jusqu'à des certificats racines Microsoft embarqués dans
l'exécutable. Aucun service n'est sollicité, le magasin de certificats de la
machine n'est pas consulté, et la révocation n'est pas contrôlée en ligne.

Seuls les signataires Microsoft pour son propre code sont acceptés
(« Microsoft Windows », « Microsoft Corporation », « Microsoft Windows
Publisher » et « .NET »). Les pilotes
\textbf{tiers} certifiés par Microsoft (signataires « Hardware Compatibility
Publisher » et « Early Launch Anti-malware Publisher ») restent prélevés : un
pilote vulnérable mais signé est un moyen d'attaque classique.

\subsection*{Applications du Store : vérifiées par leur paquet}

Les fichiers d'une application du Store ne sont pas signés un à un : c'est leur
\textbf{paquet} qui l'est. Son dossier contient la signature
(\fichier{AppxSignature.p7x}) et la carte des blocs (\fichier{AppxBlockMap.xml}),
qui donne le SHA-256 de chaque bloc de 64~Kio de chaque fichier. \wac{} vérifie
la signature, qui doit remonter à une racine Microsoft (le Store signe les
paquets qu'il distribue), vérifie qu'elle garantit bien cette carte, puis compare
chaque bloc lu à la carte. Un fichier conforme est haché sans être prélevé ; un
fichier modifié, ou ajouté au dossier après l'installation, est prélevé, de même
que tout fichier d'un paquet signé hors de la chaîne Microsoft.

Les catalogues, signatures et cartes de blocs qui ont justifié un
non-prélèvement sont consignés, pour qu'un tiers puisse refaire la vérification.
Dans les fichiers d'artefacts, le champ \champ{Signature} (ou
\champ{ServiceDllSignature}, \champ{SignatureTarget}\ldots) indique
l'authentification retenue, par exemple
« \texttt{Microsoft (catalog C:\textbackslash{}Windows\textbackslash{}\ldots.cat)} »,
« \texttt{Microsoft (embedded signature)} » ou
« \texttt{Package (package C:\textbackslash{}Program Files\textbackslash{}WindowsApps\textbackslash{}\ldots)} ».

Les \textbf{scripts} sont authentifiés de la même façon : par les catalogues de
Windows, et, pour PowerShell (\fichier{ps1}, \fichier{psm1}, \fichier{psd1},
\fichier{ps1xml}\ldots), par leur bloc de signature intégré
(\texttt{\# SIG \# Begin signature block}). Un script Windows Script Host
(\fichier{vbs}, \fichier{js}, \fichier{wsf}) signé seulement dans son texte, et
non par un catalogue, reste prélevé.

Sur la machine de test, cette vérification ramène le prélèvement de 2~131
binaires (3,3~Go) à 92 (335~Mo), sans allonger la collecte ; les scripts encore
prélevés sont ceux qui ne sont pas signés.

La lecture étant directe, prélever ne laisse aucune trace de plus que hacher :
seulement de la place sur le support.

\section{Déroulement à l'écran}

\wac{} affiche chaque étape, suivie de \texttt{OK} ou d'un message d'erreur. Une
étape en échec n'interrompt pas la collecte : tout ce qui est accessible est
recueilli, et ce qui manque est consigné.

\begin{lstlisting}
[VÉRIFICATION DES PRÉREQUIS]
 - Droits d'administrateur : OK
 - Contrôle du support de collecte : OK
[OBSERVATION DE LA MACHINE EN MARCHE]
 - Sessions : OK
 - Processus : OK
 - État du réseau : OK
[EXTRACTION BRUTE]
 - Ruches du système (lecture brute NTFS) : OK
 - Ouverture de la ruche SYSTEM : OK
 - Ouverture de la ruche SOFTWARE : OK
 - Profils des utilisateurs (ruche SOFTWARE) : OK
 - Ruches des utilisateurs (lecture brute NTFS) : OK
 - Artefacts fichiers (lecture brute NTFS) : OK
[ARTEFACTS DU REGISTRE]
 - Clés UserAssist du registre : OK
   ...
[ARTEFACTS DES FICHIERS]
 - Prefetch : OK
   ...
[JOURNAUX D'ÉVÉNEMENTS]
 - Journaux d'événements : OK
[CONSIGNE]
 - Copie des binaires prélevés dans la copie de travail : OK
 - Scellement de la consigne (manifeste + SHA-256) : OK
[JOURNAL D'INVESTIGATION]
 - Écriture de investigation.json : OK
FIN, durée : 309 s
\end{lstlisting}

Avec \opt{language=en}, ou \texttt{language: en} dans \fichier{wac.yml}, les
mêmes étapes s'affichent en anglais.

Le scellement de la consigne est la \textbf{dernière} opération sur celle-ci :
toutes les pièces, y compris celles ajoutées pendant l'analyse des journaux, sont
couvertes par le manifeste.

\section{Durée}

La durée dépend surtout du matériel (USB~2 ou~3, disque, processeur) et des
options. Sur la machine virtuelle de test (disque NVMe) : environ 70~s avec
\opt{events}, environ 310~s avec \opt{events} et \opt{binary}. Sur une clé USB,
l'écriture domine.

% ============================================================ 4 bis
\chapter{Collecter ici, convertir ailleurs}\label{chap:separation}

\section{Pourquoi séparer}

Une exécution de \wac{} fait deux choses différentes :
\begin{enumerate}
  \item \textbf{la collecte} : relever ce que seul un système en marche montre
        (processus, sessions, état des services, horloge), lire le disque
        directement, copier les preuves dans une consigne scellée. Elle a
        \textbf{besoin} de la machine examinée ;
  \item \textbf{la conversion} : décoder en JSON des milliers de fichiers
        copiés (ruches, journaux, Prefetch, raccourcis\ldots). Elle n'a besoin
        que de la consigne, \textbf{pas} de la machine.
\end{enumerate}

\wac{} ne fait \textbf{jamais les deux dans la même exécution} : il collecte
sur la machine examinée, et convertit ailleurs. Convertir sur la machine y
ajouterait un travail que rien n'impose : des minutes de traitement dans la
mémoire de \wac{}, que Windows peut paginer dans \fichier{pagefile.sys}, sur le
disque même qu'on examine ; un processus présent plus longtemps, vu et
journalisé par les protections de la machine ; en cas de plantage, un rapport
d'erreur Windows écrit sur la machine ; et des bibliothèques du système chargées
pour décoder (le moteur ESE pour rejouer les journaux d'une base).

C'est précisément ce que les bonnes pratiques demandent d'éviter :
\begin{itemize}
  \item \textbf{ISO/IEC~27037} traite de l'identification, de la collecte et de
        la préservation des preuves numériques, \textbf{ISO/IEC~27042} de leur
        analyse : deux étapes distinctes ;
  \item \textbf{principe~1 de l'ACPO} : aucune action ne doit modifier les
        données de la machine examinée ; toute opération qui s'y exécute doit
        être nécessaire ;
  \item \textbf{principe~3 de l'ACPO} : un tiers indépendant doit pouvoir
        refaire le traitement et obtenir le même résultat.
\end{itemize}

Séparer les deux étapes apporte trois choses :
\begin{itemize}
  \item \textbf{moins de traces} sur la machine examinée : seule la collecte
        s'y exécute ;
  \item \textbf{des preuves vérifiées avant toute analyse} : \opt{convert}
        contrôle d'abord le sceau et le SHA-256 de chaque pièce, et refuse une
        collecte retouchée ;
  \item \textbf{une analyse reproductible} : deux conversions d'une même
        collecte donnent des fichiers JSON identiques octet pour octet, et
        l'analyse peut être refaite plus tard, avec une version corrigée de
        \wac{} ou par un tiers, sans revenir sur la machine.
\end{itemize}

\section{Les réglages de la collecte}

\begin{itemize}
  \item \textbf{Portée de \opt{binary}} : la collecte lit \textbf{tous les
        exécutables de tous les volumes fixes NTFS}, cités ou non, car elle ne
        sait pas encore lesquels les artefacts citeront. Ce temps (8~min~30 sur
        la machine de test) est de la lecture, qui n'écrit rien sur la machine.
  \item \textbf{\opt{binary-all}} : quand le contenu de chaque exécutable doit
        figurer dans la consigne (analyse de logiciel malveillant, demande des
        fichiers eux-mêmes), pour environ 20~Go.
\end{itemize}

\section{Les deux commandes}

Sur la machine examinée, en administrateur :
\begin{lstlisting}
E:\> WAC.exe --collect --events --binary --output=poste-comptable-01
\end{lstlisting}

Sur le poste d'analyse, sans droits particuliers, avec la même configuration
(\texttt{events}, \opt{binary}) que la collecte :
\begin{lstlisting}
D:\> WAC.exe --convert=E:\poste-comptable-01 --events --binary
\end{lstlisting}

Les fichiers JSON sont écrits dans le dossier de la collecte, à côté de la
consigne. \opt{convert} exclut \opt{collect} et \opt{output}.

\section{Ce que prend \opt{collect}}

La conversion ne pouvant plus rien demander à la machine, la collecte prend
\textbf{tout ce qu'elle lira} :
\begin{itemize}
  \item les \textbf{relevés en direct}, enregistrés tels que lus dans
        \fichier{exhibits\textbackslash live} et scellés comme toute pièce :
        processus, sessions, état des services, horloge et démarrage, et les
        réglages de la machine (lecteur système, page de code, fuseau horaire,
        volumes et leurs numéros de série) ;
  \item les ruches et les fichiers d'artefacts, comme d'ordinaire ;
  \item avec les journaux d'événements (\texttt{events}), les fichiers de ressources de \textbf{tous} les
        fournisseurs d'événements déclarés (binaire, fichier de paramètres,
        satellites \fichier{.mui}), car on ne sait qu'à la conversion lesquels
        seront utiles. Cela représente environ 650~Mo sur Windows~11 ;
  \item avec \opt{binary}, \textbf{tous} les exécutables, bibliothèques,
        pilotes, scripts et documents à macros de \textbf{tous les volumes fixes
        NTFS} — un profil sur \fichier{D:} compris —, lus directement et
        vérifiés en mémoire (section~\ref{sec:binary}) :
        \begin{itemize}
          \item authentifié (catalogue, signature intégrée, paquet du Store), ou
                binaire tiers dont la signature tient face au jeu de confiance de
                la clé (\opt{update-trust} : chaîne vers une racine admise par
                Microsoft pour la signature de code, valide à la date de
                signature, non révoquée, pas un pilote vulnérable) :
                \textbf{non copié} ; sans jeu de confiance utilisable, tout
                binaire tiers est copié, et le journal d'investigation le dit ; le manifeste garde son chemin, ses empreintes
                Authenticode et le verdict (\champ{ContentStored} à \texttt{false}) ;
          \item non authentifié : \textbf{copié}, avec ses trois empreintes
                MD5, SHA-1, SHA-256, son empreinte Authenticode et la raison ;
          \item fichier différentiel de Windows (\texttt{PA30}/\texttt{PA31},
                dossiers \fichier{r} et \fichier{f} de \fichier{WinSxS}) : ce n'est
                pas un exécutable ; haché, non copié.
        \end{itemize}
        Un même contenu n'est stocké qu'une fois, et les autres noms d'un lien
        physique (\fichier{System32} et \fichier{WinSxS}) sont déclarés au
        manifeste comme le premier. Les volumes amovibles, dont le support de
        collecte, et les fichiers de mémoire (\fichier{pagefile.sys},
        \fichier{swapfile.sys}, \fichier{hiberfil.sys}) sont exclus.
\end{itemize}

\begin{attention}
Avec \opt{binary}, la collecte lit tous les volumes : sur la machine
virtuelle de test (Windows~11 sans logiciel ajouté), 44\,446~exécutables lus,
dont 41\,593 authentifiés (ou tiers à signature valide) non copiés et
2\,853 copiés (1,2~Go) ; la collecte entière dure 8~minutes~30 et la consigne,
journaux d'événements et métadonnées NTFS compris, pèse 2,7~Go (voir la
section~\ref{sec:place}). Si la place vient à manquer, les fichiers
restants sont hachés sans être copiés, et le manifeste l'indique.
\end{attention}

Aucun fichier JSON d'artefact n'est écrit : le dossier ne contient que la
consigne scellée, le répertoire de travail et \fichier{investigation.json}.

\section{Ce que fait \opt{convert}}

\begin{enumerate}
  \item \textbf{Vérification}, avant toute lecture : le sceau doit porter
        l'empreinte réelle du manifeste, et chaque pièce son empreinte SHA-256.
        Une seule pièce altérée, et la conversion est \textbf{refusée} ;
        \fichier{conversion.json} en donne la raison.
  \item \textbf{Reconstruction du répertoire de travail} depuis la consigne :
        l'ancien est supprimé plutôt que cru, les copies sont vérifiées contre
        le manifeste et les ruches réparées, comme à la collecte. Les
        exécutables et les ressources des fournisseurs sont lus dans la
        consigne même, sans copie.
  \item \textbf{Conversion}, depuis les seules pièces : aucune lecture de la
        machine qui convertit, ni de son disque, ni de son registre, ni de son
        fuseau horaire.
\end{enumerate}

La consigne n'est jamais modifiée et son sceau reste celui de la collecte.
\fichier{investigation.json} reste le journal de la collecte ;
\fichier{conversion.json} est celui de la conversion (opérateur, poste,
opérations). Deux conversions d'une même collecte donnent des fichiers JSON
identiques, ce que le harnais de test vérifie à chaque cycle.

\section{Limites}

\begin{itemize}
  \item Les \textbf{documents à macros} (\fichier{doc}, \fichier{docm},
        \fichier{xls}, \fichier{xlsm}, \fichier{ppt}\ldots, liste de la
        section~\ref{sec:binary}) sont traités comme les exécutables : lus sur
        tous les volumes fixes et \textbf{prélevés}, sauf s'ils sont, octet pour
        octet, un fichier livré par Microsoft et listé dans un catalogue — ce
        qu'un document écrit ou modifié par quiconque ne peut pas être.
  \item Un fichier cité qui ne peut porter ni code ni macro (\fichier{docx},
        \fichier{xlsx}, \fichier{pdf}, donnée) n'est pas prélevé : en
        conversion, il n'a donc \textbf{pas d'empreinte}.
  \item Un exécutable authentifié n'est pas copié : à la conversion, il a son
        empreinte Authenticode et son verdict, mais pas de MD5, SHA-1 ni
        SHA-256 — ceux-ci ne sont calculés que pour ce qui est prélevé.
  \item Un exécutable situé sur un volume amovible n'est pas prélevé par
        \opt{collect} \opt{binary}.
  \item La configuration (\texttt{events}, \opt{binary}) de la conversion doit
        être celle de la collecte : ce que la collecte n'a pas pris ne peut pas être
        converti.
\end{itemize}

% ============================================================ 5
\chapter{Les résultats}

\section{Arborescence}

\begin{lstlisting}
<sortie>/
  exhibits/            copies brutes, jamais réécrites
    MANIFEST.json     identification de chaque pièce et de la collecte
    MANIFEST.sha256   empreinte du manifeste : son sceau
    Windows/...        les pièces, sous leur chemin d'origine
    Users/...
  working/             copies de travail (ruches rejouées)
  investigation.json   journal de la collecte
  *.json               les artefacts
\end{lstlisting}

Un fichier d'un autre volume que le volume système est rangé sous
\path|_volume_D\...| ; le JSON indique toujours le chemin d'origine, avec sa
vraie lettre de lecteur.

Sous \fichier{exhibits/} et \fichier{working/}, un chemin dépasse souvent
260 caractères : celui d'origine, plus le dossier de sortie. \wac{} l'écrit en
chemin étendu (\path|\\?\|). L'Explorateur et plusieurs commandes de Windows
(\texttt{Copy-Item}, \texttt{Remove-Item} de PowerShell~5.1, une copie par
glisser-déposer) s'arrêtent sur ces fichiers, parfois sans le dire. Pour copier
ou effacer une collecte, utiliser \texttt{robocopy} ou \texttt{rd /s /q
\textbackslash\textbackslash?\textbackslash<chemin>} ; pour la vérifier,
\fichier{verifier-consigne.ps1}, qui lit chaque pièce en chemin étendu.

\section{Le manifeste de consigne}

\fichier{exhibits/MANIFEST.json} identifie la collecte et chaque pièce.

\subsection{En-tête}

\begin{description}[style=nextline,leftmargin=1.2em]
  \item[\champ{Tool}] nom de l'outil, ligne de commande, date de compilation.
  \item[\champ{Host}] nom de la machine, lecteur système, fuseau horaire du
        suspect (lu dans sa ruche \fichier{SYSTEM}) et de la machine de collecte.
  \item[\champ{Operator}] compte ayant lancé la collecte, et s'il était élevé.
  \item[\champ{Custody}] déclaration de procédure, début de l'extraction, heure
        d'écriture du manifeste, nombre de pièces et d'échecs, volume total,
        algorithmes d'empreinte, volumes lus (lettre, numéro de série, système de
        fichiers).
\end{description}

\subsection{Chaque pièce}

\begin{longtable}{@{}>{\ttfamily}L{4.2cm}L{10.2cm}@{}}
\toprule
\normalfont\bfseries Champ & \bfseries Signification \\
\midrule
\endhead
SourcePath & Chemin d'origine, avec lettre de lecteur. \\
ExhibitPath & Emplacement de la copie brute, relatif au dossier de sortie. \\
Method & Méthode de collecte (lecture directe NTFS, volume, motif). \\
Result & \texttt{OK}, ou le code et le message d'erreur. Les échecs figurent au
  manifeste : une pièce absente se lirait comme une pièce jamais cherchée. \\
MD5, SHA1, SHA256 & Les trois empreintes, calculées pendant l'écriture : elles
  portent sur ce qui a été lu du volume. \\
Bytes & Taille extraite. \\
DeclaredBytes & Taille déclarée par NTFS, présente seulement si elle diffère de
  la taille extraite (extraction tronquée). \\
ValidDataBytes & Longueur de données valides d'un fichier préalloué (journaux
  d'événements) : au-delà, la pièce contient des zéros par définition. \\
SharedExhibit & \texttt{true} si le contenu, identique octet pour octet, est déjà
  stocké sous une autre pièce : \champ{ExhibitPath} désigne celle-ci. La source
  reste une pièce à part entière, avec ses propres chemin, entrée \path|$MFT| et
  horodatages. \\
MftEntry & Numéro d'enregistrement NTFS : identifie le fichier indépendamment
  de son nom. \\
Extracted(Utc) & Instant de l'extraction de cette pièce. \\
SourceCreated, SourceModified, SourceMftModified, SourceAccessed & Les quatre
  horodatages NTFS du fichier d'origine (\path|$STANDARD_INFORMATION|), en UTC et
  en heure locale du suspect. Ce sont eux, et non les dates des copies, qui
  datent le fichier. \\
SourceFileName\newline Created, Modified, MftModified, Accessed &
  Les quatre horodatages de l'attribut \path|$FILE_NAME| du nom lu. NTFS les
  tient à part, et une modification des dates par programme ne les touche pas :
  un \champ{SourceCreated} antérieur à \champ{SourceFileNameCreated} est le signe
  classique de dates falsifiées après coup (\emph{timestomping}). \\
ContentStored & \texttt{false} : binaire authentifié lu mais non copié
  (\opt{collect} \opt{binary}) ; il n'a pas de \champ{ExhibitPath}. \\
AuthenticodeSHA1, AuthenticodeSHA256 & Empreintes Authenticode d'un PE, celles
  que couvrent sa signature et les catalogues ; seules empreintes d'un binaire
  authentifié non copié. \\
SignatureVerified & Résultat de la vérification de signature d'un exécutable :
  \texttt{true} avec \champ{Signature} (ce qui l'authentifie), \texttt{false}
  avec \champ{SignatureReason} (pourquoi). \\
PeTimeDateStamp, PeSizeOfImage, SymbolServerKey & Version d'un PE et clé
  (\texttt{nom/HORODATAGEtaille/nom}) sous laquelle le serveur de symboles de
  Microsoft la conserve : un binaire Microsoft authentique non copié peut en
  être retéléchargé, identique, et son empreinte revérifiée. \\
\bottomrule
\end{longtable}

\begin{note}[Pourquoi trois empreintes]
Des collisions MD5 se fabriquent à volonté depuis 2008, et SHA-1 depuis 2017 : une
empreinte seule ne tranche plus une contestation d'identité. Trois empreintes
indépendantes, si.
\end{note}

\section{Vérifier l'intégrité d'une collecte}\label{sec:verif}

\subsection{Le sceau}

Le manifeste ne peut pas contenir sa propre empreinte : elle est dans
\fichier{MANIFEST.sha256}. Sous Linux :

\begin{lstlisting}
$ cd sortie/exhibits
$ sha256sum -c MANIFEST.sha256
MANIFEST.json: OK
\end{lstlisting}

\subsection{Chaque pièce}

Deux scripts, livrés avec \fichier{WAC.exe} dans \fichier{bin},
vérifient le sceau puis le SHA-256 de chaque pièce :
\fichier{verifier-consigne.ps1} sous Windows, \fichier{verifier-consigne.sh}
sous Linux. Leurs équivalents anglais, \fichier{verify-exhibits.ps1} et
\fichier{verify-exhibits.sh}, font les mêmes contrôles avec des messages en
anglais. Les scripts PowerShell demandent PowerShell~4 ou plus (Windows~8.1
et suivants) : sous un PowerShell plus ancien, ils le disent et s'arrêtent
(code~2) plutôt que de lire de travers.

\textbf{Vérifier depuis une machine Linux.} C'est la voie recommandée : la
collecte peut contenir ce qui a infecté la machine examinée, et une machine
d'investigation sous Linux constitue une rupture technologique — un exécutable
ou un document à macro Windows de la collecte ne peut pas s'y exécuter. Le
script ne demande que \fichier{bash}, les coreutils et \fichier{python3} (qui
lit le JSON) ; les empreintes sont calculées par \fichier{sha256sum}, non par
\wac{} :

\begin{lstlisting}
$ bash verifier-consigne.sh /media/usb/output
Sceau conforme (DCFF53FAE709EC2E351FFFD02EC6EAD9945602265CF3363D0D928B0B452C3E32)
3360 piece(s) conforme(s), 0 differente(s)
0 binaire(s) authentifie(s), empreinte seule (pas de copie a verifier)
\end{lstlisting}

Le manifeste vient de la machine examinée : les deux scripts refusent un
chemin de pièce qui sort de la collecte (\texttt{..}, chemin absolu, et pour
le \fichier{.sh} un lien symbolique), et le comptent comme différent sans le
lire.

\textbf{Depuis Windows :}

\begin{lstlisting}
PS> powershell -ExecutionPolicy Bypass -File verifier-consigne.ps1 -Sortie E:\output
Sceau conforme (DCFF53FAE709EC2E351FFFD02EC6EAD9945602265CF3363D0D928B0B452C3E32)
3360 piece(s) conforme(s), 0 differente(s)
\end{lstlisting}

Leur code de retour vaut~0 si tout est conforme, 1~sinon.

\lstinputlisting[language={},caption={}]{../../sources/scripts/verifier-consigne.sh}

\lstinputlisting[language={},caption={}]{../../sources/scripts/verifier-consigne.ps1}

\section{Le journal d'investigation}

\fichier{investigation.json} retrace la collecte : outil et ligne de commande,
machine et fuseaux horaires, opérateur, début, fin et durée, puis la liste
\textbf{ordonnée} des opérations. Pour chacune :

\begin{description}[style=nextline,leftmargin=1.2em]
  \item[\champ{Seq}, \champ{TimestampUtc}, \champ{TimestampLocal}] rang et instant ;
  \item[\champ{Operation}, \champ{Target}] ce qui a été fait, et sur quoi ;
  \item[\champ{Result}] l'issue ;
  \item[\champ{Footprint}] la trace que cette opération est censée laisser sur la
        machine examinée.
\end{description}

Ce journal permet de confronter, collecte par collecte, les traces attendues à
celles effectivement observées. Il signale aussi un écart entre le fuseau horaire
du suspect et celui de la machine de collecte.

\section{Les fichiers d'artefacts}

Un fichier par artefact. Les champs décrits sont ceux émis ; un champ sans valeur
n'est pas émis (voir chapitre~\ref{chap:lecture}).

\subsection{Système et comptes}

\begin{description}[style=nextline,leftmargin=1.2em]
  \item[\fichier{OperatingSystem.json}] nom et domaine de la machine, édition et
        version de Windows, date d'installation, heure courante, \textbf{heure de
        démarrage}, durée d'activité, fuseau horaire.
        \champ{BootTimeSource} dit d'où vient l'heure de démarrage (noyau, ou
        estimation de repli) ; \champ{ClockAdjustedSinceBootMs} donne le cumul des
        corrections d'horloge depuis le démarrage — quelques secondes sont un
        recalage ordinaire, un grand écart signale une horloge modifiée.
        \champ{ShutdownTime} : le dernier arrêt ordonné.
        \champ{NtfsLastAccessUpdated} : faux, les dates d'accès NTFS ne sont pas
        mises à jour et ne disent rien. \champ{Settings} : les réglages qui
        conditionnent la lecture des autres artefacts, tels que stockés --- dates
        d'accès et noms courts NTFS, Prefetch actif, fichier d'échange et vidages
        mémoire, jeux de contrôle, bureau à distance, chronologie Windows,
        dernier utilisateur connecté.
  \item[\fichier{security.json}] ruche \fichier{SECURITY} : \textbf{noms} des
        secrets LSA et dates de leurs valeurs (\texttt{DefaultPassword} signale
        un mot de passe de connexion automatique enregistré,
        \texttt{\_SC\_<service>} un service qui tourne sous un compte),
        nombre d'emplacements de connexions de domaine en cache et combien sont
        occupés, politique d'audit (\champ{PolAdtEv}, en octets : son format
        n'est pas publié), et les droits des comptes (\champ{AccountRights},
        lus dans \texttt{Policy\textbackslash Accounts}) : pour chaque compte ou
        groupe, ses privilèges (\champ{Privileges}, p.~ex. \texttt{SeDebugPrivilege},
        \texttt{SeBackupPrivilege}) et ses droits de connexion (\champ{LogonRights}).
        Un privilège sensible donné à un compte inattendu est une persistance.
        \textbf{Aucune valeur de secret} n'est lue, ni aucune donnée
        d'authentification en cache. Les copies
        \path|config\RegBack| des ruches sont prélevées telles quelles.
  \item[\fichier{users.json}] comptes locaux, lus dans la ruche \fichier{SAM} :
        SID et RID, état, drapeaux, nombre de connexions et d'échecs, dates de
        dernière connexion, de dernier échec, de changement de mot de passe et de
        modification du compte, profil, et les condensats NT et LM du mot de passe
        (\champ{NtHash}, \champ{LmHash}, en hexadécimal), avec leur historique
        (\champ{NtHashHistory}, \champ{LmHashHistory}). Un condensat est à sens
        unique : il désigne un mot de passe sans le révéler. Un mot de passe vide
        ou faible utilisé par un intrus se lit comme une connexion légitime ;
        \textbf{WAC ne juge rien}, il convertit, et c'est WAE qui confronte les
        condensats, hors ligne, à ceux des mots de passe faibles connus. Les
        condensats sont déchiffrés sur place (clé de démarrage de \fichier{SYSTEM},
        clé de hachage de la \fichier{SAM}, schéma AES de Windows~10 et suivants
        ou RC4 plus ancien, dont la somme de contrôle est vérifiée : une clé
        fausse ne donne aucun condensat plutôt que des condensats faux, puis
        couche DES par RID, déchiffrée par WAC lui-même pour valoir de XP à
        Windows~11). Si le compte porte des
        questions de sécurité de réinitialisation (\champ{SecurityQuestions},
        valeur \texttt{ResetData}), elles sont converties telles quelles,
        questions ET réponses : stockées en clair, une réponse triviale est une
        faiblesse que WAE juge comme un mot de passe faible.
  \item[\fichier{local\_groups.json}] groupes locaux lus dans la \fichier{SAM} :
        groupes intégrés (\champ{Scope} \texttt{Builtin}, \texttt{S-1-5-32-…}) et
        groupes créés sur la machine (\texttt{Local}), chacun avec son RID, son
        SID, son nom, sa description et ses membres directs (SID et nom résolu
        hors ligne). C'est l'appartenance qui dit qui est administrateur ou
        membre des Bureaux à distance ; un compte ajouté à un groupe privilégié
        par un intrus y figure.
  \item[\fichier{Sessions.json}] sessions ouvertes au moment de la collecte :
        utilisateur, type d'ouverture, paquet d'authentification, heure de début.
  \item[\fichier{processes.json}] processus en cours : nom, identifiants (et
        parent), propriétaire, session, modules chargés, et, avec \opt{binary},
        empreintes de l'exécutable.
  \item[\fichier{network.json}] état réseau au moment de la collecte :
        \champ{Connections} (TCP et UDP, IPv4 et IPv6 : adresses, ports, état,
        processus propriétaire et son nom, heure de création locale et UTC),
        \champ{Neighbours} (caches ARP et NDP : les machines du réseau local
        contactées récemment), \champ{Interfaces} (adresses, passerelles,
        serveurs DNS, DHCP) et \champ{Routes}. \champ{TablesNotRead} nomme une
        table que Windows a refusé de donner, ce qui n'est pas une table vide.
        \champ{TablesIpv4Only} nomme les tables lues par les fonctions de
        Windows XP et Server 2003, qui ne connaissent qu'IPv4 : l'absence de
        ligne IPv6 n'y signifie pas l'absence de connexion IPv6. Sur ces
        versions, une connexion n'a pas d'heure de création, et un voisin
        porte un \champ{Type} (entrée ARP dynamique ou statique) à la place
        d'un \champ{State}.
  \item[\fichier{services.json}] services et pilotes, configuration lue dans la
        ruche : type, démarrage, compte, binaire et DLL hébergée, commande en cas
        d'échec, dépendances, dernière écriture de la clé ; état courant ; avec
        \opt{binary}, empreintes du binaire et de la DLL.
  \item[\fichier{ScheduledTasks.json}] tâches planifiées, lues dans leurs
        définitions XML : auteur, compte d'exécution, état, dernière exécution et
        son résultat, déclencheurs et actions (avec \opt{binary}, empreintes de la
        commande). Sous Windows~XP et Server~2003, d'après leurs fichiers
        \fichier{.job} (\champ{SourceJob}) : programme, paramètres, dossier de
        travail, auteur, commentaire, dernière exécution (heure locale
        convertie), code de sortie, déclencheurs, et état du planificateur
        (\champ{JobStatus} : \texttt{HasNotRun}, \texttt{Ready}\ldots) ; le
        compte d'exécution n'est pas dans le fichier.
\end{description}

\subsection{Exécution de programmes}

\begin{description}[style=nextline,leftmargin=1.2em]
  \item[\fichier{prefetchs.json}] fichiers Prefetch : programme, nombre et dates
        d'exécution, volumes, et liste des fichiers chargés avec leur référence
        \path|$MFT| (et leurs empreintes avec \opt{binary}). Une référence
        absente n'est pas une lecture ratée : Windows l'écrit à zéro, pour tous
        les fichiers d'un programme lancé une seule fois et pour une partie des
        autres. Versions décodées : 17 (Windows~XP, qui n'enregistre aucune
        référence \path|$MFT| et une seule date d'exécution), 23 (Vista et~7,
        une seule date), 26 (Windows~8.1), 30 et 31 (Windows~10 et 11,
        compressées) ; les autres sont déclarées « format non pris en
        charge » (section~\ref{sec:format}).
  \item[\fichier{amcache\_application\_files.json},
        \fichier{amcache\_applications.json},
        \fichier{amcache\_inventory.json}] inventaire Amcache, table par
        table et valeur par valeur, avec la date de dernière écriture de chaque
        entrée : les exécutables (chemin, éditeur, version, taille, date
        d'édition de liens, \champ{AmcacheSha1}), les applications installées, et dans
        le troisième fichier les pilotes (avec leur \champ{AmcacheSha1}), périphériques,
        conteneurs, paquets de pilotes et raccourcis, chacun avec sa
        \champ{Table}. \textbf{Attention} : pour un fichier de plus de 30~Mio,
        Amcache ne hache que les 31\,457\,280 premiers octets ;
        \champ{AmcacheSha1CoversBytes} le signale, et ce SHA-1 ne se compare pas à celui
        du fichier entier. Sous Windows~8, 8.1 et les premiers Windows~10,
        l'Amcache est faite d'autres tables, \texttt{File} et
        \texttt{Programs}, dont les valeurs sont numérotées et non nommées :
        \wac{} ne nomme que celles dont le sens a été confirmé --- pour un
        fichier, \champ{LongPath}, \champ{AmcacheSha1}, \champ{ProgramId}, sa
        référence \path|$MFT| (\champ{MftReference}, lue dans le nom de sa clé)
        et le GUID de son volume ; pour un programme, \champ{Name},
        \champ{Version}, \champ{Publisher}, \champ{UninstallKeys} et
        \champ{Files} ---, et rend les autres telles que lues
        (\champ{UnnamedValues} : nom, type, valeur). Sous Windows~7, les
        exécutables lancés depuis le dernier passage de sa tâche d'inventaire,
        que liste \fichier{RecentFileCache.bcf}, sont donnés dans
        \fichier{amcache\_application\_files.json} avec \champ{Table}
        \texttt{RecentFileCache} : leur chemin seulement, le fichier ne porte
        aucune date.
  \item[\fichier{shimcache.json}] cache de compatibilité : chemin et date de
        modification du fichier, position dans le cache (\champ{CachePosition},
        0 pour la plus récente). Une entrée prouve que le binaire était présent,
        pas qu'il a été exécuté. L'indicateur d'exécution que d'autres outils
        publient n'est \emph{pas} relevé : mesuré contre le Prefetch, son « oui »
        est presque toujours juste, mais son « non » est faux pour 9 binaires
        exécutés sur 10 — il ferait conclure qu'un programme n'a jamais tourné.
        Formats décodés : Windows~XP, Vista, 7, 8, 8.1, 10 et 11 ; sous Windows~8.1, le chemin
        commence par \texttt{SYSVOL\textbackslash}, le volume système, et il est
        publié tel quel. Windows~XP y note aussi la taille du fichier
        (\champ{FileSize}) et le moment où il a mis l'entrée à jour
        (\champ{LastUpdateUtc}). Server~2003 écrit la même signature que Vista
        (\texttt{0xBADC0FFE}) avec un autre contenu : la version examinée les
        départage, et la forme de Server~2003, pas encore mesurée, est
        déclarée non prise en charge.
  \item[\fichier{bams.json}] dernière exécution par programme et par utilisateur
        (Background Activity Moderator). Les ruches étant lues sur le disque,
        une exécution de la ou des dernières minutes peut manquer : Windows
        garde une modification du registre en mémoire avant de l'écrire (de
        56 à 82~s mesurées sur la machine de test, que \texttt{RegFlushKey}
        n'a pas avancées), et \wac{} ne force pas cette écriture, qui
        modifierait le disque examiné.
  \item[\fichier{userassists.json}] programmes lancés depuis l'interface :
        compteur, nombre et durée de premier plan (\champ{FocusTimeMs}),
        dernière exécution. Sous Windows~XP, la valeur ne garde pas le premier
        plan, et son compteur part de~5 (mesuré : deux lancements donnent~7) :
        \champ{Count} est le compteur moins~5.
  \item[\fichier{muicache.json}] noms d'affichage des programmes exécutés
        (sous Windows~XP, dans \texttt{ShellNoRoam\textbackslash MUICache} de
        NTUSER.DAT).
  \item[\fichier{run.json}] programmes lancés au démarrage ou à l'ouverture de
        session (clés Run et RunOnce), avec la dernière écriture de la clé : dans
        SOFTWARE (toute la machine), dans la ruche DEFAULT (le compte système) et
        dans le NTUSER.DAT de chaque profil (\champ{Hive}). \champ{Type}
        \texttt{REG\_EXPAND\_SZ} : la commande contient des variables
        (\texttt{\%SystemRoot\%}\ldots), données telles qu'écrites. Aussi :
        les clés Run des programmes 32~bits (\texttt{SOFTWARE (32-bit)}) ;
        \texttt{StartupApproved} (\champ{Key} \texttt{StartupApproved\textbackslash Run},
        \texttt{Run32}, \texttt{StartupFolder}), l'état d'une entrée tel que
        l'utilisateur l'a réglé dans le Gestionnaire des tâches ou les Paramètres :
        \champ{Status}, les 4 premiers octets tels que stockés (leur sens n'est pas
        mesuré, donc pas interprété), et \champ{Changed}, l'instant du
        changement ; \texttt{RunNotification}, une entrée de démarrage que Windows
        a remarquée (\champ{Notification}). Un nom sans valeur Run correspondante
        est une entrée retirée depuis.
  \item[\fichier{persistence.json}] persistance du registre au-delà des clés
        Run : IFEO (\texttt{Debugger}, \texttt{GlobalFlag}), \texttt{SilentProcessExit},
        AppInit\_DLLs et AppCertDlls, Winlogon (\texttt{Shell},
        \texttt{Userinit}, \texttt{Taskman}, \texttt{Notify}), paquets LSA,
        moniteurs d'impression, \texttt{BootExecute}, Active Setup, assistants
        NetSh, commandes d'ouverture des types exécutables (\texttt{exefile},
        \texttt{batfile}, \texttt{cmdfile}\ldots, machine et utilisateur), et
        détournements COM des classes de chaque utilisateur (UsrClass.dat). Pour chaque valeur : mécanisme (\champ{Mechanism}),
        ruche, clé, nom, type, valeur telle qu'écrite, dernière écriture de la
        clé, et avec \opt{binary} les empreintes du fichier visé. Toutes les
        valeurs sont données, celles de Windows comprises : distinguer
        l'ajouté du défaut revient à l'analyse. Les tâches BITS ne sont pas
        encore lues.
  \item[\fichier{office\_trust\_records.json}] documents Office qu'un utilisateur
        a approuvés (« Activer le contenu ») : application, version, document,
        date ; \champ{MacrosEnabled} quand il a activé les macros — la porte de
        la plupart des infections par macro.
  \item[\fichier{typed\_paths.json}, \fichier{word\_wheel\_query.json}] chemins
        tapés dans la barre d'adresse de l'Explorateur et termes qui y ont été
        cherchés, par utilisateur et par position (\champ{MruPosition}, 0 = le plus
        récent, comme dans \fichier{shellbags.json}).
  \item[\fichier{root\_certificates.json}] certificats des magasins racines de la
        machine et de chaque utilisateur : sujet, émetteur, validité ; une
        racine ajoutée fait accepter tout ce qu'elle signe (interception TLS,
        logiciel signé). \champ{ThumbprintMismatch} signale un certificat dont le
        SHA-1 n'est pas le nom de sa clé.
  \item[\fichier{internet\_settings.json}] proxy et script de configuration
        automatique de chaque utilisateur.
  \item[\fichier{intune\_logs.json}] journaux de l'extension de gestion
        Intune (\fichier{ProgramData\textbackslash Microsoft\textbackslash IntuneManagementExtension\textbackslash Logs}),
        sur une machine gérée par Microsoft Intune : les scripts PowerShell et
        les applications que le locataire a poussés et exécutés en SYSTEM,
        leur résultat, l'inscription de l'appareil. Une fiche par message
        (format CMTrace, sur une ou plusieurs lignes) : \champ{Message},
        \champ{Component}, \champ{Type}, \champ{Thread}, la date, écrite en
        heure locale, en UTC et en heure locale — exacte quand le journal
        écrit son décalage (\champ{BiasMinutes}).
  \item[\fichier{cloud\_identity.json}] inscription de la machine à Entra ID
        (\texttt{Control\textbackslash CloudDomainJoin} de SYSTEM :
        \champ{TenantId}, \champ{UserEmail}, \champ{IdpDomain}, l'empreinte du
        certificat de l'appareil en \champ{KeyName}, et les informations du
        locataire, URL d'inscription MDM comprises) et comptes professionnels
        ajoutés par chaque utilisateur (\texttt{WorkplaceJoin} de son
        NTUSER.DAT) ; toutes les valeurs telles que stockées,
        \champ{KeyLastWrite} datant l'inscription. Les caches de jetons des
        comptes (\fichier{TokenBroker}, \fichier{Microsoft.AAD.BrokerPlugin})
        sont prélevés bruts, chiffrés par DPAPI comme Windows les garde, et
        jamais déchiffrés par \wac{}. Prélevés bruts aussi, sans conversion :
        les fichiers de trace ETW que Windows écrit sur son disque
        (\opt{etw\_logs}), les journaux et réglages de OneDrive
        (\opt{onedrive}, sauf ceux de la synchronisation des listes
        SharePoint) et le dossier local de Defender for Endpoint
        (\opt{defender\_endpoint}, sans les binaires du capteur).
  \item[\fichier{firewall.json}] règles du pare-feu (nom, action, sens, ports,
        programme ; tous les champs dans \champ{Fields}) et réglages des profils
        domaine, privé et public. Sous Windows~XP SP2 et suivants et
        Server~2003 SP1, les exceptions de chaque profil --- programmes
        autorisés et ports ouverts --- en règles de \champ{Version}
        \texttt{XP}, avec leur \champ{Profile} et leur \champ{Scope} ; XP~SP1
        n'a pas de politique de pare-feu.
  \item[\fichier{firewall\_log.json}] journal du pare-feu
        (\fichier{pfirewall.log} et sa rotation \fichier{.old}), là où le
        \texttt{LogFilePath} de chaque profil le place — réglages du pare-feu et
        stratégie de groupe — et à son emplacement par défaut : une fiche par
        connexion acceptée ou rejetée, avec \champ{Action}, \champ{Protocol},
        \champ{SourceAddress}, \champ{DestinationAddress}, les ports,
        \champ{Size}, \champ{Direction} et, depuis Windows~10,
        \champ{ProcessId} ; la date, écrite en heure locale, donnée en UTC et
        en heure locale. C'est la seule trace des flux réseau qu'un Windows
        garde sur son disque ; la journalisation est désactivée par défaut.
  \item[\fichier{installed\_programs.json}] programmes installés, lus dans les
        clés \texttt{Uninstall} de la machine (64 et 32~bits) et de chaque
        utilisateur, et produits Windows Installer (\champ{Scope} \texttt{MSI
        product}, avec le paquet mis en cache \champ{LocalPackage}) : toutes leurs valeurs (nom, éditeur,
        version, dossier, commande de désinstallation, \champ{InstallDate} --- le
        jour seul, tel qu'écrit) et la date d'écriture de la clé. Un outil
        d'accès à distance installé, puis désinstallé, y laisse souvent sa trace.
  \item[\fichier{defender\_configuration.json}] configuration de Microsoft
        Defender lue dans le registre : les \textbf{exclusions} (chemins,
        extensions, processus, adresses ; \champ{Source} \texttt{Settings} ou
        \texttt{Policy}), qu'un intrus ajoute pour que ses outils ne soient pas
        analysés, et les valeurs qui règlent ou coupent la protection (temps
        réel, protection contre la falsification, envoi d'échantillons\ldots).
  \item[\fichier{explorer\_recent\_docs.json}] fichiers ouverts par chaque
        utilisateur, tels que l'Explorateur les retient (\texttt{RecentDocs}) :
        par extension (\champ{Extension}) et pour toutes extensions confondues,
        avec leur rang (\champ{MruPosition}, 0 pour le plus récent), l'élément de
        shell du raccourci créé dans Recent, et la date d'écriture de la liste,
        qui est celle de l'ouverture du fichier en position~0. Le nom survit au
        fichier et au raccourci.
  \item[\fichier{mount\_points.json}] volumes et partages réseau qu'a vus
        l'Explorateur de chaque utilisateur (\texttt{MountPoints2}) : volume
        (GUID, à rapprocher de \fichier{mounted\_device.json} pour une clé USB)
        ou partage (\champ{Share}, \texttt{\textbackslash\textbackslash serveur\textbackslash partage}),
        date d'écriture de l'entrée --- la dernière fois que l'Explorateur l'a
        vue ---, ses valeurs, et ses sous-clés avec leur commande
        (\texttt{Shell\textbackslash AutoRun\textbackslash command} : une
        persistance ancienne).
  \item[\fichier{known\_networks.json}] réseaux auxquels la machine s'est
        connectée (\texttt{NetworkList}) : nom, type, catégorie, adresse
        matérielle de la passerelle et suffixe DNS de chaque signature, première
        et dernière connexion. Ces deux dates sont écrites par Windows en heure
        locale : \wac{} les donne en UTC avec le décalage en vigueur à leur
        date (mesuré contre le journal \texttt{NetworkProfile}). Elles disent où
        la machine a été.
  \item[\fichier{program\_usage.json}] ce que chaque utilisateur a lancé ou
        utilisé, source par source (\champ{Source}) : \texttt{FeatureUsage}
        (lancements, bascules, clics sur la barre des tâches, par application,
        \champ{Count}), \texttt{JumplistData} (dernière utilisation,
        \champ{LastUsed}), le magasin de l'Assistant de compatibilité des
        programmes (programmes exécutés), les modes de compatibilité
        (\texttt{Layers}), les programmes diagnostiqués par RADAR, les éléments
        épinglés à la barre des tâches (\texttt{Taskband}), chacun décodé comme
        un raccourci complet, et les programmes choisis pour ouvrir chaque type
        de fichier (\texttt{FileExts} : \texttt{OpenWithList} avec
        \champ{Extension} et \champ{MruPosition}, 0 pour le plus récent ;
        \texttt{UserChoice}, le programme par défaut). \texttt{OpenWithProgids}
        n'est pas relevé : les installeurs le remplissent, l'utilisateur n'y
        est pour rien.
  \item[\fichier{devices.json}] périphériques que la machine a connus
        (\champ{Source}) : USB et SCSI (\texttt{Enum}), avec leurs valeurs et les
        dates d'installation, de première installation, de dernière arrivée et
        de dernier retrait (\champ{Install}, \champ{FirstInstall},
        \champ{LastArrival}, \champ{LastRemoval}, en UTC) ; périphériques
        portables (téléphone, appareil photo) ; noms de volume que la recherche
        Windows a retenus par lettre de lecteur (\texttt{VolumeInfoCache}) ;
        interfaces de disque, de volume et de périphérique portable
        (\texttt{DeviceClasses}), écrites à leur première apparition.
  \item[\fichier{network\_shares.json}] dossiers que la machine partage
        (\texttt{LanmanServer\textbackslash Shares}) : nom, chemin, réglages tels
        que stockés, et descripteur de sécurité décodé --- propriétaire, et pour
        chaque entrée de la liste de contrôle d'accès son type, son masque, le
        SID et son nom. Les partages administratifs (\texttt{C\$},
        \texttt{ADMIN\$}) n'y figurent pas : Windows ne les stocke pas.
  \item[\fichier{network\_configuration.json}] configuration réseau lue dans
        les ruches : nom d'hôte et domaine (\champ{Global}), et pour chaque carte
        sa description et son pilote, une adresse matérielle forcée
        (\texttt{NetworkAddress}), son adresse matérielle courante et
        d'origine, sa dernière configuration IPv4 et IPv6, et son bail DHCP
        (obtention et fin, en UTC). Elle reste lisible quand la machine est
        débranchée ou le bail expiré, ce que le relevé vivant ne montre plus.
  \item[\fichier{event\_log\_configuration.json}] configuration de chaque
        journal d'événements lue dans les ruches : les canaux
        (\texttt{WINEVT\textbackslash Channels} : \champ{Enabled},
        \champ{Type}, \champ{OwningPublisher}, droits d'accès
        \champ{ChannelAccess}, taille et rétention si elles sont fixées) et les
        journaux classiques (\texttt{Services\textbackslash EventLog} :
        fichier, \champ{MaxSize}, \champ{Retention}, accès) et les sessions
        ETW lancées au démarrage (\texttt{WMI\textbackslash Autologger} :
        \champ{Start}, fichier, taille). Un intrus qui coupe un canal, réduit
        un journal, en ouvre l'accès ou arrête une session comme
        \texttt{EventLog-Security} y laisse sa trace ; \champ{KeyLastWrite} date la dernière modification de chacun.
  \item[\fichier{boot\_records.json}] ce qui s'exécute avant Windows, lu en
        brut sur chaque disque physique (le disque du support de collecte
        excepté, s'il ne porte pas le système) : son premier mébioctet copié — code d'amorçage du MBR, table
        de partitions, en-tête et entrées GPT, et l'espace avant la première
        partition, où se cache le code d'un bootkit —, les 16 premiers
        secteurs de chaque partition, et chaque fichier de la partition
        système EFI, lue par le lecteur FAT de \wac{} (chargeurs, BCD). Une
        fiche par disque : \champ{PartitionStyle} (\texttt{MBR},
        \texttt{GPT}, ou \texttt{Volume} pour un disque formaté en entier,
        sans table, comme une clé USB : son premier secteur est celui du
        volume, \champ{BootSectorOemId}), \champ{DiskId}, \champ{MbrBootCodeSha256}, et chaque
        partition en secteurs comme la table la stocke (\champ{FirstSector},
        \champ{Sectors}), avec son type, le SHA-256 de ses premiers secteurs,
        l'OEM de son secteur de démarrage et, pour la partition EFI, ses
        fichiers (\champ{EfiFiles} : chemin, taille, SHA-256). Les partitions
        logiques d'une partition étendue ne sont pas suivies.
  \item[\fichier{fat\_volumes.json}] volumes FAT12, FAT16, FAT32 et exFAT, clés
        USB et cartes du suspect comprises (le volume de la collecte excepté) :
        chaque entrée, présente ou \textbf{supprimée} (\champ{Deleted}), son
        chemin, son nom 8.3, sa taille, ses attributs et ses dates. Un nom
        supprimé en FAT garde son nom long s'il n'a pas été écrasé ; sinon son
        nom 8.3, première lettre perdue (\champ{NameIncomplete}). Chaque date
        est donnée en UTC et en heure locale dérivée de l'UTC, comme dans tous
        les artefacts : un même instant se lit partout à la même heure. FAT
        stocke l'heure locale calculée par Windows au moment où il l'écrit : une
        date de modification antérieure à la création est une date conservée
        par une copie, écrite au moment de la copie, et convertie avec le
        décalage de ce moment (\champ{ModifiedWrittenAt} \texttt{Copy},
        \champ{ModifiedOffsetApplied}, \champ{ModifiedAsStored}) — sans quoi un
        fichier d'hiver copié en été serait décalé d'une heure. exFAT garde
        son décalage (\champ{ModifiedOffset}\ldots), la date UTC est exacte. La
        collecte ne copie que les métadonnées du volume (secteur de démarrage,
        FAT, répertoires), pas le contenu des fichiers ; avec \opt{binary}, un
        programme de ces volumes cité par un artefact (lancé depuis une clé,
        par exemple) est empreinté et prélevé comme sur NTFS.
  \item[\fichier{shadow\_copies.json}] clichés instantanés (« versions
        précédentes », points de restauration) de chaque volume NTFS fixe,
        \textbf{lus en brut} — catalogue et magasins de la zone de
        différences, sans le service VSS ni son interface : WAC ne crée, ne
        monte et n'expose aucun cliché. Un objet par cliché, du plus ancien au
        plus récent : \champ{Volume}, \champ{ShadowCopyId} (l'identifiant
        que montre \texttt{vssadmin}), \champ{ShadowCopySetId},
        \champ{StoreId}, \champ{CreationTimeUtc}, \champ{Context},
        \champ{Attributes}, \champ{OperatingMachine}, \champ{StoredBlocks}
        (blocs de 16~Kio dont le magasin garde l'état d'avant) et
        \champ{ExhibitFolder}, le dossier de la consigne où sont rangés ses
        fichiers (\path|vss\<instant>_<identifiant>|). Pour chaque cliché, la
        collecte relit à travers lui les ruches (système, Amcache, profils de
        \path|\Users|) et les journaux \texttt{.evtx} ; un fichier identique
        (SHA-256) à sa version plus récente — sur le volume tel qu'il est, ou
        dans le cliché suivant — n'est pas recopié : le manifeste l'inscrit
        comme partagé avec elle. Chaque pièce lue dans un cliché porte
        \champ{ShadowCopyId} au manifeste. Une ruche effacée ou modifiée par
        un intrus, un journal vidé, s'y retrouvent dans leur état d'avant.
        Windows continue d'écrire pendant la lecture : chaque fichier est lu
        deux fois, les listes de blocs relues avant chaque lecture, et refait
        si les deux lectures diffèrent. Windows~XP ne garde pas de cliché.
        Un cliché tient ce qui était \emph{sur le disque} à son instant, et
        Windows écrit le registre et les journaux d'événements en retard : une
        valeur ou un événement écrits moins d'une minute environ avant le
        cliché peuvent ne pas y être — Windows lui-même ne les y lit pas
        (mesuré : absents 15~s après, présents 90~s après).
  \item[\fichier{shadow\_copy\_events.json}] les événements que les
        journaux d'un cliché gardent et que les journaux actuels n'ont plus
        (journal vidé, enregistrements écrasés par la rotation) : même forme
        que \fichier{events.json}, plus \champ{ShadowCopyId} et
        \champ{ShadowCopyCreationTimeUtc}, le cliché le plus récent qui le
        garde. Chaque événement une fois, quel que soit le nombre de clichés.
  \item[\fichier{shadow\_copy\_registry.json}] les entrées des artefacts de
        registre (clés Run, services, comptes et groupes, ShimCache, Amcache,
        BAM, UserAssist, MUICache, ShellBags, MRU, périphériques USB,
        programmes installés, réseaux connus…) convertis depuis les ruches
        d'un cliché, que la conversion du volume ne donne pas : une valeur
        Run, un service, un compte supprimés depuis, la trace d'un outil
        effacé. \champ{Artefact} (le fichier JSON de l'artefact),
        \champ{ShadowCopyId}, \champ{ShadowCopyCreationTimeUtc} et
        \champ{Entry}, l'entrée telle que l'artefact la donne. Chaque entrée
        une fois. Une entrée modifiée depuis y figure aussi, dans son état
        d'alors. Le cliché entier n'est pas converti : il répéterait, presque
        identiques, les milliers d'entrées du volume.
  \item[\fichier{srum.json}] usage des ressources (SRUM, \path|sru\SRUDB.dat|),
        une entrée par ligne de chaque table : octets envoyés et reçus par
        application et par interface (une exfiltration s'y lit, avec le
        programme), cycles processeur, octets lus et écrits, connexions ; 30 à
        60 jours. \champ{Table} (GUID du fournisseur), \champ{Provider} (nom
        documenté : \texttt{NetworkUsage}, \texttt{ApplicationResourceUsage}\ldots),
        puis les colonnes par leur nom (lettres et chiffres seuls), et
        \champ{App}, \champ{UserSID}, \champ{UserName} résolus par
        \texttt{SruDbIdMapTable}. La base est copiée vivante : la dernière
        heure est dans ses journaux \path|SRU*.log|, copiés dans la consigne pour
        un rejeu (\texttt{esentutl /r}) sur une copie.
  \item[\fichier{bits\_jobs.json}] travaux BITS (\path|qmgr.db|) : nom,
        description, type, état, propriétaire, \textbf{commande de
        notification} (lancée à la fin ou à l'échec du transfert : une
        persistance), délais, dates, et chaque fichier : URL source,
        destination, fichier temporaire, tailles. Le blob complet est donné en
        hexadécimal : BITS~10 y garde plus (le programme créateur) à des
        endroits non décrits. Avant Windows~10, les travaux sont dans les
        fichiers de file d'attente \fichier{qmgr0.dat} et
        \fichier{qmgr1.dat}, que BITS écrit tour à tour : chaque travail dit
        dans lequel il a été lu (\champ{QueueFile}), car le plus ancien peut
        garder un état antérieur du même travail. Sous Vista (BITS~3.0), le
        séparateur des travaux n'est connu d'aucun outil publié ; il est
        mesuré et lu de même.
  \item[\fichier{wmi\_persistence.json}] abonnements aux événements WMI, lus
        \textbf{hors ligne} dans le dépôt CIM copié en brut
        (\path|wbem\Repository|) : aucune requête WMI, rien exécuté sur la
        machine. Dans \emph{chaque} espace de noms, chaque filtre (requête
        WQL), consommateur (ligne de commande, script VBScript ou JScript,
        journal, courriel, et toute classe qu'un intrus aurait définie),
        liaison et minuterie, avec toutes ses propriétés (\champ{Properties} :
        nom et valeur ; \champ{CreatorSID} en clair). C'est une persistance
        invisible au registre. Les deux dates d'en-tête d'une instance ne sont
        pas émises : ce ne sont pas des dates de création. Sous Windows XP,
        le dépôt est dans \path|wbem\Repository\FS| et suit une autre
        disposition (condensats MD5, entrées de page de 4 octets, deux
        fichiers de correspondance au SP1) : il est lu de la même façon.
  \item[\fichier{web\_client.json}] navigation de chaque utilisateur dans les
        navigateurs Chromium (Edge, Chrome, Brave, Vivaldi, Chromium, Opera),
        lue dans la base \fichier{History} de chaque profil de navigateur. Un
        enregistrement par URL (\champ{Type} \texttt{Url} : titre, nombre de
        visites et de saisies, dernière visite), par visite (\texttt{Visit} :
        instant, \champ{Transition} — \texttt{typed} saisie, \texttt{link},
        \texttt{reload}\ldots{} — et ses qualificatifs, visite d'origine
        \champ{FromUrl}, durée), par téléchargement (\texttt{Download} :
        chaîne complète des redirections \champ{UrlChain}, page et onglet
        d'origine, chemin de destination, octets reçus, état, ouverture,
        \champ{DangerType} brut, SHA-256 quand le navigateur l'a calculé) et par
        terme cherché dans la barre d'adresse (\texttt{SearchTerm}). Chaque
        enregistrement porte \champ{Browser}, \champ{BrowserProfile},
        \champ{SID} et \champ{Source}. Les dates de Chromium (microsecondes
        depuis 1601, UTC) sont converties exactement. Les bases sont lues sur une
        copie privée avec leur journal, rejoué comme le ferait le navigateur. À
        côté de \fichier{History}, la consigne reçoit brutes les autres bases
        du profil (\fichier{Login Data}, \fichier{Web Data}, cookies,
        raccourcis, signets, préférences) et \fichier{Local State}, dont la clé
        protégée par DPAPI chiffre les mots de passe enregistrés ; jamais les
        caches. Firefox (\champ{Browser} \texttt{Firefox}, chaque dossier de
        \path|Profiles|) donne les mêmes enregistrements depuis
        \fichier{places.sqlite} — visites avec leur type (\texttt{typed},
        \texttt{link}, \texttt{download}\ldots), téléchargements avec leur
        destination, leur fin et leur taille —, plus \texttt{TypedInput} (ce
        qui a été tapé dans la barre d'adresse et la page atteinte) et, depuis
        \fichier{formhistory.sqlite}, \texttt{FormEntry} : chaque valeur
        saisie dans un formulaire ou la barre de recherche
        (\champ{FieldName} \texttt{searchbar-history}), avec ses dates. Sont
        copiés bruts à côté : cookies, \fichier{logins.json} et
        \fichier{key4.db} (identifiants enregistrés, chiffrés), extensions,
        préférences, session. Enfin \texttt{TypedUrl} : chaque adresse tapée
        dans Internet Explorer ou la barre d'adresse de l'Explorateur
        (\path|TypedURLs| du \fichier{NTUSER.DAT}, \champ{MruPosition} 0 pour
        la plus récente) avec l'instant de sa saisie (\path|TypedURLsTime|,
        \champ{TypedUtc}). \texttt{SavedLogin} : chaque identifiant enregistré
        (\fichier{Login Data} de Chromium, \fichier{logins.json} de Firefox) —
        site, compte (en clair chez Chromium, chiffré et donc absent chez
        Firefox), création, dernier usage, nombre d'usages, refus
        d'enregistrer (\champ{NeverSave}) ; \texttt{Cookie} : site, nom,
        chemin, dates, drapeaux. Les mots de passe et les valeurs des cookies,
        chiffrés, ne sont jamais décodés : les bases sont dans la consigne pour
        une analyse qui détiendrait la clé DPAPI de l'utilisateur. Ce qui
        était \emph{ouvert} : \texttt{SessionEntry}, chaque page de chaque onglet
        des fichiers de session --- \path|Sessions\Session_*| et
        \path|Tabs_*| de Chromium, \fichier{sessionstore.jsonlz4} et ses
        sauvegardes chez Firefox ---, avec \champ{State} (\texttt{Open},
        \texttt{Closed}, \texttt{ClosedTab}, \texttt{ClosedWindow}), l'instant
        où la page a été affichée (\champ{NavigatedUtc}, Chromium, celui de sa
        visite dans \fichier{History}) ou le dernier accès à l'onglet
        (\champ{TabLastAccessedUtc}, Firefox), l'instant de la fermeture
        (\champ{ClosedUtc}) et la page affichée (\champ{Shown}) : un onglet
        d'une fenêtre privée, ou d'un historique effacé, peut n'être plus que
        là ; \texttt{Favicon}, l'icône de chaque page visitée, conservée après
        l'effacement de l'historique ; \texttt{AddressBarPrediction}
        (Chromium), ce qui a été tapé dans la barre d'adresse et l'adresse
        alors choisie. Les sauvegardes de favoris de Firefox, le dossier
        local d'Internet Explorer et les conteneurs de WinINet
        (\path|INetCookies|, \path|IEDownloadHistory|) sont copiés bruts. Et
        \texttt{CacheEntry} (\champ{Browser}
        \texttt{WebCache}) : chaque entrée de la base
        \fichier{WebCacheV01.dat} d'Internet Explorer, d'Edge Legacy et des
        applications WinINet, par conteneur (\champ{Container} :
        \texttt{History}, \texttt{MSHist\ldots}, \texttt{Content},
        \texttt{iedownload}, \texttt{BackgroundTransferApi}) — URL, compte et
        page d'une entrée d'historique (\champ{User}, \champ{Page}), nombre
        d'accès, dates, fichier du cache (\champ{CachedFile} : son chemin, par
        le sous-dossier que l'entrée désigne). Le moteur d'Internet Explorer
        (WinINet) reste utilisé par \texttt{mshta.exe} et des scripts même
        quand IE est redirigé vers Edge : le cache de chaque profil
        (\path|INetCache\IE|) est collecté, limité aux types utiles à
        l'enquête — HTA, scripts, exécutables, archives, documents à macros,
        PDF —, et l'HTA qu'une attaque a fait exécuter s'y retrouve avec son
        entrée. Un téléchargement d'Internet Explorer 10 et 11 (conteneur
        \texttt{iedownload}) donne en outre \champ{Download} : son URL, son
        type MIME, sa copie temporaire dans le cache, le chemin où il a été
        enregistré (\champ{TargetPath}, ou seulement \champ{SuggestedName}
        s'il ne l'a pas été), sa taille et l'adresse du serveur — format non
        documenté par Microsoft, établi sur un vrai téléchargement d'IE~11
        dans les machines de test Windows~11 et Windows~8.1, dont les
        dispositions diffèrent ; une autre disposition n'est pas devinée. Sur Windows 7 et antérieurs, Internet Explorer 5 à 9 tenait ces
        mêmes données dans des \fichier{index.dat} (cache, historique et ses
        dossiers \texttt{MSHist01}, cookies, téléchargements) : ils sont
        collectés et convertis en \texttt{CacheEntry} (\champ{Browser}
        \texttt{Internet Explorer}), entrées récupérées
        (\champ{Recovered}) et abandonnées (\champ{RecordKind}
        \texttt{LEAK}) comprises ; l'heure de visite d'un historique
        hebdomadaire, locale, est convertie, l'expiration et la dernière
        vérification, dates FAT sans fuseau, sont données telles qu'écrites.
        La date de modification des
        conteneurs \texttt{MSHist}, que Windows écrit en heure locale, est
        convertie en UTC avec le décalage en vigueur à sa date. La base est
        copiée vivante avec ses journaux : ses dernières entrées peuvent n'être
        que dans ceux-ci (rejeu \texttt{esentutl /r} par l'analyste).
  \item[\fichier{dpapi\_master\_keys.json}] clés maîtresses DPAPI du système
        (\path|System32\Microsoft\Protect|, \champ{Scope} \texttt{System}) et
        de chaque profil (\path|AppData\Roaming\Microsoft\Protect|,
        \texttt{User}). Les fichiers sont copiés tels quels dans la consigne,
        toujours chiffrés par le mot de passe de l'utilisateur ou par le secret
        DPAPI du système, avec \fichier{Preferred} et \fichier{CREDHIST} :
        avec ce mot de passe, l'analyse peut déchiffrer plus tard les secrets
        enregistrés (mots de passe des navigateurs, Wi-Fi, gestionnaire
        d'identifiants). Le JSON n'en donne que les métadonnées, jamais la clé :
        propriétaire (\champ{SID}), \champ{Guid} (celui de l'en-tête, vérifié
        égal au nom du fichier), version, politique, taille de chaque bloc —
        \champ{DomainKeyLength} non nul : compte d'un domaine —, dates du
        fichier, et \champ{Preferred} avec \champ{PreferredExpiresUtc} pour la
        clé en usage.
  \item[\fichier{startup\_folders.json}] contenu des dossiers de démarrage, de
        tous les utilisateurs (\champ{Scope} \texttt{AllUsers}) et de chaque
        profil (\texttt{User}, avec \champ{SID}) : chaque fichier, ses dates et
        ses empreintes ; un raccourci est décodé (\champ{Shortcut}, comme dans
        \fichier{recentdocs.json}).
  \item[\fichier{shortcuts.json}] raccourcis posés pour être cliqués, hors du
        dossier Recent : menu Démarrer de tous les utilisateurs et de chaque
        profil, Bureau de chaque profil et Bureau public, Quick Launch et
        éléments épinglés à la barre des tâches (\champ{Location} :
        \texttt{StartMenu}, \texttt{Desktop}, \texttt{QuickLaunch},
        \texttt{TaskBar} ; \champ{Scope} et \champ{SID} comme ci-dessus).
        Chaque raccourci est décodé comme dans \fichier{recentdocs.json} :
        cible, arguments, volume et machine d'origine. Un programme désinstallé
        y laisse son raccourci ; un intrus y pose une ligne de commande, ou
        redirige un raccourci existant vers son binaire en gardant l'icône. Le
        dossier Startup du menu Démarrer n'y figure pas : il relève de
        \fichier{startup\_folders.json}. Un Bureau redirigé ailleurs sur un
        disque local (OneDrive, un autre volume, voir
        \fichier{user\_shell\_folders.json}) est lu là aussi.
  \item[\fichier{user\_shell\_folders.json}] dossiers de chaque utilisateur
        (\path|User Shell Folders| de son \fichier{NTUSER.DAT}) : nom
        (\champ{Name}, et \champ{FolderName} pour ceux qu'un GUID désigne,
        comme les Téléchargements), valeur telle qu'écrite (\champ{Value}),
        chemin résolu (\champ{Path}, absent pour un partage réseau) et date de
        dernière écriture de la clé. OneDrive déplace le Bureau et les
        Documents, une entreprise les redirige vers un partage : c'est là qu'il
        faut chercher les fichiers de l'utilisateur, et une redirection posée
        en silence est elle-même une trace.
  \item[\fichier{powershell\_history.json}] historique PSReadLine de chaque
        profil et de chaque hôte (\champ{Host} : \texttt{ConsoleHost},
        l'ISE, VS~Code) : les commandes dans l'ordre où elles ont été tapées,
        une commande sur plusieurs lignes réunie. Aucune date par commande :
        \champ{SourceModified} date la dernière. Cherché dès Windows~7 : par
        défaut depuis Windows~10, il existe sous 7, 8 et 8.1 dès que
        PowerShell~3 ou plus et PSReadLine y ont été installés. Vista et XP
        s'arrêtent à PowerShell~2, sans historique sur disque ; leurs
        journaux PowerShell sont dans \fichier{events.json}.
  \item[\fichier{powershell\_profiles.json}] profils PowerShell, les scripts
        que PowerShell exécute à chaque démarrage — une persistance qui ne
        passe ni par le registre, ni par une tâche, ni par un service. Ceux de
        la machine, pour tous les utilisateurs (dossier de Windows PowerShell
        dans \path|System32| et \path|SysWOW64|, de PowerShell~6 et~7 dans
        \path|Program Files\PowerShell|), et ceux de chaque utilisateur
        (\path|Documents\WindowsPowerShell| et \path|Documents\PowerShell|,
        là où User Shell Folders place les documents s'ils sont déplacés) :
        \champ{Scope} (\texttt{Machine}, \texttt{User}), \champ{Edition},
        \champ{Host} (\texttt{AllHosts} pour \fichier{profile.ps1}, sinon la
        console, l'ISE, VS~Code), \champ{Path}, dates, \champ{Size},
        \champ{Sha256} et \champ{Content}, le texte du script. Cherchés sur
        toutes les versions : sous XP et Vista, PowerShell n'est présent qu'une
        fois installé.
  \item[\fichier{activities\_cache.json}] Chronologie Windows
        (\fichier{ActivitiesCache.db} de chaque compte de chaque profil, avec
        son journal \texttt{-wal}) : chaque ligne des tables
        \texttt{Activity} et \texttt{ActivityOperation} (\champ{Table}),
        toutes ses colonnes. Sous Windows~10, les applications ouvertes, les
        fichiers et pages vus (\champ{Payload} : titre, URI), le temps passé
        au premier plan (\champ{Payload} \texttt{UserEngaged}, de
        \champ{StartTime} à \champ{EndTime}) et, si l'historique du
        presse-papiers est actif, ce qui a été copié
        (\champ{ClipboardPayload}) ; sous Windows~11, qui a retiré la
        Chronologie, les seuls réglages synchronisés. Les dates sont des temps
        Unix en UTC, donnés en UTC et en heure locale ; les GUID dans l'ordre
        d'octets de .NET ; une date ou un GUID nul n'est pas émis. Les codes
        (\champ{ActivityType}, \champ{ActivityStatus},
        \champ{OperationType}) sont donnés tels quels : Microsoft ne les
        documente pas.
  \item[\fichier{search\_index.json}] index de recherche Windows
        (Windows~11 : \fichier{Windows.db} et \fichier{Windows-gather.db},
        avec leurs journaux \texttt{-wal}). Une fiche par élément indexé du
        magasin de propriétés (\champ{Table} \texttt{PropertyStore},
        \champ{WorkId}) : chemin (\champ{ItemUrl}, et
        \champ{ItemPathDisplay} dans la langue de l'interface), dates de
        création, de modification, d'accès et d'indexation, taille,
        propriétaire, type et, pour un document, le début de son texte
        (\champ{SearchAutoSummary}) ; chaque propriété sous son nom Windows
        sans le préfixe \texttt{System.}. Une fiche par document vu par le
        collecteur (\champ{Table} \texttt{Gather}, \champ{DocumentId}, le
        même numéro que \champ{WorkId}) : son URL et sa dernière modification
        au passage du collecteur (sous Vista, son chemin \champ{ItemPath} en
        capitales, tiré de la table des chemins, et toutes ses dates en
        FILETIME). L'enregistrement survit au fichier jusqu'au
        passage suivant : un document supprimé peut encore y paraître.
        Jusqu'à Windows~10, l'index est une base ESE, \fichier{Windows.edb} :
        ses journaux y sont rejoués dans la copie de travail, puis il est
        converti sous la même forme ; le type d'une propriété, que cette base
        ne dit pas, vient de la table que Windows~11 déclare, et une propriété
        qu'elle ne connaît pas est donnée en hexadécimal. Mesuré sous
        Windows~8.1, et sous Windows~7, dont la base nomme autrement ses tables
        et ses colonnes, écrit ses nombres et ses dates en gros-boutiste et
        brouille son texte ; une valeur de Windows~10 compressée autrement que sur
        7~bits (XPRESS) est signalée, non décodée.
        \textbf{Attention} : sous Windows~7 et~8, l'indexeur garde en mémoire
        plusieurs minutes ce qu'il vient d'indexer avant de l'écrire dans sa
        base ; \wac{} lit la base telle qu'elle est sur le disque --- il ne
        redémarre aucun service sur la machine examinée --- et les dernières
        minutes peuvent y manquer.
  \item[\fichier{rdp\_client.json}] connexions Bureau à distance \emph{sortantes} de
        chaque profil : connexions récentes (\champ{RecentConnection},
        \champ{MruPosition} 0 pour la plus récente) et serveurs contactés
        (\champ{Server} : nom d'utilisateur proposé \champ{UsernameHint},
        empreinte du certificat accepté \champ{CertHash}, dernière écriture de
        la clé), fichiers \fichier{.rdp} des Documents (\champ{ConnectionFile},
        \fichier{Default.rdp} compris, chaque réglage), et cache bitmap
        (\champ{BitmapCache}) : chaque tuile de \fichier{Cache*.bin}, morceau de
        l'écran distant tel que le client l'a dessiné, est exportée en image
        dans \fichier{rdp\_cache\textbackslash<SID>\textbackslash<fichier>}, avec
        sa taille, ses clés et le SHA-256 de ses pixels. L'ancien format
        \fichier{bcache22.bmc}, celui du client de Windows~XP, est décodé de
        même : tuiles compressées par le RLE entrelacé de RDP, 16~bits par
        pixel (\champ{BitsPerPixel}), \champ{UnreadTiles} comptant celles qui
        n'ont pu l'être. Ce cache garde ce qui s'est affiché à l'écran,
        y compris un texte saisi en clair dans un champ visible. Un autre
        \fichier{.bmc}, d'une profondeur dont WAC n'a pas d'échantillon, est
        compté sans être décodé (\champ{DecodingStatus}
        \texttt{UnsupportedFormat}).
  \item[\fichier{rdp\_server.json}] connexions Bureau à distance \emph{entrantes},
        lues dans les journaux collectés : connexion reçue et authentification
        (RemoteConnectionManager 261 et 1149), connexion TCP acceptée (RdpCoreTS
        131), sessions (LocalSessionManager 21 à 25, 39 à 41), ouvertures et
        fermetures de session de type~10 et reconnexions (Security 4624, 4625,
        4634, 4778, 4779). Chaque événement garde ses données et reçoit
        \champ{User}, \champ{Domain}, \champ{SourceAddress} et
        \champ{SessionId}, quel que soit le nom que lui donne son journal. Peu
        utile sur un poste, c'est la première question sur un serveur.
  \item[\fichier{emails.json}] courriels d'Outlook de chaque profil : les PST
        (\fichier{Documents\textbackslash Outlook Files}) et les OST, copie
        hors connexion d'un compte Exchange ou Microsoft~365
        (\fichier{AppData\textbackslash Local\textbackslash Microsoft\textbackslash Outlook}),
        lus par le code de \wac{} — PST Unicode d'Outlook 2003 et suivants, OST
        d'Outlook 2013 et suivants (pages de 4~Kio, blocs compressés). Une fiche
        par boîte (type, version, nombre de messages lus, de messages
        illisibles) et une par message : dossier, classe, objet, expéditeur,
        destinataires (\champ{Recipients} : nom, adresse, To/Cc/Bcc), dates
        d'envoi, de remise, de création et de modification, identifiant
        Internet du message, en-têtes de transport. Le corps (texte, HTML) et
        chaque pièce jointe sont écrits à côté du JSON
        (\fichier{emails\textbackslash<SID>\textbackslash<fichier>\textbackslash<message>\textbackslash}),
        avec leur taille et leur SHA-256 ; le nom d'une pièce jointe, choisi
        par l'expéditeur, n'est jamais pris pour un chemin : le fichier porte
        son rang et, si elle est simple, son extension. Un message joint (un
        rendez-vous, un courriel transféré en pièce jointe) donne son objet et
        sa classe. Les fichiers qu'Outlook a ouverts depuis un message
        (\fichier{INetCache\textbackslash Content.Outlook}) et les données
        locales du nouvel Outlook (\fichier{AppData\textbackslash Local\textbackslash Microsoft\textbackslash Olk})
        sont prélevés tels quels. Non décodés : le chiffrement « élevé »
        (cyclique) et le format ANSI d'Outlook~97 à 2002, sans échantillon pour
        les vérifier — la fiche de la boîte le dit. Thunderbird : ses profils
        (\fichier{AppData\textbackslash Roaming\textbackslash Thunderbird\textbackslash Profiles})
        sont prélevés entiers — carnet d'adresses et réglages des comptes
        compris — et chaque dossier de courrier, un fichier mbox, est décodé :
        en-têtes (mots encodés compris), date d'envoi en UTC, corps texte et
        HTML, pièces jointes (base64, quoted-printable, noms RFC~2231), dans
        les mêmes champs que pour Outlook, avec \champ{Offset} (position du
        message dans le fichier) au lieu de \champ{NodeId}.
  \item[\fichier{sysmon.json}] Sysmon (Sysinternals), là où il est
        installé : une fiche par installation, reconnue au lien entre le
        service (\texttt{Parameters\textbackslash DriverName}) et la clé de son
        pilote qui porte des règles — trouvée même renommée. Le service et
        le pilote avec leur \texttt{ImagePath}, la date de dernière écriture
        de \texttt{Parameters} (dernier changement de configuration), toutes
        les valeurs de cette clé telles que stockées (\champ{Parameters} :
        \texttt{Rules}, \texttt{HashingAlgorithm}, \texttt{Options},
        \texttt{ArchiveDirectory}, \texttt{ConfigFile}, \texttt{ConfigHash}),
        \champ{RulesSha256}. Le format binaire des règles n'est pas documenté
        par Sysinternals : elles sont données en hexadécimal, non décodées.
        Le dossier d'archives (\champ{ArchiveFolder}, \fichier{C:\textbackslash Sysmon}
        par défaut) est prélevé entier : ce sont les fichiers supprimés sous
        une règle \texttt{FileDelete}, nommés par leur empreinte — le contenu
        même d'un outil effacé —, listés dans \champ{ArchivedFiles} avec leur
        taille et leur SHA-256 ; le fichier de configuration XML, s'il est
        encore là, est prélevé et son SHA-256 donné (\champ{ConfigFileSha256},
        à comparer à \texttt{ConfigHash}). Le journal de Sysmon est converti
        avec les autres dans \fichier{events.json}, même déplacé ; sa
        version y est (événement 4).
  \item[\fichier{thumbcache.json}] caches des miniatures et des icônes de
        l'Explorateur, de chaque profil (\fichier{thumbcache\_*.db},
        \fichier{iconcache\_*.db}) : ce que l'utilisateur a eu sous les yeux,
        même si le fichier a été supprimé depuis. Pour chaque entrée : son
        cache et sa taille, sa clé (\champ{Hash}), son identifiant, le
        \champ{ThumbnailCacheId} que Windows donne pour le fichier, ce qui
        relie une miniature à son fichier (la clé depuis Windows~7,
        l'identifiant sous Vista), ses dimensions, son format (\texttt{png},
        \texttt{jpg}, \texttt{bmp}), le SHA-256 de l'image, et l'image
        elle-même, exportée à côté du JSON
        (\fichier{thumbcache\textbackslash<SID>\textbackslash<cache>\textbackslash<ThumbnailCacheId>.<ext>}).
        Les deux sommes de contrôle que Windows écrit pour chaque entrée sont
        recalculées (\champ{HeaderChecksumValid}, \champ{DataChecksumValid}) :
        une entrée lue de travers les échoue. Vista à Windows~11 (Vista et 7
        ne donnent pas les dimensions ; Vista donne l'extension du fichier
        d'origine, \champ{SourceExtension}). Après la première place libre
        que l'en-tête annonce, des entrées d'un état antérieur peuvent rester
        entières : elles sont rendues comme vestiges (\champ{Remnant}) si leur
        somme d'en-tête est juste, leur image seulement si sa somme l'est aussi
        (\fichier{<ThumbnailCacheId>-remnant-<position>.<ext>}) ; la lecture
        s'arrête à la première qui ne l'est pas.
        Sous Windows~XP et Server~2003, les \fichier{Thumbs.db} que
        l'Explorateur écrit dans chaque dossier affiché en miniatures,
        cherchés par leur nom sur chaque volume NTFS fixe : pour chaque
        miniature, le nom et le chemin de l'image, sa date de modification
        au moment de la miniature (\champ{SourceModifiedUtc}), même une fois
        l'image effacée, la taille et le SHA-256 de la miniature, et ses
        octets exportés tels quels (\texttt{.bin}) : un JPEG à quatre
        composantes sans ses tables, qu'aucun visualiseur ordinaire n'affiche.
  \item[\fichier{error\_reports.json}] rapports d'erreurs Windows
        (\fichier{Report.wer}) de la machine et de chaque profil : programme,
        chemin, instant (\champ{EventTimeUtc}), signature du problème
        (module fautif, version), toutes les valeurs du rapport dans
        \champ{Values}.
  \item[\fichier{defender\_log.json}] journaux de support de Microsoft Defender
        (\fichier{MPLog-*.log}) : chaque ligne telle qu'écrite, datée quand elle
        commence par une date UTC. Analyses à l'exécution, détections,
        exclusions, mises à jour. Les traces WPP du même dossier
        (\fichier{MpWpp*.bin}, traces ETW binaires que seuls les symboles
        privés de Microsoft décodent) sont copiées brutes dans la consigne, sans
        conversion.
  \item[\fichier{defender\_detections.json}] une entrée par fichier
        DetectionHistory, qui survit à l'effacement du journal d'événements :
        menace, ressource (fichier, clé), compte, processus, date de détection
        initiale, entrée de quarantaine, et les propriétés nommées
        \champ{ThreatTracking} (empreintes, taille, heure de début). Le format
        n'est pas documenté : seuls les champs dont la position a été confirmée
        contre \texttt{MSFT\_MpThreatDetection} sont nommés ; les autres sont
        rendus tels que lus (\champ{UnnamedFields} : position, type, valeur),
        sans signification inventée. Defender sous Windows~8.1 écrit les mêmes
        champs sans le marqueur \champ{MagicVersion} de Windows~10 et~11 : ils
        sont nommés de la même façon, confirmés de la même façon.
  \item[\fichier{defender\_quarantine.json}] entrées de la quarantaine :
        menace, instant de mise en quarantaine, chemin d'origine du fichier
        retiré, ses dates de création et de modification, ses attributs, sa
        taille, et l'empreinte de son contenu (\champ{ContentSha256}). La
        quarantaine est copiée \textbf{chiffrée}, telle que Defender la garde :
        rien de malveillant n'est écrit en clair, ni sur le support de collecte
        ni sur le poste d'analyse. Le contenu n'est déchiffré qu'en mémoire,
        pour son empreinte.
  \item[\fichier{hosts.json}] entrées du fichier \fichier{hosts} : adresse et
        noms, numéro de ligne ; un Windows intact n'en a aucune. Un fichier de plus
        de 16~Mio n'est pas lu (journalisé) : il reste dans la consigne.
\end{description}

\subsection{Fichiers et navigation}

\begin{description}[style=nextline,leftmargin=1.2em]
  \item[\fichier{mft.json}] chaque fichier et répertoire de chaque volume NTFS
        fixe (\champ{Volume} : \texttt{C:}, \texttt{G:}\ldots{} ; un numéro
        d'enregistrement n'a de sens que dans son volume),
        présent ou \textbf{supprimé} (\champ{Deleted}) tant que NTFS n'a pas
        réutilisé son enregistrement : chemin reconstitué, taille, attributs, et
        deux jeux de quatre dates en UTC — \champ{CreatedUtc}\ldots{} (attribut
        \path|$STANDARD_INFORMATION|, ce que montre l'Explorateur et que
        réécrit un outil d'antidatage) et \champ{FileNameCreatedUtc}\ldots{}
        (attribut \path|$FILE_NAME|, écrit par NTFS seul). Une date
        \path|$STANDARD_INFORMATION| antérieure à celle de \path|$FILE_NAME|
        trahit un antidatage. \champ{PathIncomplete} marque un chemin dont un
        répertoire parent a été réutilisé depuis : le chemin n'est pas inventé.
        Un fichier à liens physiques (\champ{HardLinks}) est un seul
        enregistrement sous plusieurs noms — \path|System32\crypt32.dll| et son
        original sous \path|WinSxS| : \champ{HardLinkPaths} donne ses autres
        chemins, chacun cherchable. \champ{SecurityId} est le numéro de son
        descripteur de sécurité dans \fichier{security\_descriptors.json}.
  \item[\fichier{security\_descriptors.json}] descripteurs de sécurité de
        chaque volume NTFS fixe (\path|$Secure:$SDS|) : chacun une seule fois,
        quel que soit le nombre de fichiers qui le partagent, repéré par
        \champ{Volume} et \champ{SecurityId} — la clé que porte chaque fichier
        de \fichier{mft.json}. Il donne le \textbf{propriétaire}
        (\champ{Owner}, \champ{OwnerName}), le groupe, et chaque entrée de la
        DACL et de la SACL (type, drapeaux, masque d'accès, SID et nom). L'outil
        d'un intrus déposé par son compte, un fichier système qui n'appartient
        plus à TrustedInstaller, une ACL ouverte à Everyone s'y lisent, fichiers
        supprimés compris. NTFS range avec chaque descripteur un condensat de
        ses octets : WAC le recalcule, et \champ{HashMismatch} signale un
        descripteur qui ne lui correspond pas.
  \item[\fichier{deleted\_files.json}] le \textbf{contenu} des fichiers
        supprimés, là où \fichier{mft.json} n'en donne que l'enregistrement :
        l'outil, le script ou l'archive qu'un intrus a effacés restent dans
        leurs clusters tant que NTFS ne les a pas donnés à un autre fichier.
        Lus en brut, sans rien ouvrir ni écrire sur la machine examinée. Le
        contenu résident (quelques centaines d'octets, dans l'enregistrement
        même) est rendu tel quel ; chaque cluster d'un contenu non résident
        est vérifié libre dans \path|$Bitmap| avant et après sa lecture :
        donné entre-temps à un autre fichier, il devient un trou de zéros. Un
        cluster pris puis rendu depuis la suppression se relit libre : les
        plages de tous les enregistrements supprimés de la \path|$MFT| sont
        donc croisées, et un cluster qu'un fichier supprimé plus récent a
        tenu devient un trou (\champ{Reason} « clusters claimed since by a
        later deleted file »). Un fichier intermédiaire dont l'enregistrement
        a été réutilisé ne laisse pas de plages : ce qu'il a écrit ne se voit
        pas ainsi.
        \champ{State} : \texttt{Recovered} (tout le contenu),
        \texttt{Partial} (trous listés dans \champ{MissingRanges}, octet de
        début et longueur), \texttt{Overwritten} (plus rien), \texttt{Zeroed}
        (clusters relus à zéro : un disque SSD qui a reçu TRIM),
        \texttt{NotRecoverable} et \champ{Reason} (compressé, chiffré par EFS,
        plages en partie dans un enregistrement d'extension, plages
        incohérentes), \texttt{RecordReused} (l'enregistrement a servi à un
        autre fichier depuis la copie de la \path|$MFT|), \texttt{NotRead}
        (non lu : \texttt{max\_total\_mib} atteint). Avec le chemin
        (\champ{Path}), les dates de l'enregistrement, l'heure de la
        suppression (\champ{DeletedUtc}, lue dans le journal des
        modifications : aucun enregistrement ne la porte),
        \champ{RecoveredBytes} et les empreintes \champ{RecoveredSHA256}\ldots{}
        du contenu \emph{tel que récupéré}, trous compris : celles du fichier
        d'origine seulement si l'état est \texttt{Recovered}. Le contenu copié
        est dans la consigne sous \path|X:\$DeletedFiles\<enregistrement>-<séquence>_<nom>|
        (\champ{Exhibit}), jamais sous son ancien chemin, où un autre fichier
        peut se trouver. Depuis Windows~10, un fichier supprimé encore ouvert
        est d'abord renommé dans \path|$Extend\$Deleted| : son
        enregistrement l'y nomme, et \champ{FormerPath} donne le chemin
        qu'il avait, lu dans le journal des modifications. Un exécutable récupéré suit la règle des autres
        (\texttt{binary}) : authentique, il est relevé avec son
        \champ{Signature} sans être copié (\champ{Copied}~\texttt{false},
        \champ{NotCopied}) ; sous \texttt{binary: none}, rien n'est jugé et il est
        copié. Les plus récents d'abord (dernière date de
        l'enregistrement, celle de la suppression n'y figurant pas), jusqu'à
        \texttt{max\_total\_mib} ; les dossiers \texttt{excluded\_folders}
        (mises à jour de Windows) ne sont pas lus, ni le dossier de sortie.
        Le réglage est une structure sous \texttt{artefacts:} ;
        \texttt{deleted\_files: false} suffit pour l'écarter.
        \champ{Source} dit d'où vient l'entrée : \texttt{MftRecord}, un
        enregistrement supprimé de la \path|$MFT| (niveau~1) ;
        \texttt{IndexSlack} (niveau~2), un nom resté dans la marge d'un bloc
        d'index de répertoire (\path|$I30|) ou dans un bloc que le répertoire
        a libéré : le nom, le répertoire (\champ{Path}), l'enregistrement du
        fichier (\champ{Entry}, \champ{Sequence}), sa taille et les quatre
        dates de son \path|$FILE_NAME| (\champ{FileNameCreatedUtc}\ldots),
        \champ{State} \texttt{NameOnly} — aucun contenu. Seul un nom
        qu'aucun répertoire ne porte plus est donné, et pas un nom déjà rendu
        par le niveau~1. \champ{CurrentPath} : le même fichier existe encore
        sous ce chemin, déplacé ou renommé et non supprimé.
        \texttt{Carved} (niveau~3, \texttt{level: 3} seulement : tout
        l'espace libre est relu, aussi longtemps qu'il faut) : un fichier
        taillé dans les clusters libres que le niveau~1 n'a pas rendus, sans
        nom ni date — \champ{Format} (\texttt{PE}, \texttt{ZIP} et les
        documents Office, \texttt{PDF}), premier cluster (\champ{Lcn}),
        \champ{Size}. Seuls ces formats, dont la structure donne la taille
        exacte et se vérifie, sont taillés ; un fichier n'est pris que si tous
        ses clusters sont contigus, libres et non rendus par le niveau~1 —
        un fichier fragmenté n'est pas trouvé, jamais coupé faux. Pièce sous
        \path|X:\$CarvedFiles\<cluster>.<ext>|.
        \texttt{FatEntry} : une entrée supprimée d'un volume FAT ou exFAT (une
        clé du suspect) : son premier cluster (\champ{FirstCluster}), sa
        taille, ses dates telles que stockées (\champ{ModifiedAsStored},
        heure locale de la machine qui l'a écrite, ou \champ{ModifiedUtc}
        quand exFAT en donne le décalage). La chaîne de la FAT est effacée :
        le contenu est relu comme les clusters qui suivent le premier, chacun
        vérifié libre — \champ{State} \texttt{RecoveredIfContiguous} et
        \champ{ContiguityAssumed}, sauf quand exFAT savait le fichier d'un
        seul tenant (\texttt{Recovered}). Le niveau~3 taille aussi l'espace
        libre de ces volumes.
        Sur un volume BitLocker, la lecture brute voit le volume déchiffré par
        le système, comme pour les autres artefacts : les clusters libres
        aussi.
  \item[\fichier{usn\_journal.json}] journal des modifications de chaque
        volume NTFS fixe (\path|$UsnJrnl:$J|, \champ{Volume} ; les numéros
        \champ{Usn} croissent au sein d'un volume) : chaque création, écriture, renommage et
        \textbf{suppression}, daté, avec ses motifs (\champ{Reasons}, un bit
        inconnu rendu en hexadécimal) et le chemin reconstitué par la
        \path|$MFT|. Il garde la trace de fichiers déposés puis effacés, que la
        \path|$MFT| a déjà oubliés.
        La \path|$MFT|, le journal des modifications, \path|$LogFile| (journal
        de transactions, non converti) et \path|$TxfLog| (journal de NTFS
        transactionnel, journal CLFS, ses conteneurs et le flux \path|$Tops:$T|, non converti) sont en
        outre copiés bruts dans la consigne, pour refaire plus tard la
        chronologie avec un autre outil.
  \item[\fichier{ntfs\_volumes.json}] chaque volume NTFS fixe, lu dans son
        secteur d'amorçage (\path|$Boot|, reconnu à ses signatures avant tout
        champ) : numéro de série (\champ{VolumeSerial}, les 32 bits que les
        raccourcis, listes de sauts et fichiers Prefetch relèvent du volume de
        leur cible, et \champ{VolumeSerial64} en entier), tailles de secteur,
        de cluster, d'enregistrement et de bloc d'index, emplacement de la
        \path|$MFT| et de son miroir. Puis l'en-tête du journal des
        modifications (\path|$UsnJrnl:$Max|, \champ{UsnJournal}) : taille
        maximale, pas d'allocation, \champ{UsnJournalId}, qui est la date de
        \textbf{création} du journal (\champ{Created}), et premier numéro
        encore présent. Un journal supprimé puis recréé (\texttt{fsutil usn
        deletejournal}) se lit à une création bien postérieure à
        l'installation.
  \item[\fichier{recentdocs.json}] raccourcis de documents récents : cible,
        horodatages du raccourci et de la cible, taille de la cible, volume
        (type, numéro de série, nom), chemin réseau, liste d'éléments du shell,
        et les blocs de données supplémentaires : le bloc de suivi
        (\champ{Tracker} : nom NetBIOS de la machine où était la cible,
        identifiants de suivi, et ce qu'en révèle le GUID de version 1 — la date
        de début du suivi et l'adresse matérielle de la machine), chemins avec
        variables d'environnement, dossier connu, propriétés de la cible.
  \item[\fichier{jumplistAutomaticDestinations.json},
        \fichier{jumplistCustomDestinations.json}] listes de sauts par
        application, avec les raccourcis qu'elles contiennent ; pour les listes
        automatiques, chaque raccourci est accompagné de son entrée DestList
        (dernier accès, nom de la machine, identifiants de suivi avec leur date
        et leur adresse matérielle, épinglage, nombre d'ouvertures
        \champ{InteractionCount}, rang \champ{MruPosition} — 0 pour le plus
        récent —, nom du dossier connu visé) ; pour les listes personnalisées,
        leurs catégories dans l'ordre (\champ{Categories} : une catégorie
        propre et son nom, une catégorie connue — fréquents ou récents —, les
        tâches).
  \item[\fichier{shellbags.json}] dossiers parcourus dans l'Explorateur,
        arborescence reconstituée ; pour chaque dossier, son rang parmi ses
        voisins (\champ{MruPosition}, 0 pour le plus récent), la date d'écriture
        de la clé qui le porte (\champ{LastWriteTime}) et de sa propre clé
        (\champ{KeyLastWriteTime} : son ajout, ou le parcours d'un sous-dossier),
        et son \champ{NodeSlot}. Les « première » et « dernière interactions »
        d'autres outils ne sont que des lectures de ces deux dates. Sous
        Windows~XP, les dossiers de l'Explorateur sont dans NTUSER.DAT
        (\texttt{ShellNoRoam\textbackslash BagMRU}), lus de même.
  \item[\fichier{mrus.json}, \fichier{mruApps.json}] listes de fichiers et
        d'applications récemment utilisés. Sous Windows~XP, en texte : le chemin
        du fichier (\champ{Path}), le programme et son dossier (\champ{Folder}),
        rangés par \champ{ID}, 0 pour le plus récent.
  \item[\fichier{zone\_identifiers.json}] flux \fichier{Zone.Identifier}
        (« marque du web ») que navigateurs et messageries écrivent sur ce qu'ils
        téléchargent : fichier, zone (\champ{ZoneId}, \champ{ZoneName} —
        \texttt{Internet} pour~3), page d'origine (\champ{ReferrerUrl}) et
        adresse du fichier (\champ{HostUrl}), toutes les valeurs dans
        \champ{Values}. Lus dans la copie de la \path|$MFT|, où ce petit flux
        est résident : aucune lecture de plus sur la machine examinée ; un flux
        non résident, rare, est signalé (\champ{ContentNotResident}).
  \item[\fichier{recycle\_bin.json}] fichiers mis à la corbeille, sur chaque
        volume NTFS fixe, d'après leur index \path|$I| : chemin d'origine,
        taille, instant de la suppression (\champ{DeletedUtc}), et l'utilisateur
        (\champ{SID}). Seuls les \path|$I| sont prélevés ; les fichiers eux-mêmes
        (\path|$R|) figurent dans \fichier{mft.json}. Sous Windows~XP et
        Server~2003, d'après \fichier{INFO2} de \path|RECYCLER\<SID>|, une
        entrée par enregistrement, avec le nom de la copie dans la corbeille
        (\champ{RecycledName}, \texttt{D<lecteur><n>.<ext>}, \texttt{Dc1.txt}
        pour un fichier de C:) ; la taille est l'espace
        qu'occupait le fichier.
\end{description}

\subsection{Périphériques}

\begin{description}[style=nextline,leftmargin=1.2em]
  \item[\fichier{setupapi.json}] sections de \fichier{setupapi.dev.log} et de
        ses rotations : opération (\champ{Operation}), périphérique ou pilote
        (\champ{Subject}, avec le numéro de série d'une clé USB), début et fin
        convertis depuis l'heure locale de la machine examinée, résultat. La
        première connexion d'un périphérique s'y date, même quand la clé
        \fichier{USBSTOR} a été nettoyée. Sous Windows~XP, d'après
        \fichier{setupapi.log} : une section par ligne datée, son premier
        message en opération, le périphérique qu'une ligne \texttt{\#I121}
        nomme en sujet, \texttt{SUCCESS} par cette ligne ou la première erreur
        en résultat ; pas de date de fin.
  \item[\fichier{Usbstor.json}] périphériques de stockage USB connus du système :
        \champ{Vendor}, \champ{Product}, \champ{Revision} et \champ{Type}, tirés
        de \champ{DeviceId} ; \champ{SerialNumber}, tiré de \champ{InstanceId}
        — absent quand l'identifiant a été fabriqué par Windows (deuxième
        caractère « \& »), le périphérique n'en déclarant aucun ; dates de
        première et de dernière insertion.
  \item[\fichier{mounted\_device.json}] correspondance entre lettres de lecteur et
        périphériques. \champ{Drive} nomme la lettre ou le volume ;
        \champ{Type} dit la forme de la valeur : \texttt{GPT partition}
        (\champ{PartitionGuid}), \texttt{MBR partition}
        (\champ{DiskSignature} et \champ{PartitionOffset}, en octets),
        \texttt{Device path} (\champ{Device}, chemin du périphérique, avec le
        numéro de série d'une clé USB) ou \texttt{Unknown} (\champ{Data}, les
        octets en hexadécimal).
\end{description}

\subsection{Journaux d'événements (\opt{events})}

\fichier{events.json} contient un objet par événement, de la même forme
quel que soit le format du journal : les journaux \fichier{.evt} de
Windows~XP et Server~2003 sont lus dans les champs que Windows rend pour une
source classique dans un \fichier{.evtx} (identifiant et ses qualificateurs,
niveau d'après le type d'événement, catégorie en tâche, chaînes en
\champ{Data}, données binaires en \champ{Binary}), leur zone circulaire
parcourue depuis l'enregistrement le plus ancien. Un journal que le registre
de la machine place hors de son dossier (valeur \texttt{File} de son canal,
\texttt{wevtutil sl /lfn}, ou de son journal classique) — déplacé pour la
place, ou par un intrus pour qu'on ne le cherche pas — est prélevé là où il
est et converti avec les autres :

\begin{description}[style=nextline,leftmargin=1.2em]
  \item[\champ{EvtSystem*}] la section système : fournisseur, identifiant,
        niveau, tâche, mots-clés, horodatage de création, numéro d'enregistrement,
        processus et fil, canal, machine, utilisateur ;
  \item[\champ{EvtEventData}] les données de l'événement, en liste de paires
        \champ{Name} / \champ{Value} ;
  \item[\champ{EvtEventMessage}] le message lisible, tel que l'Observateur
        d'événements l'afficherait, reconstitué à partir des ressources du
        fournisseur prélevées sur la machine : manifeste (\texttt{WEVT\_TEMPLATE},
        ou son jumeau \fichier{SystemResources\textbackslash*.mun} sous
        Windows~11) et table de messages dans la langue de l'interface de la
        machine, ou fichiers \texttt{EventMessageFile} d'une source classique.
        Chaque donnée y est montrée comme Windows la montre : libellé de sa
        table de valeurs ou de ses drapeaux, hexadécimal, texte d'un code
        d'erreur Win32, HRESULT ou NTSTATUS, adresse réseau ;
  \item[\champ{EvtSourceLog}] le fichier journal dont provient l'événement. Un
        même canal peut être porté par plusieurs fichiers (journal courant et
        archives) dont les numéros d'enregistrement se recouvrent ;
  \item[\champ{EvtRecovered}] \texttt{true} pour un événement d'un journal
        \emph{vidé} : Windows remet l'en-tête du bloc à zéro sans effacer les
        enregistrements, que WAC relit. L'Observateur d'événements ne les montre
        plus ; ils restent la trace de ce que l'effacement voulait faire
        disparaître, et ne doivent jamais être lus comme des événements vivants.
\end{description}

\begin{note}[Messages incomplets]
Une marque \texttt{\%1}, \texttt{\%2}\ldots{} restée dans un message signifie que
l'événement ne porte pas la donnée attendue : elle est laissée visible pour que la
phrase ne paraisse pas complète. Tous les fournisseurs ne livrent pas de
ressources de messages : leurs événements n'ont alors pas de
\champ{EvtEventMessage}, mais leurs données restent complètes.
\end{note}

% ============================================================ 6
\chapter{Lire les données}\label{chap:lecture}

\section{Horodatages}

Chaque date est émise deux fois : \champ{X} en heure locale et \champ{XUtc} en
UTC. L'heure locale est celle du \textbf{suspect}, calculée avec le fuseau lu
dans sa ruche \fichier{SYSTEM} — pas celle de la machine de collecte. Les dates
portent la fraction de seconde quand la source la fournit, et seulement dans
la précision de cette source : sept chiffres pour un \texttt{FILETIME}, trois
pour une date à la milliseconde, aucun pour une date à la seconde (date FAT,
précise à deux secondes, date lue en texte, temps Unix).

Chaque date porte le décalage en vigueur \textbf{à cette date}, heure d'été
comprise : une date de janvier relevée en été à Paris est à \texttt{+01:00},
une date de juillet à \texttt{+02:00}. \wac{} applique les règles de la machine
examinée année par année (clé \fichier{Dynamic DST} de la ruche
\fichier{SOFTWARE}) : une date de 2005 aux États-Unis suit les règles de 2005.
L'heure locale écrite est toujours dérivée de l'heure UTC : un même instant se
lit à la même heure dans tous les artefacts. Deux dates n'ont aucun fuseau
connaissable et sont données \textbf{telles qu'écrites}, sans décalage ni UTC,
avec \champ{TimeZone} qui le dit : la date d'un élément vu dans une archive
ZIP (\champ{ModifiedAsWritten} ; une entrée ZIP garde l'heure d'horloge de la
machine qui a fait l'archive — les archives que Windows~11 lit lui-même, 7z
ou tar, portent au contraire un instant UTC, donné comme tel) et une date OLE d'un magasin de propriétés
(\champ{AsWritten}). Les placer sur une chronologie demande de connaître
leur origine : \wac{} ne l'invente pas. Pour une valeur que l'artefact
enregistre en heure locale (volume FAT, \fichier{setupapi.dev.log}), l'heure répétée à l'automne est lue comme sa première occurrence, et
l'heure sautée au printemps avec le décalage d'avant le changement, comme le
fait Windows.

\section{Noms des comptes}

Les champs qui nomment un compte (\champ{SIDName}, \champ{Owner},
\champ{RunAs}) sont tirés des preuves seules, sans interroger le système ni le
réseau : comptes et groupes locaux de la base SAM (dans la langue de la
machine examinée), SID bien connus sous leur nom canonique
(\texttt{SYSTEM}, \texttt{LOCAL SERVICE}\ldots), SID de services
(\texttt{NT SERVICE\textbackslash\ldots}), et dossier de profil pour un
compte de domaine. Un SID qu'aucune de ces sources ne nomme est publié sans
nom : le SID reste la preuve, à rapprocher du domaine par l'analyste.

\paragraph{Le profil Default.} \path|C:\Users\Default| est le modèle copié
pour chaque nouveau profil : une valeur Run, un raccourci ou une association
posés là atteignent tous les comptes créés ensuite. WAC le lit comme un profil
(ruches \fichier{NTUSER.DAT} et \fichier{UsrClass.dat}, Bureau, menu
Démarrer, Recent\ldots). N'appartenant à aucun compte, il n'a pas de SID : ses
entrées portent le SID \texttt{S-0-0-0}, fait de zéros, au format d'un SID
pour que les outils les lient, mais qui n'en est pas un (révision 0), et
\champ{Profile}~: \texttt{Default} à la place de \champ{SIDName}. Ce n'est pas
\texttt{S-1-0-0}, le vrai SID nul de Windows.

\section{Noms des GUID}

Les champs qui nomment un GUID (classe COM, dossier connu, gestionnaire d'une
tâche) viennent d'abord d'une table de référence, puis des ruches de la machine
examinée : ses classes (\fichier{SOFTWARE\textbackslash Classes\textbackslash
CLSID}), ses dossiers connus (\fichier{FolderDescriptions}), et les classes de
chaque utilisateur (\fichier{UsrClass.dat}). Un composant installé sur la
machine y est donc nommé, avec le module qui le sert. Un GUID que rien ne nomme
est publié sans nom.

\section{Un champ absent n'est pas un champ vide}

Une valeur sans source n'est pas émise du tout. Une clé vide se lirait comme une
lecture ratée, et ne se distinguerait pas de « la ruche ne contient pas cette
valeur ». Ce qui n'a pas pu être lu est consigné dans
\fichier{investigation.json}.

\section{Artefact non collecté}

Quand un artefact entier n'a pas pu être lu, son fichier le dit explicitement :

\begin{lstlisting}
{
  "Artifact": "USBSTOR (registre)",
  "CollectionStatus": "NotCollected",
  "Error": "0x02 Le fichier spécifié est introuvable.",
  "Note": "La lecture de cet artefact a échoué : l'absence de données
           ci-dessus ne signifie PAS qu'aucune trace n'existe sur le système."
}
\end{lstlisting}

\begin{attention}
Un fichier d'artefact vide et un fichier « NotCollected » ne disent pas la même
chose. Le premier indique qu'aucune entrée n'a été trouvée ; le second, que
la source n'a pas pu être lue.
\end{attention}

\section{Artefact absent de la version de Windows}\label{sec:versions}

Un artefact n'existe pas sur toutes les versions de Windows : pas de SRUM avant
Windows~8, pas de jump lists avant Windows~7. \wac{} lit la version de la machine
examinée dans sa ruche \fichier{SOFTWARE}, l'annonce au début de la collecte et
la consigne dans \fichier{investigation.json}. Un artefact demandé que cette
version ne porte pas n'est ni lancé ni compté en échec, et son fichier le dit :

\begin{lstlisting}
{
  "Artifact": "RESOURCE USAGE (SRUM)",
  "CollectionStatus": "AbsentFromThisVersion",
  "WindowsRelease": "Windows 7",
  "WindowsBuild": 7601,
  "Note": "This artefact does not exist on this version of Windows
           (artefacts.srum): there was nothing to collect. Its absence is
           normal and says nothing about what happened on the system."
}
\end{lstlisting}

Une version plus récente que celles que \wac{} connaît est traitée comme la
dernière connue, et signalée comme telle : la collecte se fait. Le tableau
ci-dessous donne, pour chaque clé de \fichier{wac.yml}, le premier Windows qui
porte l'artefact, son emplacement et son format selon la version. Il est
produit depuis la table du programme, la seule que \wac{} consulte. La
version~2 de \wac{} n'est éprouvée que sous Windows~10 et 11 ; les formes plus
anciennes y figurent pour la suite.

Les emplacements sont écrits comme Windows~Vista et suivants rangent leurs
dossiers. Sous Windows~XP et Server~2003, \wac{} les cherche là où ces versions
gardent les mêmes données : \path|AppData\Local| devient
\path|Local Settings\Application Data|, \path|AppData\Roaming| devient
\path|Application Data|, les documents récents et le menu Démarrer sont à la
racine du profil, \path|Documents| devient \path|My Documents|, et
\path|ProgramData| devient \path|All Users\Application Data|, à côté des
profils.

% Generated by doc/artefact-table.py from sources/artefact_catalogue.cpp — do not edit.
{\small
\begin{longtable}{@{}L{2.8cm}L{2.0cm}L{5.0cm}L{3.3cm}L{1.5cm}@{}}
\toprule
\textbf{Artefact} & \textbf{Depuis} & \textbf{Emplacement} & \textbf{Format} & \textbf{Converti} \\
\midrule\endhead
\texttt{activities\_cache} & Windows 10 & \texttt{Users\textbackslash{}\textless{}user\textgreater{}\textbackslash{}AppData\textbackslash{}Local\textbackslash{}ConnectedDevicesPlatform\textbackslash{}\textless{}account\textgreater{}\textbackslash{}ActivitiesCache.db} & SQLite & oui \\
 & \multicolumn{4}{@{}L{12.8cm}@{}}{\itshape Toujours écrit par Windows 11 (mesuré, build 26200), bien que la Chronologie n'y soit plus affichée. Éditions serveur à mesurer.} \\
\addlinespace
\texttt{amcache} & Windows 7 & \texttt{Windows\textbackslash{}AppCompat\textbackslash{}Programs\textbackslash{}RecentFileCache.bcf} & BCF & oui \\
 & Windows 8 & \texttt{Windows\textbackslash{}AppCompat\textbackslash{}Programs\textbackslash{}Amcache.hve} & ruche, clés File et Programs & oui \\
 & Windows 10 & \texttt{Windows\textbackslash{}AppCompat\textbackslash{}Programs\textbackslash{}Amcache.hve} & ruche, clés Inventory* & oui \\
 & \multicolumn{4}{@{}L{12.8cm}@{}}{\itshape toujours présent depuis Windows 10 ; Amcache.hve aussi sous Windows 7 avec la mise à jour KB2952664.} \\
\addlinespace
\texttt{bam} & Windows 10 & \texttt{SYSTEM\textbackslash{}CurrentControlSet\textbackslash{}Services\textbackslash{}bam\textbackslash{}UserSettings} & valeurs de registre, FILETIME & oui \\
 & Windows 10 & \texttt{SYSTEM\textbackslash{}CurrentControlSet\textbackslash{}Services\textbackslash{}bam\textbackslash{}State\textbackslash{}UserSettings} & valeurs de registre, FILETIME & oui \\
 & \multicolumn{4}{@{}L{12.8cm}@{}}{\itshape Éditions serveur à mesurer.} \\
\addlinespace
\texttt{bits} & Windows XP & \texttt{ProgramData\textbackslash{}Microsoft\textbackslash{}Network\textbackslash{}Downloader\textbackslash{}qmgr0.dat, qmgr1.dat} & fichier d'état de BITS & oui \\
 & Windows 10 & \texttt{ProgramData\textbackslash{}Microsoft\textbackslash{}Network\textbackslash{}Downloader\textbackslash{}qmgr.db} & ESE & oui \\
 & \multicolumn{4}{@{}L{12.8cm}@{}}{\itshape Avant Vista : Documents and Settings\textbackslash{}All Users\textbackslash{}Application Data\textbackslash{}Microsoft\textbackslash{}Network\textbackslash{}Downloader. Les fichiers de file d'attente mesurés sous Windows 8.1 ; les versions antérieures lues comme bits\_parser (ANSSI) les décrit, sans mesure.} \\
\addlinespace
\texttt{boot\_records} & Windows XP & \texttt{chaque disque physique : son premier mébioctet, les 16 premiers secteurs de chaque partition ; les fichiers de la partition système EFI} & MBR, GPT, secteurs de démarrage ; FAT & oui \\
 & \multicolumn{4}{@{}L{12.8cm}@{}}{\itshape La partition système EFI n'existe que sur une machine démarrée en UEFI (Windows 8 et suivants, certains 7 et Vista).} \\
\addlinespace
\texttt{cloud\_identity} & Windows 10 & \texttt{SYSTEM\textbackslash{}CurrentControlSet\textbackslash{}Control\textbackslash{}CloudDomainJoin, NTUSER.DAT\textbackslash{}…\textbackslash{}WorkplaceJoin ; TokenBroker, AAD.BrokerPlugin de chaque profil} & clés du registre ; caches de jetons copiés chiffrés & oui \\
 & \multicolumn{4}{@{}L{12.8cm}@{}}{\itshape Présent seulement sur une machine inscrite à Entra ID ou portant un compte professionnel.} \\
\addlinespace
\texttt{defender} & Windows 8 & \texttt{ProgramData\textbackslash{}Microsoft\textbackslash{}Windows Defender\textbackslash{}Support\textbackslash{}MPLog-*.log} & texte UTF-16 & oui \\
 & Windows 10 & \texttt{ProgramData\textbackslash{}Microsoft\textbackslash{}Windows Defender\textbackslash{}Scans\textbackslash{}History\textbackslash{}Service\textbackslash{}DetectionHistory} & enregistrements DetectionHistory & oui \\
 & \multicolumn{4}{@{}L{12.8cm}@{}}{\itshape Server 2012 et 2012 R2 n'ont pas de Defender intégré ; la quarantaine est copiée chiffrée sur toutes les versions.} \\
\addlinespace
\texttt{defender\_configuration} & Windows Vista & \texttt{SOFTWARE\textbackslash{}Microsoft\textbackslash{}Windows Defender, SOFTWARE\textbackslash{}Policies\textbackslash{}Microsoft\textbackslash{}Windows Defender} & valeurs de registre & oui \\
\addlinespace
\texttt{defender\_endpoint} & Windows 10 & \texttt{ProgramData\textbackslash{}Microsoft\textbackslash{}Windows Defender Advanced Threat Protection} & fichiers, copie brute & non \\
 & \multicolumn{4}{@{}L{12.8cm}@{}}{\itshape Présent seulement là où Defender for Endpoint est inscrit ; ses formats ne sont pas documentés : copié brut, non converti.} \\
\addlinespace
\texttt{deleted\_files} & Windows XP & \texttt{\$MFT, \$Bitmap, \$I30, clusters libres ; volumes FAT et exFAT} & enregistrements de fichier NTFS 3.1, data runs, index blocks; entrées de répertoire FAT; PE, ZIP, PDF carved & oui \\
 & \multicolumn{4}{@{}L{12.8cm}@{}}{\itshape Ce qu'un fichier supprimé garde dépend du disque : un SSD qui a reçu TRIM rend des zéros pour les clusters libérés, déclarés comme tels, jamais comme le contenu.} \\
\addlinespace
\texttt{devices} & Windows XP & \texttt{SYSTEM\textbackslash{}CurrentControlSet\textbackslash{}Enum\textbackslash{}USB, SCSI} & clés de registre & oui \\
 & Windows Vista & \texttt{…et SOFTWARE\textbackslash{}Microsoft\textbackslash{}Windows Portable Devices, propriétés des périphériques avec leurs dates (Windows 8)} & clés de registre & oui \\
\addlinespace
\texttt{dpapi\_master\_keys} & Windows XP & \texttt{\textless{}profile\textgreater{}\textbackslash{}Application Data\textbackslash{}Microsoft\textbackslash{}Protect\textbackslash{}\textless{}SID\textgreater{}} & fichier de clé maîtresse, version 2 & oui \\
 & Windows Vista & \texttt{Users\textbackslash{}\textless{}user\textgreater{}\textbackslash{}AppData\textbackslash{}Roaming\textbackslash{}Microsoft\textbackslash{}Protect\textbackslash{}\textless{}SID\textgreater{}} & fichier de clé maîtresse, version 2 & oui \\
\addlinespace
\texttt{emails} & Windows XP & \texttt{PST et OST d'Outlook, mbox de Thunderbird, dans chaque profil} & PST/OST (ANSI et Unicode), mbox & oui \\
 & \multicolumn{4}{@{}L{12.8cm}@{}}{\itshape Dépend du logiciel de messagerie installé, non de Windows.} \\
\addlinespace
\texttt{error\_reports} & Windows Vista & \texttt{ProgramData\textbackslash{}Microsoft\textbackslash{}Windows\textbackslash{}WER, Users\textbackslash{}\textless{}user\textgreater{}\textbackslash{}AppData\textbackslash{}Local\textbackslash{}Microsoft\textbackslash{}Windows\textbackslash{}WER} & Report.wer, texte UTF-16 & oui \\
 & \multicolumn{4}{@{}L{12.8cm}@{}}{\itshape Windows XP : Dr Watson (drwtsn32.log, user.dmp), non collecté.} \\
\addlinespace
\texttt{etw\_logs} & Windows Vista & \texttt{Windows\textbackslash{}System32\textbackslash{}LogFiles\textbackslash{}WMI, Windows\textbackslash{}System32\textbackslash{}WDI, Windows\textbackslash{}Logs, ProgramData\textbackslash{}Microsoft\textbackslash{}Diagnosis\textbackslash{}ETLLogs} & fichiers de trace ETL, copie brute & non \\
 & \multicolumn{4}{@{}L{12.8cm}@{}}{\itshape Copié brut : non converti (événements du noyau au format MOF, manifestes des fournisseurs).} \\
\addlinespace
\texttt{event\_log\_configuration} & Windows XP & \texttt{SYSTEM\textbackslash{}CurrentControlSet\textbackslash{}Services\textbackslash{}EventLog} & clés de registre & oui \\
 & Windows Vista & \texttt{…et SOFTWARE\textbackslash{}Microsoft\textbackslash{}Windows\textbackslash{}CurrentVersion\textbackslash{}WINEVT\textbackslash{}Channels, Control\textbackslash{}WMI\textbackslash{}Autologger} & clés de registre & oui \\
 & \multicolumn{4}{@{}L{12.8cm}@{}}{\itshape toujours présent depuis Windows XP} \\
\addlinespace
\texttt{events} & Windows XP & \texttt{Windows\textbackslash{}System32\textbackslash{}config\textbackslash{}*.evt} & EVT (enregistrements LfLe) & oui \\
 & Windows Vista & \texttt{Windows\textbackslash{}System32\textbackslash{}winevt\textbackslash{}Logs\textbackslash{}*.evtx} & EVTX, XML binaire & oui \\
 & \multicolumn{4}{@{}L{12.8cm}@{}}{\itshape toujours présent depuis Windows XP} \\
\addlinespace
\texttt{explorer\_recent\_docs} & Windows XP & \texttt{NTUSER.DAT\textbackslash{}Software\textbackslash{}Microsoft\textbackslash{}Windows\textbackslash{}CurrentVersion\textbackslash{}Explorer\textbackslash{}RecentDocs} & MRUListEx et éléments de shell & oui \\
\addlinespace
\texttt{fat\_volumes} & Windows XP & \texttt{volumes FAT12, FAT16, FAT32} & entrées de répertoire FAT & oui \\
 & Windows Vista & \texttt{…et volumes exFAT (Vista SP1 ; XP avec KB955704)} & entrées de répertoire exFAT & oui \\
\addlinespace
\texttt{firewall} & Windows XP & \texttt{SYSTEM\textbackslash{}CurrentControlSet\textbackslash{}Services\textbackslash{}SharedAccess\textbackslash{}Parameters\textbackslash{}FirewallPolicy (XP SP2)} & valeurs de registre, profils et listes de ports & oui \\
 & Windows Vista & \texttt{…\textbackslash{}FirewallPolicy\textbackslash{}FirewallRules} & valeurs de registre, règles en chaînes v2.x & oui \\
\addlinespace
\texttt{firewall\_log} & Windows XP & \texttt{Windows\textbackslash{}pfirewall.log (XP SP2), ou le LogFilePath du profil} & texte W3C, heure locale & oui \\
 & Windows Vista & \texttt{Windows\textbackslash{}System32\textbackslash{}LogFiles\textbackslash{}Firewall\textbackslash{}pfirewall.log, ou le LogFilePath du profil} & texte W3C, heure locale ; processus depuis Windows 10 & oui \\
 & \multicolumn{4}{@{}L{12.8cm}@{}}{\itshape Désactivé par défaut : présent seulement là où la journalisation du pare-feu a été activée.} \\
\addlinespace
\texttt{hosts} & Windows XP & \texttt{Windows\textbackslash{}System32\textbackslash{}drivers\textbackslash{}etc\textbackslash{}hosts} & texte & oui \\
 & \multicolumn{4}{@{}L{12.8cm}@{}}{\itshape toujours présent depuis Windows XP} \\
\addlinespace
\texttt{installed\_programs} & Windows XP & \texttt{SOFTWARE et NTUSER.DAT …\textbackslash{}CurrentVersion\textbackslash{}Uninstall, Installer\textbackslash{}UserData} & clés de registre & oui \\
 & \multicolumn{4}{@{}L{12.8cm}@{}}{\itshape toujours présent depuis Windows XP} \\
\addlinespace
\texttt{internet\_settings} & Windows XP & \texttt{NTUSER.DAT\textbackslash{}Software\textbackslash{}Microsoft\textbackslash{}Windows\textbackslash{}CurrentVersion\textbackslash{}Internet Settings} & valeurs de registre & oui \\
\addlinespace
\texttt{intune\_logs} & Windows 10 & \texttt{ProgramData\textbackslash{}Microsoft\textbackslash{}IntuneManagementExtension\textbackslash{}Logs\textbackslash{}*.log} & texte CMTrace, heure locale & oui \\
 & \multicolumn{4}{@{}L{12.8cm}@{}}{\itshape Présent seulement sur une machine gérée par Microsoft Intune, une fois l'extension de gestion installée.} \\
\addlinespace
\texttt{jump\_lists} & Windows 7 & \texttt{Users\textbackslash{}\textless{}user\textgreater{}\textbackslash{}AppData\textbackslash{}Roaming\textbackslash{}Microsoft\textbackslash{}Windows\textbackslash{}Recent\textbackslash{}AutomaticDestinations, CustomDestinations} & fichier composé OLE, DestList version 1 & oui \\
 & Windows 10 & \texttt{les mêmes dossiers} & fichier composé OLE, DestList version 3 et 4 & oui \\
 & Windows 11 & \texttt{les mêmes dossiers} & fichier composé OLE, DestList version 6 & oui \\
 & \multicolumn{4}{@{}L{12.8cm}@{}}{\itshape DestList version 6 mesurée sur le build 26200 ; les listes personnalisées (CustomDestinations) ne portent pas de version.} \\
\addlinespace
\texttt{known\_networks} & Windows Vista & \texttt{SOFTWARE\textbackslash{}Microsoft\textbackslash{}Windows NT\textbackslash{}CurrentVersion\textbackslash{}NetworkList} & clés de registre, dates SYSTEMTIME & oui \\
\addlinespace
\texttt{local\_groups} & Windows XP & \texttt{SAM\textbackslash{}SAM\textbackslash{}Domains\textbackslash{}Builtin\textbackslash{}Aliases} & enregistrements d'alias de la SAM & oui \\
 & \multicolumn{4}{@{}L{12.8cm}@{}}{\itshape toujours présent depuis Windows XP} \\
\addlinespace
\texttt{log\_file} & Windows XP & \texttt{\$LogFile} & journal NTFS, copie brute & non \\
 & Windows Vista & \texttt{…et \$Extend\textbackslash{}\$RmMetadata\textbackslash{}\$TxfLog} & journal NTFS transactionnel, copie brute & non \\
 & \multicolumn{4}{@{}L{12.8cm}@{}}{\itshape Copié brut sur toutes les versions : non converti.} \\
\addlinespace
\texttt{mft} & Windows XP & \texttt{\$MFT} & enregistrements de fichier NTFS 3.1 & oui \\
\addlinespace
\texttt{mounted\_devices} & Windows XP & \texttt{SYSTEM\textbackslash{}MountedDevices} & valeurs de registre & oui \\
 & \multicolumn{4}{@{}L{12.8cm}@{}}{\itshape toujours présent depuis Windows XP} \\
\addlinespace
\texttt{mount\_points} & Windows XP & \texttt{NTUSER.DAT\textbackslash{}Software\textbackslash{}Microsoft\textbackslash{}Windows\textbackslash{}CurrentVersion\textbackslash{}Explorer\textbackslash{}MountPoints2} & clés de registre & oui \\
\addlinespace
\texttt{mru} & Windows XP & \texttt{NTUSER.DAT …\textbackslash{}Explorer\textbackslash{}ComDlg32\textbackslash{}OpenSaveMRU, LastVisitedMRU} & MRUList et chaînes & oui \\
 & Windows Vista & \texttt{NTUSER.DAT …\textbackslash{}Explorer\textbackslash{}ComDlg32\textbackslash{}OpenSavePidlMRU, LastVisitedPidlMRU} & MRUListEx et éléments de shell & oui \\
\addlinespace
\texttt{mru\_apps} & Windows XP & \texttt{NTUSER.DAT …\textbackslash{}Explorer\textbackslash{}ComDlg32\textbackslash{}LastVisitedMRU} & MRUList et chaînes & oui \\
 & Windows Vista & \texttt{NTUSER.DAT …\textbackslash{}Explorer\textbackslash{}ComDlg32\textbackslash{}LastVisitedPidlMRU} & MRUListEx et éléments de shell & oui \\
\addlinespace
\texttt{muicache} & Windows XP & \texttt{NTUSER.DAT\textbackslash{}Software\textbackslash{}Microsoft\textbackslash{}Windows\textbackslash{}ShellNoRoam\textbackslash{}MUICache} & valeurs de registre & oui \\
 & Windows Vista & \texttt{UsrClass.dat\textbackslash{}Local Settings\textbackslash{}Software\textbackslash{}Microsoft\textbackslash{}Windows\textbackslash{}Shell\textbackslash{}MuiCache} & valeurs de registre & oui \\
\addlinespace
\texttt{network} & Windows XP & \texttt{en direct : tables TCP/IP du système en marche} & observation & oui \\
\addlinespace
\texttt{network\_configuration} & Windows XP & \texttt{SYSTEM\textbackslash{}CurrentControlSet\textbackslash{}Services\textbackslash{}Tcpip\textbackslash{}Parameters\textbackslash{}Interfaces, Control\textbackslash{}Network} & valeurs de registre & oui \\
 & Windows 10 & \texttt{…et Control\textbackslash{}NetworkSetup2\textbackslash{}Interfaces\textbackslash{}\textless{}guid\textgreater{}\textbackslash{}Kernel : adresses matérielles} & valeurs de registre & oui \\
 & \multicolumn{4}{@{}L{12.8cm}@{}}{\itshape toujours présent depuis Windows XP ; Avant Windows 10, les adresses matérielles des cartes ne sont dans aucune clé documentée des ruches : non données.} \\
\addlinespace
\texttt{network\_shares} & Windows XP & \texttt{SYSTEM\textbackslash{}CurrentControlSet\textbackslash{}Services\textbackslash{}LanmanServer\textbackslash{}Shares} & valeurs de registre, descripteurs de sécurité & oui \\
\addlinespace
\texttt{ntfs\_volumes} & Windows XP & \texttt{\$Boot, \$Extend\textbackslash{}\$UsnJrnl:\$Max} & secteur de démarrage NTFS, en-tête du journal & oui \\
\addlinespace
\texttt{office\_trust\_records} & Windows XP & \texttt{NTUSER.DAT\textbackslash{}Software\textbackslash{}Microsoft\textbackslash{}Office\textbackslash{}\textless{}version\textgreater{}\textbackslash{}\textless{}program\textgreater{}\textbackslash{}Security\textbackslash{}Trusted Documents\textbackslash{}TrustRecords} & valeurs de registre & oui \\
 & \multicolumn{4}{@{}L{12.8cm}@{}}{\itshape Dépend d'Office (2010 et suivants), non de Windows.} \\
\addlinespace
\texttt{onedrive} & Windows 7 & \texttt{Users\textbackslash{}\textless{}user\textgreater{}\textbackslash{}AppData\textbackslash{}Local\textbackslash{}Microsoft\textbackslash{}OneDrive\textbackslash{}logs et settings} & journaux ODL, base de synchronisation, copie brute & non \\
 & \multicolumn{4}{@{}L{12.8cm}@{}}{\itshape Présent seulement là où OneDrive a tourné ; journaux de la synchronisation des listes SharePoint (ListSync) écartés. Copié brut : non converti.} \\
\addlinespace
\texttt{operating\_system} & Windows XP & \texttt{SOFTWARE\textbackslash{}Microsoft\textbackslash{}Windows NT\textbackslash{}CurrentVersion, SYSTEM\textbackslash{}…\textbackslash{}Control} & valeurs de registre & oui \\
 & \multicolumn{4}{@{}L{12.8cm}@{}}{\itshape toujours présent depuis Windows XP} \\
\addlinespace
\texttt{persistence} & Windows XP & \texttt{SOFTWARE, SYSTEM et NTUSER.DAT : IFEO, AppInit\_DLLs, Winlogon, LSA, Active Setup…} & valeurs de registre & oui \\
 & \multicolumn{4}{@{}L{12.8cm}@{}}{\itshape toujours présent depuis Windows XP} \\
\addlinespace
\texttt{powershell\_history} & Windows 7 & \texttt{Users\textbackslash{}\textless{}user\textgreater{}\textbackslash{}AppData\textbackslash{}Roaming\textbackslash{}Microsoft\textbackslash{}Windows\textbackslash{}PowerShell\textbackslash{}PSReadLine\textbackslash{}ConsoleHost\_history.txt} & texte UTF-8 & oui \\
 & \multicolumn{4}{@{}L{12.8cm}@{}}{\itshape Écrit par PSReadLine : par défaut depuis Windows 10 ; sous Windows 7, 8 et 8.1 une fois PowerShell 3 ou plus et PSReadLine installés. Vista et XP s'arrêtent à PowerShell 2, sur lequel PSReadLine ne fonctionne pas.} \\
\addlinespace
\texttt{powershell\_profiles} & Windows XP & \texttt{Windows\textbackslash{}System32\textbackslash{}WindowsPowerShell\textbackslash{}v1.0\textbackslash{}*profile.ps1, Windows\textbackslash{}SysWOW64\textbackslash{}WindowsPowerShell\textbackslash{}v1.0, Users\textbackslash{}\textless{}user\textgreater{}\textbackslash{}Documents\textbackslash{}WindowsPowerShell} & script PowerShell & oui \\
 & Windows 7 & \texttt{…et Program Files\textbackslash{}PowerShell\textbackslash{}\textless{}version\textgreater{}, Users\textbackslash{}\textless{}user\textgreater{}\textbackslash{}Documents\textbackslash{}PowerShell} & script PowerShell & oui \\
 & \multicolumn{4}{@{}L{12.8cm}@{}}{\itshape Windows XP et Vista : seulement une fois PowerShell installé, absent par défaut. Les documents d'un utilisateur aussi là où User Shell Folders les déplace.} \\
\addlinespace
\texttt{prefetch} & Windows XP & \texttt{Windows\textbackslash{}Prefetch\textbackslash{}*.pf} & SCCA version 17 & oui \\
 & Windows Vista & \texttt{Windows\textbackslash{}Prefetch\textbackslash{}*.pf} & SCCA version 23 & oui \\
 & Windows 8 & \texttt{Windows\textbackslash{}Prefetch\textbackslash{}*.pf} & SCCA version 26 & oui \\
 & Windows 10 & \texttt{Windows\textbackslash{}Prefetch\textbackslash{}*.pf} & SCCA version 30 et 31, compressé (MAM) & oui \\
 & \multicolumn{4}{@{}L{12.8cm}@{}}{\itshape Désactivé par défaut sur les éditions serveur.} \\
\addlinespace
\texttt{processes} & Windows XP & \texttt{en direct : processus du système en marche} & observation & oui \\
\addlinespace
\texttt{program\_usage} & Windows XP & \texttt{NTUSER.DAT …\textbackslash{}AppCompatFlags\textbackslash{}Layers} & valeurs de registre & oui \\
 & Windows Vista & \texttt{…et AppCompatFlags\textbackslash{}Compatibility Assistant\textbackslash{}Store, Explorer\textbackslash{}Taskband} & valeurs de registre & oui \\
 & Windows 10 & \texttt{…et Explorer\textbackslash{}FeatureUsage, Search\textbackslash{}JumplistData} & valeurs de registre & oui \\
\addlinespace
\texttt{rdp\_client} & Windows XP & \texttt{NTUSER.DAT\textbackslash{}Software\textbackslash{}Microsoft\textbackslash{}Terminal Server Client; …\textbackslash{}Terminal Server Client\textbackslash{}Cache\textbackslash{}bcache*.bmc} & valeurs de registre; cache d'images BMC & oui \\
 & Windows 7 & \texttt{…et Cache\textbackslash{}Cache????.bin} & cache d'images BIN & oui \\
 & \multicolumn{4}{@{}L{12.8cm}@{}}{\itshape bcache22.bmc décodé (16 bits par pixel, mesuré sous XP SP1) ; un autre .bmc est compté, ses tuiles non exportées.} \\
\addlinespace
\texttt{rdp\_server} & Windows Vista & \texttt{Security.evtx, journaux TerminalServices-LocalSessionManager et RemoteConnectionManager} & événements EVTX & oui \\
 & \multicolumn{4}{@{}L{12.8cm}@{}}{\itshape Windows XP et Server 2003 : événements 528 et 682 de Security.evt, non convertis.} \\
\addlinespace
\texttt{recent\_documents} & Windows XP & \texttt{\textless{}profile\textgreater{}\textbackslash{}Recent\textbackslash{}*.lnk} & raccourci & oui \\
 & Windows Vista & \texttt{Users\textbackslash{}\textless{}user\textgreater{}\textbackslash{}AppData\textbackslash{}Roaming\textbackslash{}Microsoft\textbackslash{}Windows\textbackslash{}Recent\textbackslash{}*.lnk} & raccourci & oui \\
\addlinespace
\texttt{recycle\_bin} & Windows XP & \texttt{RECYCLER\textbackslash{}\textless{}SID\textgreater{}\textbackslash{}INFO2} & enregistrements INFO2, version 5 & oui \\
 & Windows Vista & \texttt{\$Recycle.Bin\textbackslash{}\textless{}SID\textgreater{}\textbackslash{}\$I*} & \$I version 1 & oui \\
 & Windows 10 & \texttt{\$Recycle.Bin\textbackslash{}\textless{}SID\textgreater{}\textbackslash{}\$I*} & \$I version 2 & oui \\
\addlinespace
\texttt{registry\_backups} & Windows XP & \texttt{Windows\textbackslash{}repair} & ruches & non \\
 & Windows Vista & \texttt{Windows\textbackslash{}System32\textbackslash{}config\textbackslash{}RegBack} & ruches, copie brute & non \\
 & \multicolumn{4}{@{}L{12.8cm}@{}}{\itshape Vide depuis Windows 10 1803, sauf si la tâche de sauvegarde a été réactivée. Copié brut : non converti.} \\
\addlinespace
\texttt{root\_certificates} & Windows XP & \texttt{SOFTWARE et NTUSER.DAT …\textbackslash{}SystemCertificates\textbackslash{}Root, AuthRoot} & certificats (blobs) & oui \\
 & \multicolumn{4}{@{}L{12.8cm}@{}}{\itshape toujours présent depuis Windows XP} \\
\addlinespace
\texttt{run\_keys} & Windows XP & \texttt{SOFTWARE et NTUSER.DAT …\textbackslash{}CurrentVersion\textbackslash{}Run, RunOnce} & valeurs de registre & oui \\
 & \multicolumn{4}{@{}L{12.8cm}@{}}{\itshape toujours présent depuis Windows XP} \\
\addlinespace
\texttt{scheduled\_tasks} & Windows XP & \texttt{Windows\textbackslash{}Tasks\textbackslash{}*.job} & binaire JOB & oui \\
 & Windows Vista & \texttt{Windows\textbackslash{}System32\textbackslash{}Tasks, SOFTWARE …\textbackslash{}Schedule\textbackslash{}TaskCache} & XML, valeurs de registre & oui \\
 & \multicolumn{4}{@{}L{12.8cm}@{}}{\itshape toujours présent depuis Windows Vista} \\
\addlinespace
\texttt{search\_index} & Windows Vista & \texttt{ProgramData\textbackslash{}Microsoft\textbackslash{}Search\textbackslash{}Data\textbackslash{}Applications\textbackslash{}Windows\textbackslash{}Windows.edb} & ESE & oui \\
 & Windows 11 & \texttt{…\textbackslash{}Windows\textbackslash{}Windows.db, Windows-gather.db} & SQLite & oui \\
 & \multicolumn{4}{@{}L{12.8cm}@{}}{\itshape Windows.edb converti, avec son texte compressé sur 7 bits (mesuré sous Windows 8.1), et la disposition de Windows 7 : ses propres noms de table et de colonnes, ses nombres gros-boutistes, son texte brouillé (mesurés) ; une valeur compressée autrement (XPRESS) est signalée, non décodée. Windows XP : Windows Desktop Search, une installation facultative.} \\
\addlinespace
\texttt{security} & Windows XP & \texttt{ruche SECURITY : Policy\textbackslash{}Secrets, Cache, PolAdtEv} & valeurs de registre & oui \\
 & \multicolumn{4}{@{}L{12.8cm}@{}}{\itshape toujours présent depuis Windows XP} \\
\addlinespace
\texttt{security\_descriptors} & Windows XP & \texttt{\$Secure:\$SDS} & flux des descripteurs de sécurité NTFS & oui \\
\addlinespace
\texttt{services} & Windows XP & \texttt{SYSTEM\textbackslash{}CurrentControlSet\textbackslash{}Services} & clés de registre & oui \\
 & \multicolumn{4}{@{}L{12.8cm}@{}}{\itshape toujours présent depuis Windows XP} \\
\addlinespace
\texttt{sessions} & Windows XP & \texttt{en direct : sessions ouvertes sur le système en marche} & observation & oui \\
\addlinespace
\texttt{setupapi} & Windows XP & \texttt{Windows\textbackslash{}setupapi.log} & texte, disposition de XP & oui \\
 & Windows Vista & \texttt{Windows\textbackslash{}INF\textbackslash{}setupapi.dev.log} & texte, sections & oui \\
 & \multicolumn{4}{@{}L{12.8cm}@{}}{\itshape toujours présent depuis Windows XP} \\
\addlinespace
\texttt{shadow\_copies} & Windows Vista & \texttt{System Volume Information\textbackslash{}\{3808876b-c176-4e48-b7ae-04046e6cc752\}, \{GUID\}\{3808876b-…\}} & catalogue et magasins VSS, blocs de 16 Kio & oui \\
 & \multicolumn{4}{@{}L{12.8cm}@{}}{\itshape Windows XP ne garde aucun cliché persistant : ses points de restauration sont des dossiers à part. Server 2003 ne catalogue aucun magasin (pas d'entrée de type 3) : ses clichés sont listés, pas lus.} \\
\addlinespace
\texttt{shellbags} & Windows XP & \texttt{NTUSER.DAT\textbackslash{}Software\textbackslash{}Microsoft\textbackslash{}Windows\textbackslash{}Shell et ShellNoRoam} & éléments de shell de BagMRU & oui \\
 & Windows Vista & \texttt{UsrClass.dat\textbackslash{}Local Settings\textbackslash{}Software\textbackslash{}Microsoft\textbackslash{}Windows\textbackslash{}Shell, et NTUSER.DAT} & éléments de shell de BagMRU & oui \\
\addlinespace
\texttt{shimcache} & Windows XP & \texttt{SYSTEM\textbackslash{}…\textbackslash{}Control\textbackslash{}Session Manager\textbackslash{}AppCompatibility\textbackslash{}AppCompatCache} & en-tête 0xDEADBEEF, 96 entrées & oui \\
 & Windows XP & \texttt{SYSTEM\textbackslash{}…\textbackslash{}Control\textbackslash{}Session Manager\textbackslash{}AppCompatCache\textbackslash{}AppCompatCache} & en-tête 0xBADC0FFE, taille du fichier & non \\
 & Windows Vista & \texttt{la même valeur} & en-tête 0xBADC0FFE, entrées de 24 ou 32 octets & oui \\
 & Windows 7 & \texttt{la même valeur} & en-tête 0xBADC0FEE & oui \\
 & Windows 8 & \texttt{la même valeur} & en-tête 0x80, entrées ''00ts'' & oui \\
 & Windows 8.1 & \texttt{la même valeur} & en-tête de 0x80 octets commençant par 0, entrées ''10ts'' & oui \\
 & Windows 10 & \texttt{la même valeur} & en-tête 0x30 ou 0x34, entrées ''10ts'' & oui \\
 & \multicolumn{4}{@{}L{12.8cm}@{}}{\itshape toujours présent depuis Windows XP ; Windows 8 et 8.1 ont un en-tête de 0x80 octets : ils se distinguent par la marque de leur première entrée, qui tient ici lieu de signature. Windows 8.1 mesuré sur le build 9600. Server 2003 et Vista partagent leur signature : la version choisit la forme ; Vista mesuré sur SP2 x86, Server 2003 pas encore.} \\
\addlinespace
\texttt{shortcuts} & Windows XP & \texttt{menu Démarrer, bureaux, Lancement rapide de chaque profil et de All Users} & raccourci & oui \\
\addlinespace
\texttt{srum} & Windows 8 & \texttt{Windows\textbackslash{}System32\textbackslash{}sru\textbackslash{}SRUDB.dat} & ESE & oui \\
 & \multicolumn{4}{@{}L{12.8cm}@{}}{\itshape toujours présent depuis Windows 8} \\
\addlinespace
\texttt{startup\_folders} & Windows XP & \texttt{Start Menu\textbackslash{}Programs\textbackslash{}Startup de chaque profil et de All Users} & fichiers & oui \\
\addlinespace
\texttt{sysmon} & Windows 7 & \texttt{SYSTEM\textbackslash{}CurrentControlSet\textbackslash{}Services\textbackslash{}\textless{}pilote de Sysmon\textgreater{}\textbackslash{}Parameters, son ArchiveDirectory} & clés du registre ; Rules telles que stockées & oui \\
 & \multicolumn{4}{@{}L{12.8cm}@{}}{\itshape Absent de Windows : présent seulement là où Sysmon a été installé. Le format binaire de ses règles n'est pas documenté : donné tel que stocké, avec son empreinte.} \\
\addlinespace
\texttt{thumbcache} & Windows XP & \texttt{Thumbs.db dans chaque dossier} & fichier composé OLE & oui \\
 & Windows Vista & \texttt{Users\textbackslash{}\textless{}user\textgreater{}\textbackslash{}AppData\textbackslash{}Local\textbackslash{}Microsoft\textbackslash{}Windows\textbackslash{}Explorer\textbackslash{}thumbcache\_*.db} & CMMM, disposition de Vista & oui \\
 & Windows 7 & \texttt{le même dossier} & CMMM, disposition de Windows 7 & oui \\
 & Windows 8 & \texttt{le même dossier, avec iconcache\_*.db} & CMMM, disposition de Windows 8 et suivants & oui \\
\addlinespace
\texttt{typed\_paths} & Windows Vista & \texttt{NTUSER.DAT\textbackslash{}Software\textbackslash{}Microsoft\textbackslash{}Windows\textbackslash{}CurrentVersion\textbackslash{}Explorer\textbackslash{}TypedPaths} & valeurs de registre & oui \\
\addlinespace
\texttt{usbstor} & Windows XP & \texttt{SYSTEM\textbackslash{}CurrentControlSet\textbackslash{}Enum\textbackslash{}USBSTOR} & clés de registre & oui \\
 & Windows 8 & \texttt{…avec les dates des propriétés du périphérique (première installation, dernière connexion, dernier retrait)} & clés de registre, FILETIME & oui \\
\addlinespace
\texttt{userassist} & Windows XP & \texttt{NTUSER.DAT …\textbackslash{}Explorer\textbackslash{}UserAssist\textbackslash{}\{5E6AB780-…\}, \{75048700-…\}} & noms en ROT13, valeurs de 16 octets & oui \\
 & Windows 7 & \texttt{NTUSER.DAT …\textbackslash{}Explorer\textbackslash{}UserAssist\textbackslash{}\{CEBFF5CD-…\}, \{F4E57C4B-…\}} & noms en ROT13, valeurs de 72 octets & oui \\
\addlinespace
\texttt{users} & Windows XP & \texttt{SAM\textbackslash{}SAM\textbackslash{}Domains\textbackslash{}Account\textbackslash{}Users} & enregistrements F et V, condensats sous RC4 (clé de démarrage) & oui \\
 & Windows 10 & \texttt{les mêmes clés} & enregistrements F et V, condensats sous AES & oui \\
 & \multicolumn{4}{@{}L{12.8cm}@{}}{\itshape toujours présent depuis Windows XP} \\
\addlinespace
\texttt{user\_shell\_folders} & Windows XP & \texttt{NTUSER.DAT …\textbackslash{}Explorer\textbackslash{}User Shell Folders} & valeurs de registre & oui \\
\addlinespace
\texttt{usn\_journal} & Windows XP & \texttt{\$Extend\textbackslash{}\$UsnJrnl:\$J} & USN\_RECORD\_V2 & oui \\
 & Windows 8 & \texttt{le même flux} & USN\_RECORD\_V2, V3 (références de 128 bits) & oui \\
 & Windows 8.1 & \texttt{le même flux} & USN\_RECORD\_V2, V3, V4 (suivi des plages) & oui \\
 & \multicolumn{4}{@{}L{12.8cm}@{}}{\itshape Le journal n'existe que là où il a été créé : toujours sur le volume système à partir de Vista.} \\
\addlinespace
\texttt{web\_client} & Windows XP & \texttt{Internet Explorer 5 à 9 : index.dat de l'historique, du cache, des cookies, des téléchargements} & MSIE Cache File 5.2 & oui \\
 & Windows 8 & \texttt{Internet Explorer 10 et 11, Edge Legacy : …\textbackslash{}WebCache\textbackslash{}WebCacheV01.dat} & ESE & oui \\
 & Windows 10 & \texttt{…et navigateurs Chromium, Firefox : History, places.sqlite, fichiers de session} & SQLite, SNSS, mozLz4 & oui \\
 & \multicolumn{4}{@{}L{12.8cm}@{}}{\itshape Chromium et Firefox dépendent du navigateur installé, non de Windows.} \\
\addlinespace
\texttt{wmi} & Windows XP & \texttt{Windows\textbackslash{}System32\textbackslash{}wbem\textbackslash{}Repository\textbackslash{}FS} & dépôt CIM de Windows XP (INDEX.BTR, OBJECTS.DATA, MAPPING, dans FS) & oui \\
 & Windows Vista & \texttt{Windows\textbackslash{}System32\textbackslash{}wbem\textbackslash{}Repository} & dépôt CIM (INDEX.BTR, OBJECTS.DATA, MAPPING) & oui \\
 & \multicolumn{4}{@{}L{12.8cm}@{}}{\itshape toujours présent depuis Windows XP ; Windows XP range le dépôt dans un sous-dossier FS, avec des condensats MD5 et des entrées de page de 4 octets (mesuré, SP1 et SP3).} \\
\addlinespace
\texttt{word\_wheel\_query} & Windows 7 & \texttt{NTUSER.DAT …\textbackslash{}Explorer\textbackslash{}WordWheelQuery} & MRUListEx et chaînes UTF-16 & oui \\
 & \multicolumn{4}{@{}L{12.8cm}@{}}{\itshape Windows XP : Search Assistant\textbackslash{}ACMru, non collecté.} \\
\addlinespace
\texttt{zone\_identifiers} & Windows XP & \texttt{\textless{}fichier\textgreater{}:Zone.Identifier (XP SP2)} & flux de données alterné, texte INI & oui \\
\bottomrule
\end{longtable}
}

\noindent{\small Une ligne par forme de chaque artefact : le premier Windows qui la porte, son emplacement, son format, et si WAC convertit cette forme (une forme non convertie est au plus collectée brute).}


\section{Couverture de MITRE ATT\&CK}\label{sec:attack}

ATT\&CK décrit, technique par technique, comment la détecter : des
analytiques, chacune nommant les sources dont elle a besoin — un composant de
données (« Process Creation ») et la source de journal qui le porte
(« WinEventLog:Sysmon EventCode=1 »). C'est une liste tenue hors de \wac{},
contre laquelle sa couverture se mesure. Les tableaux ci-dessous en sont
produits : d'abord ce que \wac{} fait de chaque source de journal, puis les
artefacts qui gardent sur le disque la trace de chaque composant de données
hors des journaux, enfin, pour chaque technique Windows, si un journal que
\wac{} prend ou un artefact du disque en porte une trace, et les sources qu'il
ne prend pas, avec leur raison. Une trace portée n'est pas une détection : la rechercher dans les
fichiers JSON relève de l'analyse.

% Generated by doc/attack-coverage.py from ATT&CK Enterprise 20.0 — do not edit.
{\small
\begin{longtable}{@{}L{5.2cm}L{9.4cm}@{}}
\toprule
\textbf{Source de journal} & \textbf{WAC} \\
\midrule\endhead
\texttt{wineventlog:sysmon…} & \texttt{events}, \texttt{sysmon} \\
\texttt{wineventlog:iis…} & rôle serveur (IIS, AD FS, contrôleur de domaine) : version 4.0 \\
\texttt{wineventlog:microsoft-iis…} & rôle serveur (IIS, AD FS, contrôleur de domaine) : version 4.0 \\
\texttt{wineventlog:adfs…} & rôle serveur (IIS, AD FS, contrôleur de domaine) : version 4.0 \\
\texttt{wineventlog:…} & \texttt{events} \\
\texttt{autoruns:registryscan…} & \texttt{persistence}, \texttt{run\_keys}, \texttt{services}, \texttt{scheduled\_tasks}, \texttt{startup\_folders} \\
\texttt{windows firewall log…} & \texttt{firewall\_log} \\
\texttt{applicationlog:intune…} & \texttt{intune\_logs}, \texttt{events} \\
\texttt{applicationlog:iis…} & rôle serveur (IIS, AD FS, contrôleur de domaine) : version 4.0 \\
\texttt{applicationlog:callrecords…} & service en nuage : dans les journaux du locataire ; le côté de la machine est pris (cloud\_identity, onedrive, intune\_logs) \\
\texttt{etw…} & télémétrie vivante (ETW, appels d'API, compteurs) : lue par un consommateur à l'instant ; les fichiers de trace que Windows écrit sur son disque sont pris bruts (etw\_logs) \\
\texttt{windows:perfmon…} & télémétrie vivante (ETW, appels d'API, compteurs) : lue par un consommateur à l'instant ; les fichiers de trace que Windows écrit sur son disque sont pris bruts (etw\_logs) \\
\texttt{windows:osquery…} & télémétrie vivante (ETW, appels d'API, compteurs) : lue par un consommateur à l'instant ; les fichiers de trace que Windows écrit sur son disque sont pris bruts (etw\_logs) \\
\texttt{edr…} & télémétrie propre à un EDR : dans sa console ; le dossier local de Defender for Endpoint est pris brut (defender\_endpoint) \\
\texttt{m365defender…} & télémétrie propre à un EDR : dans sa console ; le dossier local de Defender for Endpoint est pris brut (defender\_endpoint) \\
\texttt{m365:defender…} & télémétrie propre à un EDR : dans sa console ; le dossier local de Defender for Endpoint est pris brut (defender\_endpoint) \\
\texttt{m365:…} & service en nuage : dans les journaux du locataire ; le côté de la machine est pris (cloud\_identity, onedrive, intune\_logs) \\
\texttt{azure:…} & service en nuage : dans les journaux du locataire ; le côté de la machine est pris (cloud\_identity, onedrive, intune\_logs) \\
\texttt{aws:…} & service en nuage : dans les journaux du locataire ; le côté de la machine est pris (cloud\_identity, onedrive, intune\_logs) \\
\texttt{nsm:…} & capture réseau : sur le réseau, pas sur la machine \\
\texttt{network traffic…} & capture réseau : sur le réseau, pas sur la machine \\
\texttt{dns:query…} & capture réseau : sur le réseau, pas sur la machine \\
\texttt{firmware:…} & micrologiciel : non lu, WAC lit le disque \\
\bottomrule
\end{longtable}

\begin{longtable}{@{}L{5.2cm}L{9.4cm}@{}}
\toprule
\textbf{Composant de données} & \textbf{Sur le disque, hors des journaux} \\
\midrule\endhead
Process Creation & \texttt{prefetch}, \texttt{amcache}, \texttt{shimcache}, \texttt{bam}, \texttt{userassist}, \texttt{muicache}, \texttt{srum}, \texttt{program\_usage} \\
Process Access & — \\
Process Modification & — \\
Process Metadata & \texttt{processes} \\
Process Termination & — \\
Command Execution & \texttt{powershell\_history}, \texttt{prefetch} \\
Script Execution & \texttt{powershell\_history}, \texttt{powershell\_profiles} \\
OS API Execution & — \\
Module Load & — \\
File Creation & \texttt{mft}, \texttt{usn\_journal}, \texttt{zone\_identifiers} \\
File Modification & \texttt{mft}, \texttt{usn\_journal}, \texttt{log\_file} \\
File Deletion & \texttt{usn\_journal}, \texttt{deleted\_files}, \texttt{recycle\_bin} \\
File Access & \texttt{jump\_lists}, \texttt{recent\_documents}, \texttt{shortcuts}, \texttt{shellbags}, \texttt{explorer\_recent\_docs} \\
File Metadata & \texttt{mft}, \texttt{security\_descriptors}, \texttt{zone\_identifiers} \\
Network Connection Creation & \texttt{srum}, \texttt{network}, \texttt{firewall\_log}, \texttt{rdp\_client}, \texttt{bits}, \texttt{web\_client} \\
Network Traffic Content & — \\
Network Traffic Flow & \texttt{srum}, \texttt{firewall\_log} \\
Network Share Access & \texttt{mount\_points}, \texttt{network\_shares}, \texttt{shellbags} \\
Windows Registry Key Modification & \texttt{run\_keys}, \texttt{persistence}, \texttt{services}, \texttt{registry\_backups}, \texttt{shadow\_copies} \\
Windows Registry Key Creation & \texttt{run\_keys}, \texttt{persistence}, \texttt{services}, \texttt{registry\_backups}, \texttt{shadow\_copies} \\
Windows Registry Key Access & — \\
Service Creation & \texttt{services} \\
Service Modification & \texttt{services} \\
Service Metadata & \texttt{services} \\
Driver Load & \texttt{services} \\
Scheduled Job Creation & \texttt{scheduled\_tasks} \\
Scheduled Job Modification & \texttt{scheduled\_tasks} \\
Scheduled Job Metadata & \texttt{scheduled\_tasks} \\
WMI Creation & \texttt{wmi} \\
Logon Session Creation & \texttt{users}, \texttt{rdp\_server} \\
Logon Session Metadata & \texttt{users}, \texttt{sessions}, \texttt{rdp\_server} \\
User Account Authentication & \texttt{users} \\
User Account Creation & \texttt{users} \\
User Account Modification & \texttt{users}, \texttt{local\_groups} \\
User Account Deletion & — \\
User Account Metadata & \texttt{users}, \texttt{local\_groups}, \texttt{security} \\
Group Enumeration & — \\
Firewall Rule Modification & \texttt{firewall} \\
Drive Creation & \texttt{usbstor}, \texttt{mounted\_devices}, \texttt{devices}, \texttt{setupapi}, \texttt{mount\_points} \\
Drive Access & \texttt{mount\_points}, \texttt{shellbags}, \texttt{fat\_volumes} \\
Drive Modification & \texttt{ntfs\_volumes}, \texttt{boot\_records} \\
Volume Creation & \texttt{mounted\_devices}, \texttt{ntfs\_volumes}, \texttt{shadow\_copies} \\
Host Status & — \\
Named Pipe Metadata & — \\
Application Log Content & \texttt{defender}, \texttt{setupapi}, \texttt{error\_reports}, \texttt{intune\_logs}, \texttt{onedrive} \\
Web Credential Creation & \texttt{web\_client}, \texttt{dpapi\_master\_keys}, \texttt{cloud\_identity} \\
Firmware Modification & — \\
Active Directory Object Modification & — \\
Active Directory Object Access & — \\
Active Directory Object Creation & — \\
Active Directory Object Deletion & — \\
Active Directory Credential Request & — \\
Cloud Service Metadata & \texttt{cloud\_identity} \\
Cloud Storage Access & \texttt{onedrive} \\
Instance Creation & — \\
Instance Start & — \\
Domain Registration & — \\
\bottomrule
\end{longtable}

\begin{longtable}{@{}L{1.9cm}L{6.3cm}L{1.6cm}L{1.6cm}L{3.2cm}@{}}
\toprule
\textbf{ID} & \textbf{Technique} & \textbf{Journaux} & \textbf{Disque} & \textbf{Sources non prises} \\
\midrule\endhead
T1001 & Data Obfuscation & — & — & NSM \\
T1001.001 & Junk Data & oui & oui & NSM \\
T1001.002 & Steganography & oui & oui & NSM \\
T1001.003 & Protocol or Service Impersonation & oui & oui & NSM \\
T1003 & OS Credential Dumping & oui & oui &  \\
T1003.001 & LSASS Memory & oui & oui &  \\
T1003.002 & Security Account Manager & oui & oui &  \\
T1003.003 & NTDS & oui & oui &  \\
T1003.004 & LSA Secrets & oui & oui &  \\
T1003.005 & Cached Domain Credentials & oui & oui &  \\
T1003.006 & DCSync & oui & — & NSM \\
T1005 & Data from Local System & oui & oui &  \\
T1006 & Direct Volume Access & oui & oui &  \\
T1007 & System Service Discovery & oui & oui &  \\
T1008 & Fallback Channels & oui & oui & NSM \\
T1010 & Application Window Discovery & oui & oui &  \\
T1011 & Exfiltration Over Other Network Medium & oui & oui &  \\
T1011.001 & Exfiltration Over Bluetooth & oui & oui &  \\
T1012 & Query Registry & oui & oui &  \\
T1014 & Rootkit & oui & oui &  \\
T1016 & System Network Configuration Discovery & oui & oui &  \\
T1016.001 & Internet Connection Discovery & oui & oui &  \\
T1016.002 & Wi-Fi Discovery & oui & oui &  \\
T1018 & Remote System Discovery & oui & oui &  \\
T1020 & Automated Exfiltration & oui & oui &  \\
T1021 & Remote Services & oui & oui &  \\
T1021.001 & Remote Desktop Protocol & oui & oui &  \\
T1021.002 & SMB/Windows Admin Shares & oui & oui &  \\
T1021.003 & Distributed Component Object Model & oui & oui &  \\
T1021.004 & SSH* & oui & — &  \\
T1021.005 & VNC & oui & oui & NSM \\
T1021.006 & Windows Remote Management & oui & oui & NSM \\
T1025 & Data from Removable Media & oui & oui &  \\
T1027 & Obfuscated Files or Information & oui & oui &  \\
T1027.001 & Binary Padding & oui & oui &  \\
T1027.002 & Software Packing & oui & oui &  \\
T1027.003 & Steganography & oui & oui &  \\
T1027.004 & Compile After Delivery & oui & oui &  \\
T1027.005 & Indicator Removal from Tools & oui & oui &  \\
T1027.006 & HTML Smuggling & oui & oui & EDR, NSM \\
T1027.007 & Dynamic API Resolution & oui & oui & ETW \\
T1027.008 & Stripped Payloads & oui & oui & EDR \\
T1027.009 & Embedded Payloads & oui & oui & EDR \\
T1027.010 & Command Obfuscation & oui & oui &  \\
T1027.011 & Fileless Storage & oui & oui &  \\
T1027.012 & LNK Icon Smuggling & oui & oui &  \\
T1027.013 & Encrypted/Encoded File & oui & oui &  \\
T1027.014 & Polymorphic Code & oui & oui &  \\
T1027.015 & Compression & oui & oui &  \\
T1027.016 & Junk Code Insertion & oui & oui &  \\
T1027.017 & SVG Smuggling & oui & oui &  \\
T1027.018 & Invisible Unicode & oui & oui &  \\
T1029 & Scheduled Transfer & oui & oui &  \\
T1030 & Data Transfer Size Limits & oui & oui & NSM \\
T1033 & System Owner/User Discovery & oui & oui &  \\
T1036 & Masquerading & oui & oui &  \\
T1036.001 & Invalid Code Signature & oui & oui &  \\
T1036.002 & Right-to-Left Override & oui & oui &  \\
T1036.003 & Rename Legitimate Utilities & oui & oui & EDR \\
T1036.004 & Masquerade Task or Service & oui & oui &  \\
T1036.005 & Match Legitimate Resource Name or Location & oui & oui &  \\
T1036.007 & Double File Extension & oui & oui &  \\
T1036.008 & Masquerade File Type & oui & oui &  \\
T1036.010 & Masquerade Account Name & oui & oui & ETW \\
T1036.012 & Browser Fingerprint & oui & oui & ETW, NSM \\
T1037 & Boot or Logon Initialization Scripts & oui & oui &  \\
T1037.001 & Logon Script (Windows) & oui & oui &  \\
T1037.003 & Network Logon Script & oui & oui &  \\
T1039 & Data from Network Shared Drive & oui & oui &  \\
T1040 & Network Sniffing & oui & oui &  \\
T1041 & Exfiltration Over C2 Channel & oui & oui & NSM \\
T1046 & Network Service Discovery & oui & oui &  \\
T1047 & Windows Management Instrumentation & oui & oui &  \\
T1048 & Exfiltration Over Alternative Protocol & oui & oui &  \\
T1048.001 & Exfiltration Over Symmetric Encrypted Non-C2 Protocol & oui & oui &  \\
T1048.002 & Exfiltration Over Asymmetric Encrypted Non-C2 Protocol & oui & oui & NSM \\
T1048.003 & Exfiltration Over Unencrypted Non-C2 Protocol & oui & oui & NSM \\
T1049 & System Network Connections Discovery & oui & oui &  \\
T1052 & Exfiltration Over Physical Medium & oui & oui &  \\
T1052.001 & Exfiltration over USB & oui & oui &  \\
T1053 & Scheduled Task/Job & oui & oui &  \\
T1053.002 & At & oui & oui &  \\
T1053.005 & Scheduled Task & oui & oui &  \\
T1055 & Process Injection & oui & oui & ETW \\
T1055.001 & Dynamic-link Library Injection & oui & oui &  \\
T1055.002 & Portable Executable Injection & oui & oui &  \\
T1055.003 & Thread Execution Hijacking & oui & oui & ETW \\
T1055.004 & Asynchronous Procedure Call & oui & oui & ETW \\
T1055.005 & Thread Local Storage & oui & — & EDR \\
T1055.011 & Extra Window Memory Injection & oui & oui & ETW \\
T1055.012 & Process Hollowing & oui & oui & ETW \\
T1055.013 & Process Doppelgänging & oui & oui & ETW \\
T1055.015 & ListPlanting & oui & oui & ETW \\
T1056 & Input Capture & oui & oui &  \\
T1056.001 & Keylogging & oui & oui &  \\
T1056.002 & GUI Input Capture & oui & oui &  \\
T1056.003 & Web Portal Capture & oui & oui & 4.0 \\
T1056.004 & Credential API Hooking & oui & — &  \\
T1057 & Process Discovery & oui & oui &  \\
T1059 & Command and Scripting Interpreter & oui & oui &  \\
T1059.001 & PowerShell & oui & oui &  \\
T1059.003 & Windows Command Shell & oui & oui & EDR \\
T1059.005 & Visual Basic & oui & oui &  \\
T1059.006 & Python & oui & oui & EDR \\
T1059.007 & JavaScript & oui & oui & EDR \\
T1059.010 & AutoHotKey \& AutoIT & oui & oui &  \\
T1059.011 & Lua & oui & oui &  \\
T1068 & Exploitation for Privilege Escalation & oui & oui &  \\
T1069 & Permission Groups Discovery & oui & oui &  \\
T1069.001 & Local Groups & oui & oui &  \\
T1069.002 & Domain Groups & oui & oui &  \\
T1070 & Indicator Removal & oui & oui &  \\
T1070.003 & Clear Command History & oui & oui &  \\
T1070.004 & File Deletion & oui & oui &  \\
T1070.005 & Network Share Connection Removal & oui & oui & NSM \\
T1070.006 & Timestomp & oui & oui & EDR \\
T1070.007 & Clear Network Connection History and Configurations & oui & oui & EDR \\
T1070.008 & Clear Mailbox Data & oui & oui & cloud \\
T1070.009 & Clear Persistence & oui & oui &  \\
T1070.010 & Relocate Malware & oui & oui &  \\
T1071 & Application Layer Protocol & oui & oui & NSM \\
T1071.001 & Web Protocols & oui & oui & NSM \\
T1071.002 & File Transfer Protocols & oui & oui & NSM \\
T1071.003 & Mail Protocols & oui & oui & NSM \\
T1071.004 & DNS & oui & oui & NSM \\
T1071.005 & Publish/Subscribe Protocols & oui & oui & NSM \\
T1072 & Software Deployment Tools & oui & oui &  \\
T1074 & Data Staged & oui & oui &  \\
T1074.001 & Local Data Staging & oui & oui &  \\
T1074.002 & Remote Data Staging & oui & oui &  \\
T1078 & Valid Accounts & oui & oui &  \\
T1078.001 & Default Accounts & oui & oui &  \\
T1078.002 & Domain Accounts & oui & oui &  \\
T1078.003 & Local Accounts & oui & oui &  \\
T1080 & Taint Shared Content & oui & oui &  \\
T1082 & System Information Discovery & oui & oui &  \\
T1083 & File and Directory Discovery & oui & oui &  \\
T1087 & Account Discovery & oui & oui &  \\
T1087.001 & Local Account & oui & oui &  \\
T1087.002 & Domain Account & oui & oui & NSM \\
T1087.003 & Email Account & oui & oui &  \\
T1090 & Proxy & oui & oui & NSM \\
T1090.001 & Internal Proxy & oui & oui &  \\
T1090.002 & External Proxy & oui & oui &  \\
T1090.003 & Multi-hop Proxy & oui & oui & NSM \\
T1090.004 & Domain Fronting & oui & oui & NSM \\
T1091 & Replication Through Removable Media & oui & oui &  \\
T1092 & Communication Through Removable Media & oui & oui &  \\
T1095 & Non-Application Layer Protocol & oui & oui & NSM \\
T1098 & Account Manipulation & oui & oui &  \\
T1098.002 & Additional Email Delegate Permissions & oui & oui & cloud \\
T1098.005 & Device Registration & oui & oui &  \\
T1098.007 & Additional Local or Domain Groups & oui & oui &  \\
T1102 & Web Service & oui & oui & NSM \\
T1102.001 & Dead Drop Resolver & oui & oui & ETW \\
T1102.002 & Bidirectional Communication & oui & oui & ETW \\
T1102.003 & One-Way Communication & oui & oui & ETW \\
T1104 & Multi-Stage Channels & oui & oui &  \\
T1105 & Ingress Tool Transfer & oui & oui &  \\
T1106 & Native API & oui & oui &  \\
T1110 & Brute Force & oui & oui &  \\
T1110.001 & Password Guessing & oui & oui &  \\
T1110.002 & Password Cracking & oui & oui &  \\
T1110.003 & Password Spraying & oui & oui &  \\
T1110.004 & Credential Stuffing & oui & oui &  \\
T1111 & Multi-Factor Authentication Interception & oui & oui &  \\
T1112 & Modify Registry & oui & oui &  \\
T1113 & Screen Capture & oui & oui &  \\
T1114 & Email Collection & oui & oui &  \\
T1114.001 & Local Email Collection & oui & oui &  \\
T1114.002 & Remote Email Collection & oui & oui & cloud \\
T1114.003 & Email Forwarding Rule & oui & oui & cloud \\
T1115 & Clipboard Data & oui & oui &  \\
T1119 & Automated Collection & oui & oui &  \\
T1120 & Peripheral Device Discovery & oui & oui &  \\
T1123 & Audio Capture & oui & oui &  \\
T1124 & System Time Discovery & oui & oui & EDR, ETW \\
T1125 & Video Capture & oui & oui &  \\
T1127 & Trusted Developer Utilities Proxy Execution & oui & oui &  \\
T1127.001 & MSBuild & oui & oui & EDR \\
T1127.002 & ClickOnce & oui & oui &  \\
T1127.003 & JamPlus & oui & oui &  \\
T1129 & Shared Modules & oui & oui &  \\
T1132 & Data Encoding & oui & oui & NSM \\
T1132.001 & Standard Encoding & oui & oui & EDR \\
T1132.002 & Non-Standard Encoding & oui & oui & EDR \\
T1133 & External Remote Services & oui & oui &  \\
T1134 & Access Token Manipulation & oui & oui & ETW \\
T1134.001 & Token Impersonation/Theft & oui & oui & ETW \\
T1134.002 & Create Process with Token & oui & oui & ETW \\
T1134.003 & Make and Impersonate Token & oui & oui & ETW \\
T1134.004 & Parent PID Spoofing & oui & oui & ETW \\
T1134.005 & SID-History Injection & oui & oui & ETW \\
T1135 & Network Share Discovery & oui & oui & ETW \\
T1136 & Create Account & oui & oui &  \\
T1136.001 & Local Account & oui & oui &  \\
T1136.002 & Domain Account & oui & oui &  \\
T1137 & Office Application Startup & oui & oui &  \\
T1137.001 & Office Template Macros & oui & oui &  \\
T1137.002 & Office Test & oui & oui &  \\
T1137.003 & Outlook Forms & oui & oui &  \\
T1137.004 & Outlook Home Page & oui & oui &  \\
T1137.005 & Outlook Rules & oui & oui &  \\
T1137.006 & Add-ins & oui & oui &  \\
T1140 & Deobfuscate/Decode Files or Information & oui & oui &  \\
T1176 & Software Extensions & oui & oui &  \\
T1176.001 & Browser Extensions & oui & oui &  \\
T1176.002 & IDE Extensions & oui & oui &  \\
T1185 & Browser Session Hijacking & oui & oui &  \\
T1187 & Forced Authentication & oui & oui & NSM \\
T1189 & Drive-by Compromise & oui & oui & ETW, NSM \\
T1190 & Exploit Public-Facing Application & oui & oui & 4.0 \\
T1195 & Supply Chain Compromise & oui & oui & NSM \\
T1195.001 & Compromise Software Dependencies and Development Tools & oui & oui & NSM \\
T1195.002 & Compromise Software Supply Chain & oui & oui & NSM \\
T1195.003 & Compromise Hardware Supply Chain & oui & oui &  \\
T1197 & BITS Jobs & oui & oui &  \\
T1199 & Trusted Relationship & oui & oui &  \\
T1200 & Hardware Additions & oui & oui &  \\
T1201 & Password Policy Discovery & oui & oui &  \\
T1202 & Indirect Command Execution & oui & oui &  \\
T1203 & Exploitation for Client Execution & oui & oui &  \\
T1204 & User Execution & oui & oui &  \\
T1204.001 & Malicious Link & oui & oui & NSM \\
T1204.002 & Malicious File & oui & oui &  \\
T1204.004 & Malicious Copy and Paste & oui & oui & NSM \\
T1204.005 & Malicious Library & oui & oui &  \\
T1205 & Traffic Signaling & oui & oui &  \\
T1205.001 & Port Knocking & oui & oui &  \\
T1205.002 & Socket Filters & oui & oui & NSM \\
T1207 & Rogue Domain Controller & oui & — & cloud, NSM \\
T1210 & Exploitation of Remote Services & oui & oui & NSM \\
T1211 & Exploitation for Stealth & oui & oui &  \\
T1212 & Exploitation for Credential Access & oui & oui &  \\
T1213 & Data from Information Repositories & oui & oui & cloud \\
T1213.002 & Sharepoint & — & oui & cloud \\
T1213.006 & Databases & oui & oui &  \\
T1216 & System Script Proxy Execution & oui & oui &  \\
T1216.001 & PubPrn & oui & oui &  \\
T1216.002 & SyncAppvPublishingServer & oui & oui &  \\
T1217 & Browser Information Discovery & oui & oui &  \\
T1218 & System Binary Proxy Execution & oui & oui &  \\
T1218.001 & Compiled HTML File & oui & oui &  \\
T1218.002 & Control Panel & oui & oui &  \\
T1218.003 & CMSTP & oui & oui &  \\
T1218.004 & InstallUtil & oui & oui &  \\
T1218.005 & Mshta & oui & oui &  \\
T1218.007 & Msiexec & oui & oui &  \\
T1218.008 & Odbcconf & oui & oui &  \\
T1218.009 & Regsvcs/Regasm & oui & oui &  \\
T1218.010 & Regsvr32 & oui & oui &  \\
T1218.011 & Rundll32 & oui & oui &  \\
T1218.012 & Verclsid & oui & oui &  \\
T1218.013 & Mavinject & oui & oui &  \\
T1218.014 & MMC & oui & oui &  \\
T1218.015 & Electron Applications & oui & oui &  \\
T1219 & Remote Access Tools & oui & oui &  \\
T1219.001 & IDE Tunneling & oui & oui & NSM \\
T1219.002 & Remote Desktop Software & oui & oui &  \\
T1219.003 & Remote Access Hardware & oui & oui &  \\
T1220 & XSL Script Processing & oui & oui &  \\
T1221 & Template Injection & oui & oui &  \\
T1222 & File and Directory Permissions Modification & oui & oui &  \\
T1222.001 & Windows Permissions & oui & oui &  \\
T1480 & Execution Guardrails & oui & oui &  \\
T1480.001 & Environmental Keying & oui & oui &  \\
T1480.002 & Mutual Exclusion & oui & oui &  \\
T1482 & Domain Trust Discovery & oui & oui &  \\
T1484 & Domain or Tenant Policy Modification & oui & oui &  \\
T1484.001 & Group Policy Modification & oui & oui &  \\
T1484.002 & Trust Modification & oui & oui &  \\
T1485 & Data Destruction & oui & oui &  \\
T1486 & Data Encrypted for Impact & oui & oui &  \\
T1489 & Service Stop & oui & oui &  \\
T1490 & Inhibit System Recovery & oui & oui &  \\
T1491 & Defacement & oui & oui &  \\
T1491.001 & Internal Defacement & oui & oui &  \\
T1491.002 & External Defacement & oui & oui & NSM \\
T1495 & Firmware Corruption & oui & oui &  \\
T1496 & Resource Hijacking & oui & oui & ETW \\
T1496.001 & Compute Hijacking & oui & oui &  \\
T1496.002 & Bandwidth Hijacking & oui & oui &  \\
T1497 & Virtualization/Sandbox Evasion & oui & oui &  \\
T1497.001 & System Checks & oui & oui &  \\
T1497.002 & User Activity Based Checks & oui & oui &  \\
T1497.003 & Time Based Checks & oui & oui &  \\
T1498 & Network Denial of Service & oui & oui &  \\
T1498.001 & Direct Network Flood & oui & oui &  \\
T1498.002 & Reflection Amplification & oui & oui & ETW \\
T1499 & Endpoint Denial of Service & oui & oui &  \\
T1499.001 & OS Exhaustion Flood & oui & oui & NSM \\
T1499.002 & Service Exhaustion Flood & oui & oui & ETW \\
T1499.003 & Application Exhaustion Flood & oui & oui & ETW \\
T1499.004 & Application or System Exploitation & oui & oui &  \\
T1505 & Server Software Component & oui & oui &  \\
T1505.001 & SQL Stored Procedures & oui & oui &  \\
T1505.002 & Transport Agent & oui & oui &  \\
T1505.003 & Web Shell & oui & oui & NSM \\
T1505.004 & IIS Components & oui & oui & 4.0 \\
T1505.005 & Terminal Services DLL & oui & oui &  \\
T1518 & Software Discovery & oui & oui &  \\
T1518.001 & Security Software Discovery & oui & oui &  \\
T1518.002 & Backup Software Discovery & oui & oui &  \\
T1529 & System Shutdown/Reboot & oui & oui &  \\
T1531 & Account Access Removal & oui & oui &  \\
T1534 & Internal Spearphishing & oui & oui & cloud \\
T1539 & Steal Web Session Cookie & oui & oui &  \\
T1542 & Pre-OS Boot & oui & oui &  \\
T1542.001 & System Firmware & oui & oui &  \\
T1542.002 & Component Firmware & oui & oui & firmware \\
T1542.003 & Bootkit & oui & oui &  \\
T1543 & Create or Modify System Process & oui & oui &  \\
T1543.003 & Windows Service & oui & oui &  \\
T1546 & Event Triggered Execution & oui & oui &  \\
T1546.001 & Change Default File Association & oui & oui &  \\
T1546.002 & Screensaver & oui & oui &  \\
T1546.003 & Windows Management Instrumentation Event Subscription & oui & oui &  \\
T1546.007 & Netsh Helper DLL & oui & oui &  \\
T1546.008 & Accessibility Features & oui & oui &  \\
T1546.009 & AppCert DLLs & oui & oui &  \\
T1546.010 & AppInit DLLs & oui & oui &  \\
T1546.011 & Application Shimming & oui & oui &  \\
T1546.012 & Image File Execution Options Injection & oui & oui &  \\
T1546.013 & PowerShell Profile & oui & oui &  \\
T1546.015 & Component Object Model Hijacking & oui & oui &  \\
T1546.016 & Installer Packages & oui & oui &  \\
T1546.018 & Python Startup Hooks* & — & oui &  \\
T1547 & Boot or Logon Autostart Execution & oui & oui &  \\
T1547.001 & Registry Run Keys / Startup Folder & oui & oui &  \\
T1547.002 & Authentication Package & oui & oui &  \\
T1547.003 & Time Providers & oui & oui &  \\
T1547.004 & Winlogon Helper DLL & oui & oui &  \\
T1547.005 & Security Support Provider & oui & oui &  \\
T1547.008 & LSASS Driver & oui & oui &  \\
T1547.009 & Shortcut Modification & oui & oui &  \\
T1547.010 & Port Monitors & oui & oui &  \\
T1547.012 & Print Processors & oui & oui &  \\
T1547.014 & Active Setup & oui & oui &  \\
T1548 & Abuse Elevation Control Mechanism & oui & oui &  \\
T1548.002 & Bypass User Account Control & oui & oui &  \\
T1550 & Use Alternate Authentication Material & oui & oui &  \\
T1550.002 & Pass the Hash & oui & oui &  \\
T1550.003 & Pass the Ticket & oui & oui &  \\
T1552 & Unsecured Credentials & oui & oui &  \\
T1552.001 & Credentials In Files & oui & oui &  \\
T1552.002 & Credentials in Registry & oui & oui & EDR \\
T1552.003 & Shell History* & — & oui &  \\
T1552.004 & Private Keys & oui & oui &  \\
T1552.006 & Group Policy Preferences & oui & oui &  \\
T1553 & Subvert Trust Controls & oui & oui &  \\
T1553.002 & Code Signing & oui & oui &  \\
T1553.003 & SIP and Trust Provider Hijacking & oui & oui &  \\
T1553.004 & Install Root Certificate & oui & oui &  \\
T1553.005 & Mark-of-the-Web Bypass & oui & oui &  \\
T1553.006 & Code Signing Policy Modification & oui & oui &  \\
T1554 & Compromise Host Software Binary & oui & oui &  \\
T1555 & Credentials from Password Stores & oui & oui &  \\
T1555.003 & Credentials from Web Browsers & oui & oui &  \\
T1555.004 & Windows Credential Manager & oui & oui &  \\
T1555.005 & Password Managers & oui & oui &  \\
T1556 & Modify Authentication Process & oui & oui &  \\
T1556.001 & Domain Controller Authentication & oui & oui &  \\
T1556.002 & Password Filter DLL & oui & oui &  \\
T1556.005 & Reversible Encryption & oui & oui &  \\
T1556.006 & Multi-Factor Authentication & oui & oui &  \\
T1556.007 & Hybrid Identity & oui & oui &  \\
T1556.008 & Network Provider DLL & oui & oui &  \\
T1557 & Adversary-in-the-Middle & oui & oui &  \\
T1557.001 & Name Resolution Poisoning and SMB Relay & oui & oui & NSM \\
T1557.002 & ARP Cache Poisoning & oui & oui &  \\
T1557.003 & DHCP Spoofing & oui & oui & NSM \\
T1558 & Steal or Forge Kerberos Tickets & oui & oui &  \\
T1558.001 & Golden Ticket & oui & oui &  \\
T1558.002 & Silver Ticket & oui & oui &  \\
T1558.003 & Kerberoasting & oui & oui &  \\
T1558.004 & AS-REP Roasting & oui & oui &  \\
T1559 & Inter-Process Communication & oui & — &  \\
T1559.001 & Component Object Model & oui & oui &  \\
T1559.002 & Dynamic Data Exchange & oui & oui &  \\
T1560 & Archive Collected Data & oui & oui &  \\
T1560.001 & Archive via Utility & oui & oui &  \\
T1560.002 & Archive via Library & oui & oui &  \\
T1560.003 & Archive via Custom Method & oui & oui &  \\
T1561 & Disk Wipe & oui & oui &  \\
T1561.001 & Disk Content Wipe & oui & oui &  \\
T1561.002 & Disk Structure Wipe & oui & oui &  \\
T1563 & Remote Service Session Hijacking & oui & oui &  \\
T1563.002 & RDP Hijacking & oui & oui &  \\
T1564 & Hide Artifacts & oui & oui &  \\
T1564.001 & Hidden Files and Directories & oui & oui &  \\
T1564.002 & Hidden Users & oui & oui &  \\
T1564.003 & Hidden Window & oui & oui &  \\
T1564.004 & NTFS File Attributes & oui & oui & ETW \\
T1564.005 & Hidden File System & oui & oui & ETW \\
T1564.006 & Run Virtual Instance & oui & oui &  \\
T1564.007 & VBA Stomping & oui & oui &  \\
T1564.008 & Email Hiding Rules & oui & oui & cloud \\
T1564.010 & Process Argument Spoofing & oui & oui &  \\
T1564.011 & Ignore Process Interrupts & oui & oui &  \\
T1564.012 & File/Path Exclusions & oui & oui &  \\
T1565 & Data Manipulation & oui & oui &  \\
T1565.001 & Stored Data Manipulation & oui & oui &  \\
T1565.002 & Transmitted Data Manipulation & oui & oui &  \\
T1565.003 & Runtime Data Manipulation & oui & oui &  \\
T1566 & Phishing & oui & oui & cloud \\
T1566.001 & Spearphishing Attachment & oui & oui & cloud \\
T1566.002 & Spearphishing Link & oui & oui & cloud \\
T1566.003 & Spearphishing via Service & oui & oui &  \\
T1566.004 & Spearphishing Voice & — & oui & cloud \\
T1567 & Exfiltration Over Web Service & oui & oui &  \\
T1567.001 & Exfiltration to Code Repository & oui & oui &  \\
T1567.002 & Exfiltration to Cloud Storage & oui & oui &  \\
T1567.003 & Exfiltration to Text Storage Sites & oui & oui &  \\
T1567.004 & Exfiltration Over Webhook & oui & oui &  \\
T1568 & Dynamic Resolution & oui & oui &  \\
T1568.001 & Fast Flux DNS & oui & oui &  \\
T1568.002 & Domain Generation Algorithms & oui & oui &  \\
T1568.003 & DNS Calculation & oui & oui &  \\
T1569 & System Services & oui & oui &  \\
T1569.002 & Service Execution & oui & oui &  \\
T1570 & Lateral Tool Transfer & oui & oui &  \\
T1571 & Non-Standard Port & oui & oui &  \\
T1572 & Protocol Tunneling & oui & oui &  \\
T1573 & Encrypted Channel & oui & oui &  \\
T1573.001 & Symmetric Cryptography & oui & oui &  \\
T1573.002 & Asymmetric Cryptography & oui & oui &  \\
T1574 & Hijack Execution Flow & oui & oui &  \\
T1574.001 & DLL & oui & oui &  \\
T1574.005 & Executable Installer File Permissions Weakness & oui & oui &  \\
T1574.007 & Path Interception by PATH Environment Variable & oui & oui &  \\
T1574.008 & Path Interception by Search Order Hijacking & oui & oui &  \\
T1574.009 & Path Interception by Unquoted Path & oui & oui &  \\
T1574.010 & Services File Permissions Weakness & oui & oui &  \\
T1574.011 & Services Registry Permissions Weakness & oui & oui &  \\
T1574.012 & COR\_PROFILER & oui & oui &  \\
T1574.013 & KernelCallbackTable & oui & oui & ETW \\
T1574.014 & AppDomainManager & oui & oui &  \\
T1606 & Forge Web Credentials & oui & oui &  \\
T1606.001 & Web Cookies & oui & oui &  \\
T1606.002 & SAML Tokens & oui & oui & 4.0 \\
T1611 & Escape to Host & oui & oui &  \\
T1614 & System Location Discovery & oui & oui & ETW \\
T1614.001 & System Language Discovery & oui & oui & ETW \\
T1615 & Group Policy Discovery & oui & oui & NSM \\
T1620 & Reflective Code Loading & oui & oui & ETW \\
T1621 & Multi-Factor Authentication Request Generation & oui & oui &  \\
T1622 & Debugger Evasion & oui & oui & ETW \\
T1649 & Steal or Forge Authentication Certificates & oui & — &  \\
T1652 & Device Driver Discovery & oui & oui &  \\
T1653 & Power Settings & oui & oui &  \\
T1654 & Log Enumeration & oui & oui &  \\
T1657 & Financial Theft & oui & oui &  \\
T1659 & Content Injection & oui & oui & NSM \\
T1665 & Hide Infrastructure & oui & oui & NSM \\
T1667 & Email Bombing & oui & oui & cloud \\
T1668 & Exclusive Control & oui & oui &  \\
T1669 & Wi-Fi Networks & oui & oui &  \\
T1673 & Virtual Machine Discovery & oui & oui &  \\
T1674 & Input Injection & oui & oui &  \\
T1678 & Delay Execution & oui & oui &  \\
T1679 & Selective Exclusion & oui & oui &  \\
T1680 & Local Storage Discovery & oui & oui &  \\
T1684 & Social Engineering & oui & oui &  \\
T1684.001 & Impersonation & oui & oui & cloud \\
T1684.002 & Email Spoofing & — & oui & cloud \\
T1685 & Disable or Modify Tools & oui & oui &  \\
T1685.001 & Disable or Modify Windows Event Log & oui & oui &  \\
T1685.003 & Modify or Spoof Tool UI & oui & oui &  \\
T1685.005 & Clear Windows Event Logs & oui & oui &  \\
T1685.007 & Impair Autonomous Defenses* & — & — & IA \\
T1686 & Disable or Modify System Firewall & oui & oui &  \\
T1686.003 & Windows Host Firewall & oui & oui &  \\
T1687 & Exploitation for Defense Impairment & oui & oui &  \\
T1688 & Safe Mode Boot & oui & oui &  \\
T1689 & Downgrade Attack & oui & oui &  \\
T1690 & Prevent Command History Logging & oui & oui &  \\
T1698 & Indirect Prompt Injection* & — & — & IA \\
T1699 & Direct Prompt Injection* & — & — & IA \\
T1700 & AI-Enabled Exfiltration & oui & oui &  \\
\bottomrule
\end{longtable}
}

\noindent{\small ATT\&CK Enterprise 20.0 (2026-10-13, SHA-256 81113BA98AFA2928…) : 479 techniques et sous-techniques Windows ; 475 avec au moins une trace que WAC prélève, 4 dont toutes les sources sont hors du disque ; 6 (marquées *) n'ont pas d'analytique Windows dans ATT\&CK et sont décidées d'après la technique elle-même. \emph{Journaux} : un journal que WAC prend porte une trace de la technique ; \emph{Disque} : un artefact sur le disque en porte une, hors des journaux. Un journal d'événements n'est là que si la machine l'a tenu : WAC relève la configuration des canaux et la politique d'audit, qui le disent.}

\noindent{\small Sources non prises, et pourquoi : ETW, télémétrie vivante (ETW, appels d'API, compteurs) : lue par un consommateur à l'instant ; les fichiers de trace que Windows écrit sur son disque sont pris bruts (etw\_logs) ; NSM, capture réseau : sur le réseau, pas sur la machine ; cloud, service en nuage : dans les journaux du locataire ; le côté de la machine est pris (cloud\_identity, onedrive, intune\_logs) ; EDR, télémétrie propre à un EDR : dans sa console ; le dossier local de Defender for Endpoint est pris brut (defender\_endpoint) ; firmware, micrologiciel : non lu, WAC lit le disque ; 4.0, rôle serveur (IIS, AD FS, contrôleur de domaine) : version 4.0 ; IA, comportement d'un modèle d'IA : ATT\&CK ne lui connaît aucune trace sous Windows}


\section{Source attendue mais introuvable}\label{sec:attendu}

À l'inverse, certaines sources sont sur toute installation d'une version
donnée, et Windows continue de fonctionner sans elles : le fichier
\fichier{hosts}, \fichier{Amcache.hve} et \fichier{SRUDB.dat} depuis qu'ils
existent, les journaux d'événements, la clé des services. Le tableau ci-dessus
les signale par « toujours présent depuis ». Quand la collecte ne contient pas
une telle source, ce n'est ni une lecture ratée ni une absence d'activité :
c'est un constat, que le fichier de l'artefact énonce.

\begin{lstlisting}
{
  "Artifact": "Hosts file",
  "CollectionStatus": "ExpectedButMissing",
  "WindowsRelease": "Windows 11",
  "WindowsBuild": 26200,
  "ExpectedLocation": "Windows\\System32\\drivers\\etc\\hosts",
  "Error": "0x80070002 Le fichier spécifié est introuvable.",
  "Note": "This version of Windows always has this artefact's source
           (artefacts.hosts) and the collection does not hold it. [...]"
}
\end{lstlisting}

La conversion l'annonce en console et le consigne dans
\fichier{conversion.json}. Elle sait que la collecte scellée ne contient pas la
source ; que la machine ne l'avait pas se lit dans \fichier{investigation.json},
où l'extraction du dossier est consignée avec son résultat et le nombre de
fichiers pris. Un fichier effacé, un service arrêté ou un système endommagé
sont alors à expliquer.

\begin{attention}
Le statut n'est posé que là où la présence est certaine. Il ne l'est pas pour
ce qui n'existe qu'après un événement (la clé \fichier{USBSTOR} avant le premier
support USB), pour ce qui dépend d'un programme installé (navigateurs,
messagerie), ni sur les éditions serveur, qui restent à mesurer : l'absence y
reste un artefact non collecté ou une liste vide.
\end{attention}

\section{Données non décodées}

Un élément que \wac{} ne sait pas décoder (élément du shell inconnu, bloc
d'extension non reconnu) est accompagné de ses octets bruts en hexadécimal, pour
qu'un analyste puisse le décoder plus tard. Avec \opt{dump}, la valeur brute des
shellbags et des raccourcis est ajoutée systématiquement.

\section{Format non pris en charge}\label{sec:format}

Le format d'une pièce se lit sur la pièce elle-même --- sa signature ou sa
version interne ---, jamais sur la version de Windows : un cache de miniatures
de Vista retrouvé sur un Windows~10 reste un cache de Vista. Quand \wac{} n'a pas de
décodeur pour le format rencontré, qu'il soit ancien ou qu'il n'ait encore
jamais été vu, il ne convertit rien et le dit ; la pièce, elle, est dans la
consigne, prête pour un autre outil.

\begin{lstlisting}
{
  "SID": "S-1-5-21-1004336348-1177238915-682003330-1000",
  "SIDName": "jean",
  "Cache": "C:\\Users\\jean\\AppData\\Local\\Microsoft\\Windows\\Explorer\\thumbcache_96.db",
  "FormatVersion": 20,
  "DecodingStatus": "UnsupportedFormat",
  "FormatSignature": "0x00000014",
  "FormatName": "CMMM, Vista layout",
  "Note": "Collected raw, not converted: this version of WAC has no
           decoder for this format. The exhibit store holds the source."
}
\end{lstlisting}

\champ{FormatName} nomme le format quand le tableau de la
section~\ref{sec:versions} le connaît (colonne « Converti » à « non ») ; il
est absent pour un format inconnu. Le statut porte sur une pièce
(\champ{DecodingStatus}) quand l'artefact en compte plusieurs --- fichiers
Prefetch, index \fichier{\$I} de la corbeille, caches de miniatures, jump
lists (dont les raccourcis restent décodés, sans leur entrée DestList) --- et
sur l'artefact entier (\champ{CollectionStatus}) quand sa source est unique,
comme la valeur du ShimCache.

Certains formats ne portent pas de signature : c'est alors la version de
Windows examinée qui dit la forme de l'artefact. Avant Windows~10~1607,
\fichier{Amcache.hve} n'a pas ses clés \texttt{Inventory} ; avant
Windows~10~1709, BITS tient ses travaux dans \fichier{qmgr0.dat} ; avant
Windows~11, l'index de recherche est une base ESE. La collecte prend la pièce
comme d'habitude, et la conversion l'annonce au lieu de rendre une liste vide :

\begin{lstlisting}
{
  "Artifact": "Amcache applications",
  "CollectionStatus": "UnsupportedFormat",
  "WindowsRelease": "Windows 8.1",
  "WindowsBuild": 9600,
  "FormatName": "registry hive, File and Programs keys",
  "Location": "Windows\\AppCompat\\Programs\\Amcache.hve",
  "Note": "On this version of Windows this artefact (artefacts.amcache) is
           in a format this version of WAC has no decoder for. [...]"
}
\end{lstlisting}

\begin{attention}
Une pièce « UnsupportedFormat » n'est ni une liste vide ni une lecture ratée :
la trace existe, elle est prélevée, et son contenu reste à lire.
\end{attention}

\section{Chemins}

Les chemins et valeurs sont restitués \textbf{tels qu'ils figurent sur le
disque}. Rien n'est transformé : l'échappement n'intervient qu'à l'écriture du
JSON.

% ============================================================ 7
\chapter{Traces laissées sur la machine examinée}\label{chap:traces}
\markboth{Traces sur la machine examinée}{}

À lire avant d'employer \wac{} sur une affaire. Ce chapitre énumère ce que \wac{}
fait à la machine, opération par opération, y compris ce qu'il ne peut pas
éviter.

\section{Ce que WAC fait encore}

\begin{longtable}{@{}L{4.3cm}L{10.1cm}@{}}
\toprule
\bfseries Opération & \bfseries Trace \\
\midrule
\endhead
Lecture directe du volume & Un accès par volume effectivement lu (en général un
  seul). Aucun fichier ouvert, donc aucune date d'accès modifiée. Si l'audit
  d'accès aux objets est activé, l'ouverture du volume peut être journalisée
  (Security 4656/4663). \\
Écriture des résultats & Sur le support de collecte uniquement. Rien n'est écrit
  sur le disque examiné. \\
Rejeu des journaux de ruche & Sur la copie de travail uniquement, avec journal
  d'annulation. \\
État des services & Une requête en lecture au gestionnaire de services. Aucun
  service n'est ouvert ni modifié. \\
Processus & Un instantané du noyau ; aucun processus n'est ouvert. Relevé des
  propriétaires : une requête au service Terminal Services. \\
Sessions & Une requête en lecture à LSASS. \\
État réseau & Lectures en mémoire des tables de la pile TCP/IP : aucun paquet
  émis, aucun nom résolu. \texttt{iphlpapi.dll} n'est chargée que si
  \texttt{artefacts.network} est activé. \\
Heure de démarrage & Une requête en mémoire au noyau. \\
Journaux d'événements (\opt{events}) & Lus par le même accès direct au volume :
  le service de journalisation n'est pas sollicité. \\
Fichiers cités (\opt{binary}) & Lus par le même accès direct. Aucun fichier
  ouvert. \\
Vérification d'authenticité (\opt{binary}) & Catalogues lus par le même accès
  direct ; vérification en mémoire. Aucun service sollicité, aucun accès au
  magasin de certificats, aucun accès réseau. \\
\bottomrule
\end{longtable}

\section{Ce que WAC ne fait plus}

Les versions antérieures créaient un cliché instantané (VSS), passaient par COM et
WMI, interrogeaient le planificateur de tâches et le service de journalisation,
ouvraient chaque processus et chaque service, lisaient la liste des profils dans
le registre vivant et ouvraient chaque fichier cité pour en calculer
l'empreinte. Toutes ces opérations, qui laissaient des entrées dans les journaux
ou modifiaient des dates d'accès, ont été supprimées.

\section{Ce qu'aucun outil ne peut éviter}

Ces traces sont listées pour que l'analyste reconnaisse l'empreinte de \wac{} au
lieu de la prendre pour l'activité étudiée :

\begin{itemize}
  \item \textbf{lancer \fichier{WAC.exe} est une exécution} : elle produit une
        entrée Amcache, une entrée ShimCache à l'arrêt et, si \wac{} est lancé
        du disque de la machine plutôt que d'une clé, un fichier Prefetch et une
        valeur BAM. \wac{} n'écrit lui-même aucun événement, mais une
        surveillance en place enregistre son exécution : Security 4688 et 4689
        si l'audit de création de processus est actif, Sysmon (le processus,
        ses lectures brutes du volume, les fichiers qu'il crée sur la clé), un
        EDR. Attendez-vous à retrouver \wac{} dans les artefacts qu'il
        collecte ;
  \item \textbf{brancher la clé USB} l'inscrit dans \fichier{USBSTOR},
        \fichier{MountedDevices} et les journaux associés ;
  \item \textbf{ouvrir une session} pour lancer l'outil crée une session et ses
        entrées de journal ;
  \item \wac{} \textbf{occupe de la mémoire} pendant la collecte, ce qui peut
        évincer des pages vers le fichier d'échange. Sa lecture du volume, elle,
        contourne le cache (lecture non tamponnée) : les gigaoctets lus ne
        remplacent pas le cache de fichiers de la machine.
\end{itemize}

% ============================================================ 8
\chapter{Limites connues}

\begin{description}[style=nextline,leftmargin=1.2em]
  \item[Volumes ni NTFS ni FAT] Un fichier situé sur un CD-ROM n'est pas lu :
        il est signalé comme non lu.
  \item[Variables propres à un utilisateur] Un chemin exprimé avec
        \path|%APPDATA%|, \path|%LOCALAPPDATA%| ou \path|%USERPROFILE%| (dans une
        tâche planifiée, par exemple) n'est pas résolu : le compte concerné n'est
        pas connu à cet endroit, et deviner donnerait l'empreinte d'un autre
        fichier.
  \item[Exécutable supprimé] Un processus dont l'exécutable a été effacé du disque
        n'a pas d'empreinte : il ne reste que la mémoire, hors du périmètre de
        \wac{}.
  \item[Messages d'événements] Sur la machine de test, 99,6~\% des événements
        reçoivent un message lisible, identique à celui de Windows pour 98~\% des
        sortes d'événements. Le reste est expliqué dans le dossier de validation :
        écarts de Windows lui-même par rapport au manifeste, satellites de langue
        plus anciens que leur binaire. Les données, elles, sont toujours
        complètes.
  \item[Signés mais non authentifiés] Quelques sortes de fichiers réellement
        signés restent prélevées, dans le sens prudent : les applications du
        Microsoft Store (signées au niveau du paquet, et non fichier par fichier)
        et les scripts Windows Script Host signés seulement dans leur texte.
\end{description}

% ============================================================ 9
\chapter{Dépannage}

\begin{longtable}{@{}L{5cm}L{9.4cm}@{}}
\toprule
\bfseries Message ou constat & \bfseries Cause et remède \\
\midrule
\endhead
Échec de « Check administrator rights » & L'invite n'a pas été lancée en tant
  qu'administrateur. Relancez-la avec les droits d'administrateur. \\
« antérieur à Windows 10 » sous la version de Windows examinée & La collecte se
  fait, et chaque format que \wac{} ne décode pas sur cette version est déclaré
  (section~\ref{sec:format}) ; le cycle de test complet ne couvre encore que
  Windows~10 et 11. \\
Échec de « Checking the collection medium », répertoire non vide & Le dossier de
  sortie contient une collecte antérieure. Indiquez un dossier neuf avec
  \opt{output}. \\
Échec de « Checking the collection medium », place insuffisante & Le support
  est trop petit pour les options choisies (section~\ref{sec:place}). Utilisez un
  support plus grand ou retirez \opt{binary}. \\
Un fichier d'artefact porte \texttt{"CollectionStatus": "NotCollected"} &
  La source n'a pas pu être lue ; le motif est dans \champ{Error}. Les autres
  artefacts ne sont pas affectés. \\
Des pièces en échec au manifeste & Fichiers cités absents du disque (logiciel
  désinstallé, profil de service sans ruche), volume non NTFS\ldots{} Le motif est
  dans \champ{Result}. \\
Le journal indique des binaires « hachés sans être copiés » & La place a manqué
  pendant le prélèvement : les empreintes sont complètes, les copies non. \\
Besoin de comprendre un échec de lecture & Relancez avec \opt{loglevel=2} (et, pour
  la lecture NTFS, \opt{debug}) et consultez \fichier{WAC.exe.log}. \\
\bottomrule
\end{longtable}

% ============================================================ APPENDICES
\appendix

\chapter{Glossaire}

\begin{description}[style=nextline,leftmargin=1.2em]
  \item[Artefact] Trace laissée par l'activité du système ou d'un utilisateur, et
        exploitable en investigation.
  \item[Consigne] Répertoire des copies brutes, scellé, jamais réécrit.
  \item[Empreinte] Valeur calculée sur le contenu d'un fichier (MD5, SHA-1,
        SHA-256) : deux contenus différents ont des empreintes différentes.
  \item[Journal de transaction] Fichier \fichier{.LOG1}/\fichier{.LOG2} contenant
        les modifications récentes d'une ruche, pas encore écrites dans celle-ci.
  \item[Manifeste] Fichier qui identifie chaque pièce de la consigne.
  \item[\path|$MFT|] Table des fichiers d'un volume NTFS.
  \item[Ruche] Fichier constituant une partie de la base de registre
        (\fichier{SYSTEM}, \fichier{SOFTWARE}, \fichier{ntuser.dat}\ldots).
  \item[Sceau] Empreinte du manifeste, conservée à part : elle révèle toute
        retouche du manifeste.
  \item[Travail] Répertoire des copies de travail, sur lesquelles portent l'analyse
        et les éventuelles modifications.
  \item[WOF / Compact OS] Mode de stockage compressé des fichiers système de
        Windows~10 et~11.
\end{description}

\end{document}
